% Lesson 36 — Writing: Writes compositions on garbage collection, disposal, and recycling services — Anglais Tle A — Cours signé OnBuch+
% Programme officiel MINESEC. Compiler avec : tectonic EN36-writing-writes-compositions-on-garbage-collection-dispo.tex
\newcommand{\DOCMATIERE}{Anglais}
\newcommand{\DOCNIVEAU}{Tle A}
\newcommand{\DOCMODULE}{Chapter 3 : Environment, Well-being and Health}
\newcommand{\DOCLECON}{EN36}
\newcommand{\DOCTITRE}{Lesson 36 — Writing: Writes compositions on garbage collection, disposal, and recycling services}
% OnBuch+ — préambule « cours signés » ENRICHI (compilé avec tectonic / XeLaTeX)
% Système visuel « Pop » d'OnBuch+ : fond crème (jamais blanc), bordures encre
% épaisses, ombres « dures » décalées, titres Archivo Black en capitales.
% Version MATHÉMATIQUES : graphes (pgfplots/tikz), boîtes pédagogiques
% (définition, propriété, méthode, attention, le savais-tu, exercice résolu,
% exercice, TP, bilan), page de garde et signature OnBuch+ ancrée en bas.
%
% Macros attendues dans main.tex AVANT \input{preamble} :
%   \DOCMATIERE  \DOCNIVEAU  \DOCMODULE  \DOCLECON (numéro)  \DOCTITRE
\documentclass[11pt]{article}

\usepackage{fontspec}
\usepackage[english,french]{babel} % français = langue principale ; l'anglais via \eng{..} / textebox
\usepackage[a4paper,margin=2cm,top=3.4cm,bottom=2.8cm,headheight=1.4cm,footskip=1.3cm]{geometry}
\usepackage{amsmath,amssymb,amsfonts}
\usepackage{mathtools} % \xmapsto, \coloneqq... souvent utilisés par les modèles
\usepackage{cancel} % simplifications barrées (bilans de forces)
% \abs{z} (module, valeur absolue) et \norm{u} (norme), utilisés par les modèles
\providecommand{\abs}{}\let\abs\relax\DeclarePairedDelimiter{\abs}{\lvert}{\rvert}
\providecommand{\norm}{}\let\norm\relax\DeclarePairedDelimiter{\norm}{\lVert}{\rVert}
\usepackage[table]{xcolor} % [table] : fournit aussi \rowcolors (alternance de lignes)
\usepackage{pagecolor}
\usepackage{tikz}
\usetikzlibrary{arrows.meta,positioning,calc,shapes.geometric,shapes.misc,shapes.arrows,decorations.pathmorphing,decorations.pathreplacing,decorations.markings,patterns,shadows,mindmap,trees,fit,backgrounds,matrix,intersections,angles,quotes}
\usepackage{pgfplots}
\usepackage[version=4]{mhchem} % équations nucléaires \ce{^{235}_{92}U}
\usepackage[european,siunitx]{circuitikz} % schémas électriques
\pgfplotsset{compat=1.18}
\usepgfplotslibrary{fillbetween,groupplots} % aires entre courbes (calcul intégral)
\usepackage{enumitem}
\usepackage{fancyhdr}
% [most] charge toutes les bibliothèques tcolorbox (listings, minted, raster,
% documentation...) et gonfle inutilement la mémoire TeX sur un document long
% (~300 boîtes) : on ne charge que ce qui est réellement utilisé.
\usepackage[skins,breakable]{tcolorbox}
\usepackage{graphicx}
\usepackage{colortbl}
\usepackage{booktabs}
\usepackage{tabularx}
\usepackage{diagbox} % case d'en-tête diagonale des tableaux à double entrée
% Colonnes extensibles centrée / gauche / droite (C, L, R), souvent utilisées par les modèles
\newcolumntype{C}{>{\centering\arraybackslash}X}
\newcolumntype{L}{>{\raggedright\arraybackslash}X}
\newcolumntype{R}{>{\raggedleft\arraybackslash}X}
\usepackage{array}
\usepackage{multicol}
\usepackage{multirow}
\usepackage{siunitx}
% parse-numbers=false : accepte aussi \SI{2,5 \times 10^{-3}}{...}
\sisetup{locale=FR,output-decimal-marker={,},group-minimum-digits=4,parse-numbers=false}
\DeclareSIUnit\ohm{\ensuremath{\symup{\Omega}}} % U+2126 absent de la police
\DeclareSIUnit\division{div} % réglages d'oscilloscope (ms/div, V/div)
\DeclareSIUnit\tr{tr} % tours (vitesse de rotation, ex. \SI{1500}{\tr\per\minute})
\DeclareSIUnit\molar{M} % concentration molaire (ex. \SI{3}{\molar}, \SI{1}{\micro\molar})
\DeclareSIUnit\an{an} % année (le modèle confond parfois avec \year, un compteur TeX) — ex. \SI{5}{\kilo\meter\per\an}
\DeclareSIUnit\FCFA{FCFA} % franc CFA (ex. \SI{500}{\FCFA\per\kilo\gram})
\DeclareSIUnit\degreeBrix{\textdegree Brix} % teneur en sucre d'un sirop/jus (réfractométrie)
\DeclareSIUnit\month{mois} % le modèle utilise parfois \month, un compteur TeX, comme unité de durée
\usepackage{qrcode}
% \degree : commande fréquemment utilisée par les modèles pour un angle ou une
% température en dehors de \SI{}{} (non fournie par les paquets chargés).
\providecommand{\degree}{^{\circ}}
% \qed (fin de démonstration), utilisé par les modèles sans amsthm
\providecommand{\qed}{\hfill$\square$}
% \barre : double barre des valeurs interdites dans un tableau de variations (tkz-tab)
\providecommand{\barre}{\ensuremath{\|}}
% environnement proof (amsthm non chargé) utilisé par les modèles
\ifdefined\proof\else\newenvironment{proof}[1][Démonstration]{\par\noindent\textit{#1.}\ }{\qed\par}\fi
\usepackage{lastpage}
\usepackage{hyperref}
\hypersetup{hidelinks}

% ============================================================
% Palette « Pop » OnBuch+ — mêmes hex que la classe P (plus_ui.dart)
% ============================================================
\definecolor{poporange}{HTML}{F59321}
\definecolor{popdark}{HTML}{E07A0C}
\definecolor{popcream}{HTML}{FFF6E8}
\definecolor{popink}{HTML}{1C1714}
\definecolor{poppurple}{HTML}{7A5AE0}
\definecolor{popgreen}{HTML}{2E9E6A}
\definecolor{poppink}{HTML}{E0457B}
\definecolor{popblue}{HTML}{2F7DE1}
\definecolor{popgold}{HTML}{C98A12}
\definecolor{popmuted}{HTML}{8A7F76}
\definecolor{popline}{HTML}{EADFD2}
\definecolor{popgray}{HTML}{8A7F76}
\colorlet{popgrey}{popgray}
% Alias tolérants (le modèle nomme parfois les couleurs différemment)
\colorlet{popwhite}{white}
\colorlet{popblack}{popink}
\colorlet{popred}{poppink}
\colorlet{popyellow}{popgold}
% Teintes claires (fonds de tableaux/encadrés) — le modèle nomme parfois
% les couleurs avec un suffixe L (ex. popblueL) sans les avoir définies.
\colorlet{poporangeL}{poporange!20}
\colorlet{popdarkL}{popdark!20}
\colorlet{poppurpleL}{poppurple!20}
\colorlet{popgreenL}{popgreen!20}
\colorlet{poppinkL}{poppink!20}
\colorlet{popblueL}{popblue!20}
\colorlet{popgoldL}{popgold!20}
\colorlet{popredL}{popred!20}
\colorlet{popgrayL}{popgray!20}
% Teintes claires pour fonds de tableaux / zones
\colorlet{popblueL}{popblue!10!white}
\colorlet{poporangeL}{poporange!14!white}
\colorlet{popgreenL}{popgreen!12!white}
\colorlet{poppurpleL}{poppurple!12!white}
\colorlet{poppinkL}{poppink!12!white}
\colorlet{popgoldL}{popgold!14!white}

\pagecolor{popcream}

% ============================================================
% Typographie
% ============================================================
\defaultfontfeatures{Ligatures=TeX}
\setmainfont{PlusJakartaSans-Medium.ttf}[Path=../../../fonts/,BoldFont=PlusJakartaSans-Bold.ttf]
% Maths en sans-serif (Fira Math) avec chiffres et lettres droites de Plus
% Jakarta, pour que les formules se fondent dans le texte.
\usepackage[math-style=ISO,bold-style=ISO]{unicode-math}
\setmathfont{FiraMath-Regular.otf}[Path=../../../fonts/]
\setmathfont{PlusJakartaSans-Medium.ttf}[Path=../../../fonts/,range={up/num,up/{Latin,latin},\mathup}]
\usepackage{icomma} % virgule décimale sans espace en mode math
% Caractères mathématiques Unicode absents des polices de texte (Plus Jakarta,
% Archivo Black) : rendus par leur équivalent mathématique, en texte comme en
% formule — sinon ils s'affichent en carré vide (ex. « Arithmétique dans ℤ »).
\usepackage{newunicodechar}
\newunicodechar{ℝ}{\ensuremath{\mathbb{R}}}
\newunicodechar{ℤ}{\ensuremath{\mathbb{Z}}}
\newunicodechar{ℂ}{\ensuremath{\mathbb{C}}}
\newunicodechar{ℕ}{\ensuremath{\mathbb{N}}}
\newunicodechar{ℚ}{\ensuremath{\mathbb{Q}}}
\newunicodechar{𝔻}{\ensuremath{\mathbb{D}}}
\newunicodechar{α}{\ensuremath{\alpha}}
\newunicodechar{β}{\ensuremath{\beta}}
\newunicodechar{γ}{\ensuremath{\gamma}}
\newunicodechar{δ}{\ensuremath{\delta}}
\newunicodechar{ε}{\ensuremath{\varepsilon}}
\newunicodechar{θ}{\ensuremath{\theta}}
\newunicodechar{λ}{\ensuremath{\lambda}}
\newunicodechar{μ}{\ensuremath{\mu}}
\newunicodechar{σ}{\ensuremath{\sigma}}
\newunicodechar{τ}{\ensuremath{\tau}}
\newunicodechar{φ}{\ensuremath{\varphi}}
\newunicodechar{ψ}{\ensuremath{\psi}}
\newunicodechar{ω}{\ensuremath{\omega}}
\newunicodechar{Σ}{\ensuremath{\Sigma}}
% majuscules produites par \MakeUppercase dans les titres (α-aminés -> Α)
\newunicodechar{Α}{\ensuremath{\alpha}}
\newunicodechar{Β}{\ensuremath{\beta}}
\newunicodechar{Γ}{\ensuremath{\Gamma}}
\newunicodechar{Δ}{\ensuremath{\Delta}}
\newunicodechar{Ε}{\ensuremath{\varepsilon}}
\newunicodechar{Θ}{\ensuremath{\Theta}}
\newunicodechar{Λ}{\ensuremath{\Lambda}}
\newunicodechar{Μ}{\ensuremath{\mu}}
\newunicodechar{Τ}{\ensuremath{\tau}}
\newunicodechar{Φ}{\ensuremath{\Phi}}
\newunicodechar{Ψ}{\ensuremath{\Psi}}
\newunicodechar{Ω}{\ensuremath{\Omega}}
\newunicodechar{↦}{\ensuremath{\mapsto}}
\newunicodechar{∈}{\ensuremath{\in}}
\newunicodechar{∉}{\ensuremath{\notin}}
\newunicodechar{∧}{\ensuremath{\wedge}}
\newunicodechar{⇔}{\ensuremath{\Leftrightarrow}}
\newunicodechar{⟺}{\ensuremath{\Longleftrightarrow}}
\newunicodechar{⇒}{\ensuremath{\Rightarrow}}
\newunicodechar{⟹}{\ensuremath{\Longrightarrow}}
\newunicodechar{∘}{\ensuremath{\circ}}
\newunicodechar{∪}{\ensuremath{\cup}}
\newunicodechar{∩}{\ensuremath{\cap}}
\newunicodechar{≡}{\ensuremath{\equiv}}
\newunicodechar{⊂}{\ensuremath{\subset}}
\newunicodechar{⊥}{\ensuremath{\perp}}
\newunicodechar{∃}{\ensuremath{\exists}}
\newunicodechar{∀}{\ensuremath{\forall}}
\newunicodechar{∛}{\ensuremath{\sqrt[3]{\,}}}
\newunicodechar{∜}{\ensuremath{\sqrt[4]{\,}}}
\newunicodechar{⃗}{\ensuremath{{}^{\to}}}
\newunicodechar{ⁿ}{\ifmmode^{n}\else\textsuperscript{\ensuremath{n}}\fi}
\newunicodechar{ˣ}{\ifmmode^{x}\else\textsuperscript{\ensuremath{x}}\fi}
\newunicodechar{ᵇ}{\ifmmode^{b}\else\textsuperscript{\ensuremath{b}}\fi}
\newunicodechar{ʳ}{\ifmmode^{r}\else\textsuperscript{\ensuremath{r}}\fi}
\newunicodechar{ᵘ}{\ifmmode^{u}\else\textsuperscript{\ensuremath{u}}\fi}
\newunicodechar{ʸ}{\ifmmode^{y}\else\textsuperscript{\ensuremath{y}}\fi}
\newunicodechar{ᵃ}{\ifmmode^{a}\else\textsuperscript{\ensuremath{a}}\fi}
\newunicodechar{ᵉ}{\ifmmode^{e}\else\textsuperscript{\ensuremath{e}}\fi}
\newunicodechar{ᵏ}{\ifmmode^{k}\else\textsuperscript{\ensuremath{k}}\fi}
\newunicodechar{ᵖ}{\ifmmode^{p}\else\textsuperscript{\ensuremath{p}}\fi}
\newunicodechar{ᴺ}{\ifmmode^{N}\else\textsuperscript{\ensuremath{N}}\fi}
\newunicodechar{ᶜ}{\ifmmode^{c}\else\textsuperscript{\ensuremath{c}}\fi}
\newunicodechar{ʰ}{\ifmmode^{h}\else\textsuperscript{\ensuremath{h}}\fi}
\newunicodechar{ᵠ}{\ifmmode^{q}\else\textsuperscript{\ensuremath{q}}\fi}
\newunicodechar{ₙ}{\ifmmode_{n}\else\textsubscript{\ensuremath{n}}\fi}
\newunicodechar{ₐ}{\ifmmode_{a}\else\textsubscript{\ensuremath{a}}\fi}
\newunicodechar{ᵢ}{\ifmmode_{i}\else\textsubscript{\ensuremath{i}}\fi}
\newunicodechar{ᵣ}{\ifmmode_{r}\else\textsubscript{\ensuremath{r}}\fi}
\newunicodechar{ₖ}{\ifmmode_{k}\else\textsubscript{\ensuremath{k}}\fi}
\newunicodechar{ᵦ}{\ifmmode_{\beta}\else\textsubscript{\ensuremath{\beta}}\fi}
\newunicodechar{ₓ}{\ifmmode_{x}\else\textsubscript{\ensuremath{x}}\fi}
\newfontfamily\popdisplay{ArchivoBlack-Regular.ttf}[Path=../../../fonts/]
\newfontfamily\poplogo{SpaceGrotesk-Bold.ttf}[Path=../../../fonts/]
\newfontfamily\popsemi{PlusJakartaSans-SemiBold.ttf}[Path=../../../fonts/]
\newfontfamily\popxbold{PlusJakartaSans-ExtraBold.ttf}[Path=../../../fonts/]
\newfontfamily\popxboldls{PlusJakartaSans-ExtraBold.ttf}[Path=../../../fonts/,LetterSpace=3.0]

\setlist[enumerate]{leftmargin=1.7em,itemsep=5pt,topsep=4pt}
\setlist[itemize]{leftmargin=1.7em,itemsep=4pt,topsep=4pt,label={\color{poporange}\textbullet}}
\setlength{\parindent}{0pt}
\setlength{\parskip}{5pt}
\renewcommand{\arraystretch}{1.25}

% pgfplots : style Pop par défaut
\pgfplotsset{
  every axis/.append style={axis line style={popink,line width=1pt},tick style={popink},
    grid style={popline,line width=0.5pt},label style={font=\small},tick label style={font=\footnotesize}},
  popcurve/.style={poporange,line width=1.8pt,no markers,smooth},
}
% Même style disponible dans un \draw TikZ simple (hors pgfplots)
\tikzset{popcurve/.style={poporange,line width=1.8pt}}
\tikzset{popgrid/.style={popline,very thin}}
\colorlet{popcurve}{poporange} % le nom du style est parfois utilisé comme couleur

% ============================================================
% Pill / squiggle
% ============================================================
\newcommand{\poppill}[3]{%
  \tikz[baseline=-0.55ex]\node[draw=popink,line width=1.2pt,rounded corners=2.6pt,%
    fill=#2,inner xsep=6.5pt,inner ysep=3pt]{%
    \popxboldls\fontsize{7}{7}\selectfont\color{#3}\MakeUppercase{#1}};%
}
\newcommand{\popsquiggle}[1]{%
  \tikz[baseline=-0.4ex]{
    \draw[#1,line width=2.6pt,line cap=round]
      plot [smooth] coordinates {(0,0.09) (0.34,0.01) (0.68,0.21) (1.02,0.01) (1.36,0.10)};
  }%
}
% Mot-clé surligné (marqueur orange sous le texte)
\newcommand{\cle}[1]{{\popxbold\color{popdark}#1}}

% ============================================================
% En-tête / pied
% ============================================================
\pagestyle{fancy}
\fancyhf{}
\renewcommand{\headrulewidth}{0pt}
\renewcommand{\footrulewidth}{0pt}
\lhead{\raisebox{-0.15ex}{\poplogo\fontsize{13}{13}\selectfont\color{popink}OnBuch\color{poporange}+}}
\rhead{\poppill{\DOCMATIERE\ \textbullet\ \DOCNIVEAU\ \textbullet\ Leçon \DOCLECON}{white}{popink}}
\lfoot{\popxbold\fontsize{7.6}{7.6}\selectfont\color{popmuted}\MakeUppercase{Cours signé OnBuch+}\ \textbullet\ {\popsemi\DOCTITRE}}
\rfoot{\tikz[baseline=-0.6ex]\node[rounded corners=8.5pt,fill=popink,minimum height=17pt,minimum width=17pt,inner xsep=5pt,inner sep=0]{\popxbold\fontsize{8}{8}\selectfont\color{popcream}\thepage\,/\,\pageref*{LastPage}};}

% ============================================================
% Titres
% ============================================================
\newcommand{\chaptitre}[2]{%
  \begin{tcolorbox}[enhanced,colback=poporange,colframe=popink,boxrule=2.6pt,arc=11pt,
      drop shadow=popink,left=15pt,right=15pt,top=12pt,bottom=11pt,before skip=3pt,after skip=18pt]
    \poppill{#2}{popink}{popcream}\par\vspace{9pt}
    {\raggedright\hyphenpenalty=10000\exhyphenpenalty=10000\popdisplay\fontsize{22}{24}\selectfont\color{popcream}\MakeUppercase{#1}\par}
  \end{tcolorbox}
}
% Section numérotée : pastille orange avec le numéro + titre Archivo
\newcounter{popsec}
\newcommand{\coursec}[1]{%
  \stepcounter{popsec}\setcounter{popsub}{0}%
  \par\vspace{16pt}\noindent
  \begin{minipage}{\linewidth}
  \tikz[baseline=-0.7ex]\node[draw=popink,line width=1.6pt,fill=poporange,rounded corners=4pt,minimum size=22pt,inner sep=0]{\popdisplay\fontsize{12}{12}\selectfont\color{popcream}\thepopsec};%
  \hspace{8pt}{\popdisplay\fontsize{15}{17}\selectfont\color{popink}\MakeUppercase{#1}}\par
  \vspace{4pt}\noindent{\color{poporange}\rule{36pt}{3.6pt}}
  \end{minipage}\par\nopagebreak\vspace{8pt}
}
\newcounter{popsub}
\newcommand{\courssub}[1]{%
  \stepcounter{popsub}%
  \par\vspace{9pt}\noindent{\popxbold\fontsize{11.5}{13}\selectfont\color{popink}{\color{poporange}\thepopsec.\thepopsub}\hspace{6pt}#1}\par\nopagebreak\vspace{4pt}
}

% ============================================================
% Boîtes pédagogiques — même recette : fond blanc, bordure épaisse colorée,
% ombre dure encre, badge plein. Titre optionnel [#1].
% ============================================================
\tcbset{popbase/.style={breakable,enhanced,colback=white,boxrule=2.3pt,arc=9pt,
  shadow={3pt}{-3pt}{0pt}{popink},left=14pt,right=14pt,top=11pt,bottom=11pt,before skip=11pt,after skip=15pt,
  fontupper=\popsemi\fontsize{10.6}{14.5}\selectfont\color{popink},
  fontlower=\popsemi\fontsize{10.6}{14.5}\selectfont\color{popink}}}
% Coche dessinée (le glyphe ✓ n'existe pas dans les polices du document)
\newcommand{\popcheck}[1]{\tikz[baseline=-0.5ex]\draw[#1,line width=1.6pt,line cap=round,line join=round](-0.13,0.0)--(-0.04,-0.09)--(0.13,0.1);}
\providecommand{\coche}{\popcheck{popgreen}}
\providecommand{\resultat}[1]{\par\smallskip\noindent\fcolorbox{popgreen}{popgreenL}{\parbox{\dimexpr\linewidth-2\fboxsep-2\fboxrule\relax}{\textbf{Résultat.} #1}}\par\smallskip}
\newcommand{\popboxhead}[3]{\poppill{#1}{#2}{white}\ifx\relax#3\relax\else\hspace{6pt}{\popxbold\fontsize{10.6}{12}\selectfont\color{popink}#3}\fi\par\vspace{7pt}}

\newtcolorbox{aretenir}[1][]{popbase,colframe=popgreen,before upper={\popboxhead{À retenir}{popgreen}{#1}}}
% Texte en anglais (document de lecture, dialogue, extrait) : cadre
% d'étude. Un titre optionnel (source, auteur) s'affiche dans l'en-tête.
\newtcolorbox{textebox}[1][]{popbase,colframe=popblue,before upper={\popboxhead{Text}{popblue}{#1}\begin{otherlanguage}{english}},after upper={\end{otherlanguage}}}
% Marques ✓ / ✗ (forme correcte / fautive) souvent utilisées par les modèles
\providecommand{\cross}{\textcolor{poppink}{\ensuremath{\times}}}
\providecommand{\xmark}{\cross}\providecommand{\crossmark}{\cross}\providecommand{\cmark}{\ensuremath{\checkmark}}
% Ligne de réponse / justification à compléter (inventée par les modèles)
\providecommand{\justification}[1]{\par\noindent\textit{Justification~:} \dotfill\par}
% \textipa (package tipa non chargé) : la police ne couvre pas l'API -> simple passage
\providecommand{\textipa}[1]{#1}
% Soulignement (ulem non chargé) : \uline, \ul
\providecommand{\uline}[1]{\underline{#1}}\providecommand{\ul}[1]{\underline{#1}}
% \glossaire{\item ...} inventée par les modèles : liste de glossaire
\providecommand{\glossaire}[1]{\begin{itemize}#1\end{itemize}}
% \alert (beamer) inventée par les modèles : texte mis en évidence
\providecommand{\alert}[1]{\textbf{\textcolor{poppink}{#1}}}
% variantes ASCII d'ordinaux écrites en macro
\providecommand{\iere}{\textsuperscript{re}}\providecommand{\ieme}{\textsuperscript{e}}\providecommand{\ere}{\textsuperscript{re}}\providecommand{\eme}{\textsuperscript{e}}\providecommand{\er}{\textsuperscript{er}}\providecommand{\ie}{\textsuperscript{e}}
% Ordinaux français écrits en macro par les modèles : 1\ière, 3\ième
\newcommand{\ière}{\textsuperscript{re}}\newcommand{\ième}{\textsuperscript{e}}
% \exemple (inventée par les modèles) : étiquette « Exemple : »
\providecommand{\exemple}{\textbf{Exemple~:}\ }
% \square utilisable aussi hors mode math (cases à cocher)
\AtBeginDocument{\let\squareMath\square\renewcommand{\square}{\ensuremath{\squareMath}}}
% Mot ou phrase en anglais à l'intérieur d'un paragraphe français
\newcommand{\eng}[1]{\ifmmode\text{\textit{#1}}\else\begingroup\selectlanguage{english}\itshape #1\endgroup\fi}
\let\esp\eng % alias toléré
% flèches d'intonation inventées par les modèles
\providecommand{\textsearrow}{}\renewcommand{\textsearrow}{\ensuremath{\searrow}}\providecommand{\textnearrow}{}\renewcommand{\textnearrow}{\ensuremath{\nearrow}}
% \fr{..} : mot français (inventé par les modèles)
\providecommand{\fr}[1]{\textit{#1}}
\newtcolorbox{exemplebox}[1][]{popbase,colframe=poporange,before upper={\popboxhead{Exemple}{poporange}{#1}}}
\newtcolorbox{definition}[1][]{popbase,colframe=popblue,before upper={\popboxhead{Définition}{popblue}{#1}}}
\newtcolorbox{propriete}[1][]{popbase,colframe=poppurple,before upper={\popboxhead{Propriété}{poppurple}{#1}}}
\newtcolorbox{methode}[1][]{popbase,colframe=popgold,before upper={\popboxhead{Méthode}{popgold}{#1}}}
\newtcolorbox{attention}[1][]{popbase,colframe=poppink,before upper={\popboxhead{Attention}{poppink}{#1}}}
\newtcolorbox{experience}[1][]{popbase,colframe=popblue,colback=popblueL,before upper={\popboxhead{Expérience}{popblue}{#1}}}
% « Le savais-tu ? » : carte violette pleine (contexte, culture, Cameroun)
\newtcolorbox{savaistu}[1][]{popbase,colback=poppurple,colframe=popink,
  fontupper=\popsemi\fontsize{10.4}{14.2}\selectfont\color{white},
  before upper={\poppill{Le savais-tu\,?}{white}{poppurple}\ifx\relax#1\relax\else\hspace{6pt}{\popxbold\color{white}#1}\fi\par\vspace{7pt}}}
% Exercice résolu : énoncé en haut, \tcblower puis la solution sur fond crème
\newcounter{popexo}\newcounter{popexr}
\newtcolorbox{exoresolu}[1][]{popbase,colframe=poporange,colbacklower=popcream,
  segmentation style={popink,line width=1.2pt,dashed},
  before upper={\stepcounter{popexr}\popboxhead{Exercice résolu \thepopexr}{poporange}{#1}},
  before lower={\poppill{Solution}{popink}{popcream}\par\vspace{6pt}}}
% Exercice (non résolu en place) — la correction est dans la partie Corrigés
\newtcolorbox{exercice}[1][]{popbase,colframe=popink,
  before upper={\stepcounter{popexo}\popboxhead{Exercice \thepopexo}{popink}{#1}}}
% Corrigé
\newtcolorbox{corrige}[1][]{popbase,colframe=popgreen,colback=popgreenL,
  before upper={\popboxhead{Corrigé}{popgreen}{#1}}}
% Objectifs / prérequis / bilan
\newtcolorbox{objectifs}{popbase,colframe=popink,colback=poporangeL,
  before upper={\popboxhead{Objectifs de la leçon}{popink}{}}}
\newtcolorbox{prerequis}{popbase,colframe=popink,colback=popblueL,
  before upper={\popboxhead{Prérequis}{popblue}{}}}
\newtcolorbox{bilan}{popbase,colframe=popink,colback=white,boxrule=2.8pt,
  before upper={\popboxhead{Fiche bilan}{poporange}{L'essentiel à connaître pour l'examen}}}
% Figure encadrée avec légende
\newtcolorbox{popfigure}[1][]{enhanced,breakable=false,colback=white,colframe=popline,boxrule=1.4pt,arc=8pt,
  left=10pt,right=10pt,top=10pt,bottom=8pt,before skip=10pt,after skip=12pt,halign=center,
  lower separated=false,halign lower=center,
  fontlower=\popsemi\fontsize{9}{11.5}\selectfont\color{popmuted},
  #1}
\newcommand{\legende}[1]{\par\vspace{6pt}{\popsemi\fontsize{9}{11.5}\selectfont\color{popmuted}#1}}

% ============================================================
% Signature OnBuch+
% ============================================================
\newcommand{\poplogoinline}[1]{{\poplogo\fontsize{#1}{#1}\selectfont\color{popink}OnBuch\color{poporange}+}}
\newcommand{\signatureonbuch}{%
  \begin{tcolorbox}[enhanced,colback=popink,colframe=popink,boxrule=2.2pt,arc=10pt,
      drop shadow=poporange,left=14pt,right=14pt,top=10pt,bottom=10pt,before skip=0pt,after skip=0pt]
    \tikz[baseline=-0.7ex]\node[circle,fill=popgreen,draw=popcream,line width=1.3pt,minimum size=22pt,inner sep=0]{\popcheck{white}};%
    \hspace{9pt}\begin{minipage}[c]{0.62\linewidth}
      {\popxbold\fontsize{12}{13}\selectfont\color{popcream}Signé\ \textbullet\ {\poplogo OnBuch\color{poporange}+}}\par\vspace{2pt}
      {\raggedright\popsemi\fontsize{8}{10.5}\selectfont\color{popline}\MakeUppercase{Équipe pédagogique OnBuch+}\\\MakeUppercase{Conforme au programme officiel MINESEC}\par}
    \end{minipage}\hfill
    {\poplogo\fontsize{22}{22}\selectfont\color{popcream}OnBuch\color{poporange}+}
  \end{tcolorbox}%
}

% ============================================================
% Page de garde
% #1 titre · #2 matière · #3 niveau · #4 module · #5 n° leçon · #6 durée ·
% #7 description / objectifs (texte ou itemize)
% ============================================================
\newcommand{\pagegarde}[7]{%
  \thispagestyle{empty}
  \newgeometry{a4paper,margin=1.7cm,top=1.6cm,bottom=1.4cm}%
  \begin{tikzpicture}[remember picture,overlay]
    \fill[poporange!18] ([xshift=-3.2cm,yshift=-3.6cm]current page.north east) circle (4.6cm);
    \fill[poppurple!14] ([xshift=2.4cm,yshift=7.6cm]current page.south west) circle (3.8cm);
  \end{tikzpicture}%
  {\poplogo\fontsize{30}{30}\selectfont\color{popink}OnBuch\color{poporange}+}\hfill
  \raisebox{0.6ex}{\poppill{Cours signé \textbullet\ Premium}{popink}{popcream}}\par
  \vspace{1pt}\popsquiggle{poppurple}\par
  \vspace{8pt}
  \begin{tcolorbox}[enhanced,colback=poporange,colframe=popink,boxrule=2.8pt,arc=13pt,
      drop shadow=popink,left=18pt,right=18pt,top=12pt,bottom=13pt,after skip=10pt]
    \poppill{#2 \textbullet\ #3}{popink}{popcream}\hspace{5pt}\poppill{#4}{white}{popink}\par\vspace{8pt}
    {\popdisplay\fontsize{14}{14}\selectfont\color{popink}LEÇON #5}\par\vspace{4pt}
    {\raggedright\hyphenpenalty=10000\exhyphenpenalty=10000\popdisplay\fontsize{25}{27}\selectfont\color{popcream}\MakeUppercase{#1}\par}
    \vspace{8pt}
    {\popxbold\fontsize{9.5}{9.5}\selectfont\color{popink}\MakeUppercase{Durée conseillée : #6}}
  \end{tcolorbox}
  \begin{tcolorbox}[enhanced,colback=white,colframe=popink,boxrule=2.2pt,arc=10pt,
      drop shadow=popink,left=16pt,right=16pt,top=10pt,bottom=10pt,after skip=8pt]
    \poppill{Dans cette leçon}{poppurple}{white}\par\vspace{5pt}
    \popsemi\fontsize{9.3}{12.6}\selectfont\color{popink}#7
  \end{tcolorbox}
  \noindent\begin{minipage}[t]{0.487\linewidth}
    \begin{tcolorbox}[enhanced,colback=popgreen,colframe=popink,boxrule=2.2pt,arc=10pt,
        drop shadow=popink,left=11pt,right=11pt,top=10pt,bottom=10pt,height=3.1cm,valign=center]
      \tcbox[colback=white,colframe=popink,boxrule=1.3pt,arc=5pt,boxsep=0pt,nobeforeafter,
        left=3pt,right=3pt,top=3pt,bottom=3pt,valign=center]{\qrcode[height=1.9cm]{https://play.google.com/store/apps/details?id=com.luvvix.onbuch&hl=fr}}%
      \hspace{8pt}\begin{minipage}[c]{0.5\linewidth}
        {\popdisplay\fontsize{11}{12.5}\selectfont\color{white}\MakeUppercase{Télécharge OnBuch}}\par\vspace{4pt}
        {\raggedright\popsemi\fontsize{8.4}{11}\selectfont\color{white}Gratuit \textbullet\ Google Play\\Tuteur IA, annales, concours, résultats\par}
      \end{minipage}
    \end{tcolorbox}
  \end{minipage}\hfill
  \begin{minipage}[t]{0.487\linewidth}
    \begin{tcolorbox}[enhanced,colback=poppurple,colframe=popink,boxrule=2.2pt,arc=10pt,
        drop shadow=popink,left=11pt,right=11pt,top=10pt,bottom=10pt,height=3.1cm,valign=center]
      \tikz[baseline=-0.7ex]\node[circle,fill=white,draw=popink,line width=1.3pt,minimum size=1.6cm,inner sep=0]{\popdisplay\fontsize{24}{24}\selectfont\color{poppurple}+};%
      \hspace{8pt}\begin{minipage}[c]{0.6\linewidth}
        {\popdisplay\fontsize{11}{12.5}\selectfont\color{white}\MakeUppercase{Passe à OnBuch+}}\par\vspace{4pt}
        {\raggedright\popsemi\fontsize{8.4}{11}\selectfont\color{white}Tous les cours signés, Tuteur IA illimité, fiches PDF — onglet OnBuch+ de l'app\par}
      \end{minipage}
    \end{tcolorbox}
  \end{minipage}
  \vspace{8pt}
  \popcontenu
  \vfill
  \signatureonbuch
  \restoregeometry
  \newpage
}

% Rangée « Ce cours contient » (page de garde)
\newcommand{\poptile}[3]{% #1 couleur #2 gros texte #3 légende
  \begin{tcolorbox}[enhanced,colback=#1,colframe=popink,boxrule=2pt,arc=9pt,shadow={2.5pt}{-2.5pt}{0pt}{popink},
      left=6pt,right=6pt,top=9pt,bottom=9pt,halign=center,valign=center,height=1.75cm,width=\linewidth,nobeforeafter]
    {\popdisplay\fontsize{12.5}{13}\selectfont\color{white}\MakeUppercase{#2}}\par\vspace{3pt}
    {\centering\popsemi\fontsize{7.8}{9.5}\selectfont\color{white}#3\par}
  \end{tcolorbox}}
\newcommand{\popcontenu}{%
  \par\noindent{\popxboldls\fontsize{7.5}{7.5}\selectfont\color{popmuted}CE COURS SIGNÉ CONTIENT}\par\vspace{7pt}
  \noindent\begin{minipage}[t]{0.235\linewidth}\poptile{popblue}{Cours}{complet, illustré, pas à pas}\end{minipage}\hfill
  \begin{minipage}[t]{0.235\linewidth}\poptile{poporange}{Méthodes}{et exercices résolus}\end{minipage}\hfill
  \begin{minipage}[t]{0.235\linewidth}\poptile{poppink}{Exercices}{type examen + corrigés}\end{minipage}\hfill
  \begin{minipage}[t]{0.235\linewidth}\poptile{popgreen}{Fiche bilan}{l'essentiel à retenir}\end{minipage}\par}

% Sommaire
\newtcolorbox{sommaire}{enhanced,colback=white,colframe=popink,boxrule=2.2pt,arc=10pt,
  drop shadow=popink,left=18pt,right=18pt,top=13pt,bottom=13pt,before skip=6pt,after skip=20pt,
  before upper={\poppill{Sommaire}{poppurple}{white}\par\vspace{9pt}\popsemi\fontsize{10.6}{14.5}\selectfont\color{popink}}}

% Signature de fin de cours — poussée tout en bas de la dernière page
\newcommand{\findecours}{%
  \par\vfill\nopagebreak
  \begin{center}\popsquiggle{poporange}\end{center}\vspace{4pt}
  \signatureonbuch
}
\AtBeginDocument{\def\llbracket{\mathopen{[\mkern-3.5mu[}}\def\rrbracket{\mathclose{]\mkern-3.5mu]}}} % intervalles d'entiers (glyphe ⟦ absent de la police)
% Glyphes absents de Fira Math : redéfinis à partir de glyphes disponibles
\AtBeginDocument{%
  \def\setminus{\mathbin{\backslash}}\let\smallsetminus\setminus
  \def\vdots{\mathord{\vbox{\baselineskip3.2pt\lineskiplimit0pt\kern4pt\hbox{.}\hbox{.}\hbox{.}}}}%
  \def\ddots{\mathinner{\mkern1mu\raise6.5pt\hbox{.}\mkern2mu\raise3.6pt\hbox{.}\mkern2mu\raise0.7pt\hbox{.}\mkern1mu}}%
  \def\triangle{\mathord{\scalebox{0.9}{$\symup{\Delta}$}}}%
  \def\nabla{\mathord{\raisebox{\depth}{\scalebox{1}[-1]{$\symup{\Delta}$}}}}%
  \def\checkmark{\popcheck{popdark}}%
  \def\star{\mathbin{\ast}}\def\bigstar{\mathord{\ast}}%
  \def\diamond{\mathbin{\scalebox{0.6}{$\Diamond$}}}%
  \def\bigcup{\mathop{\mathchoice{\raisebox{-0.3em}{\scalebox{1.7}{$\cup$}}}{\scalebox{1.25}{$\cup$}}{\cup}{\cup}}\displaylimits}%
  \def\bigcap{\mathop{\mathchoice{\raisebox{-0.3em}{\scalebox{1.7}{$\cap$}}}{\scalebox{1.25}{$\cap$}}{\cap}{\cap}}\displaylimits}%
  \def\lesseqqgtr{\mathrel{\lessgtr}}\def\bigoplus{\mathop{\oplus}\displaylimits}%
}
\providecommand{\ssi}{si et seulement si} % abréviation fréquente
\AtBeginDocument{\providecommand{\wideparen}{\overparen}} % arc de cercle

%% FIGFIX-BEGIN v4 — correctifs globaux des figures (docs/onbuch-generation/tools/figfix)
\usepackage{adjustbox}\usepackage{etoolbox}
% toute figure TikZ/pgfplots plus large que la ligne est réduite à la largeur disponible
\BeforeBeginEnvironment{tikzpicture}{\begin{adjustbox}{max width=\linewidth}}
\AfterEndEnvironment{tikzpicture}{\end{adjustbox}}
% légendes pgfplots lisibles par-dessus les courbes
\pgfplotsset{every axis/.append style={legend style={fill=white,fill opacity=0.92,text opacity=1,draw=black!15,inner sep=3pt}}}
% étiquettes d'arêtes (syntaxe edge["..."]) avec fond blanc
\tikzset{every edge quotes/.append style={fill=white,inner sep=1.2pt}}
%% FIGFIX-END
\begin{document}

\pagegarde{Lesson 36 — Writing: Writes compositions on garbage collection, disposal, and recycling services}{Anglais}{Tle A}{Chapter 3 : Environment, Well-being and Health}{EN36}{2 h}{Cette leçon d'expression écrite apprend aux élèves de Terminale A à planifier et rédiger une composition structurée (introduction, développement, conclusion) sur la collecte, l'élimination et le recyclage des déchets. Elle fournit le vocabulaire thématique, les connecteurs logiques, un modèle commenté, les erreurs à éviter et une grille d'évaluation conforme aux attentes du Bac.
\vspace{4pt}
\begin{multicols}{2}\small
\begin{itemize}
\item À la fin de cette leçon, je sais utiliser le vocabulaire thématique de la collecte, de l'élimination et du recyclage des déchets
\item À la fin de cette leçon, je sais analyser un sujet de composition et dégager la problématique
\item À la fin de cette leçon, je sais construire un plan en trois parties : introduction, développement, conclusion
\item À la fin de cette leçon, je sais rédiger une introduction avec amorce, problématique et annonce de plan
\item À la fin de cette leçon, je sais organiser mes arguments en paragraphes reliés par des connecteurs logiques
\item À la fin de cette leçon, je sais employer les temps verbaux appropriés (présent simple, présent continu, futur, passif)
\end{itemize}\end{multicols}}

\begin{sommaire}
\begin{multicols}{2}
\begin{enumerate}
\item Vocabulaire thématique : garbage, disposal, recycling
\item Analyser le sujet et planifier
\item La structure de la composition
\item Connecteurs logiques et temps verbaux
\item Modèle commenté
\item Erreurs à éviter et grille d'évaluation
\item Activité d'intégration
\item Méthodes et exercices résolus
\item Exercices
\item Corrigés des exercices
\item Fiche bilan
\end{enumerate}
\end{multicols}
\end{sommaire}

\begin{objectifs}
\begin{itemize}
\item À la fin de cette leçon, je sais utiliser le vocabulaire thématique de la collecte, de l'élimination et du recyclage des déchets
\item À la fin de cette leçon, je sais analyser un sujet de composition et dégager la problématique
\item À la fin de cette leçon, je sais construire un plan en trois parties : introduction, développement, conclusion
\item À la fin de cette leçon, je sais rédiger une introduction avec amorce, problématique et annonce de plan
\item À la fin de cette leçon, je sais organiser mes arguments en paragraphes reliés par des connecteurs logiques
\item À la fin de cette leçon, je sais employer les temps verbaux appropriés (présent simple, présent continu, futur, passif)
\item À la fin de cette leçon, je sais rédiger une conclusion qui synthétise et ouvre sur une recommandation
\item À la fin de cette leçon, je sais relire et corriger ma production à l'aide d'une grille d'évaluation
\end{itemize}
\end{objectifs}

\begin{prerequis}
\begin{itemize}
\item Vocabulaire de base de l'environnement et de la santé (pollution, disease, waste) vu en Première
\item Les connecteurs logiques simples (and, but, because, so, however)
\item Le présent simple et le présent continu pour décrire des faits et des habitudes
\item La voix passive (is collected, are recycled) vue en Première
\item La rédaction d'un paragraphe argumentatif simple
\end{itemize}
\end{prerequis}

\newpage

\coursec{Vocabulaire thématique : garbage, disposal, recycling}

Every morning in Yaoundé, Douala or Bamenda, heaps of uncollected \eng{garbage} pile up in the streets, blocking gutters and threatening public health. Yet in London or Accra, recycling trucks collect sorted waste door to door and turn it into new products. Imagine your school wants to write to the city council to demand better garbage collection services: could you convince them in a well-organised English composition? This lesson gives you the tools to plan, write and polish a composition on garbage collection, disposal and recycling — a classic topic at the Baccalauréat.

\courssub{Champ lexical de la collecte (Collection)}

Observons d'abord un texte court qui met en scène le vocabulaire de la collecte.

\begin{textebox}[Situation de collecte à Yaoundé]
Every day, garbage collectors pick up tons of rubbish in our streets. Unfortunately, much of this waste is dumped in open landfills instead of being sorted and recycled. The municipal truck passes only once a week, so residents often burn their trash or throw it into streams.
\end{textebox}

\begin{definition}[Vocabulaire clé : collecte]
\begin{itemize}
\item \cle{garbage} (US) / \cle{rubbish} (UK) / \cle{waste} (formel) : ordures, déchets (indénombrables).
\item \cle{dustbin} (UK) / \cle{trash can} (US) / \cle{bin} : poubelle.
\item \cle{garbage truck} / \cle{refuse truck} : camion-poubelle.
\item \cle{dustman} (UK) / \cle{garbage collector} / \cle{refuse collector} : éboueur.
\item \cle{to collect} / \cle{to pick up} : collecter, ramasser.
\item \cle{to pick up litter} : ramasser les détritus (jetés au sol).
\item \cle{door-to-door collection} : collecte porte-à-porte.
\end{itemize}
\end{definition}

\begin{center}
\begin{tabularx}{\linewidth}{|X|X|X|}
\hline
\rowcolor{popblueL} \textbf{Français} & \textbf{English (UK)} & \textbf{English (US)} \\ \hline
poubelle & dustbin, bin & trash can, garbage can \\ \hline
ordures (général) & rubbish & garbage, trash \\ \hline
déchets (formel) & waste & waste \\ \hline
éboueur & dustman, refuse collector & garbage collector, sanitation worker \\ \hline
camion-poubelle & dustcart, refuse truck & garbage truck \\ \hline
ramasser (déchets) & collect, pick up & collect, pick up \\ \hline
\end{tabularx}
\end{center}

\eng{The city council collects garbage twice a week.} = La mairie collecte les ordures deux fois par semaine. \\
\eng{Residents must put their bins out on Monday evening.} = Les habitants doivent sortir leurs poubelles le lundi soir. \\
\eng{Don't litter! Use the dustbin.} = Ne jetez pas de détritus ! Utilisez la poubelle.

\begin{attention}[Faux ami et piège d'orthographe]
\begin{itemize}
\item \cle{garbage}, \cle{rubbish}, \cle{waste}, \cle{litter} sont \textbf{indénombrables}. On ne dit jamais \eng{garbages} ni \eng{rubbishes}. Pour compter, on utilise \eng{a piece of rubbish}, \eng{items of waste}, \eng{heaps of garbage}.
\item \cle{litter} (comme nom) = détritus jetés dans la rue ; comme verbe \eng{to litter} = jeter des détritus. \eng{He was fined for littering.} = Il a été amendé pour avoir jeté des déchets.
\end{itemize}
\end{attention}

\begin{popfigure}
\begin{tikzpicture}[
  mindmap,
  concept color=popblue,
  level 1/.style={sibling angle=120, level distance=4.5cm},
  level 2/.style={sibling angle=60, level distance=3.5cm},
  every node/.style={concept, execute at begin node=\strut},
  text=white,
  font=\small
]
\node {WASTE MANAGEMENT}
  child[concept color=popgreen] { node {Collection}
    child { node {garbage / rubbish / waste} }
    child { node {dustbin / trash can} }
    child { node {garbage truck} }
    child { node {garbage collector} }
    child { node {door-to-door collection} }
    child { node {pick up litter} }
  }
  child[concept color=poporange] { node {Disposal}
    child { node {dispose of} }
    child { node {landfill / dump} }
    child { node {incinerator} }
    child { node {open dumping} }
    child { node {hazardous waste} }
    child { node {sewage} }
  }
  child[concept color=poppurple] { node {Recycling}
    child { node {recycle / recycling} }
    child { node {sort / separate waste} }
    child { node {reusable / biodegradable} }
    child { node {compost} }
    child { node {plastic, glass, metal} }
    child { node {reduce, reuse, recycle} }
  };
\end{tikzpicture}
\legende{Carte mentale : vocabulaire de la gestion des déchets (Collection, Disposal, Recycling)}
\end{popfigure}

\courssub{Champ lexical de l'élimination (Disposal)}

\begin{textebox}[Élimination des déchets]
After collection, waste is transported to disposal sites. In many cities, open dumping is still common: trucks unload garbage in huge landfills where it rots and produces methane. Modern incinerators burn waste at high temperatures, reducing volume but requiring strict emission controls. Hazardous waste like batteries or chemicals needs special treatment. Sewage, meanwhile, goes to treatment plants before being released into rivers.
\end{textebox}

\begin{definition}[Vocabulaire clé : élimination]
\begin{itemize}
\item \cle{to dispose of} : éliminer, se débarrasser de (suivi de \eng{of}). \eng{We must dispose of waste properly.}
\item \cle{disposal} (nom) : élimination. \eng{Waste disposal is a major challenge.}
\item \cle{dump} / \cle{landfill} : décharge (à ciel ouvert / enfouissement technique).
\item \cle{to dump} : déverser, jeter (souvent illégalement).
\item \cle{incinerator} : incinérateur.
\item \cle{to burn} : brûler.
\item \cle{open dumping} : décharge sauvage.
\item \cle{sewage} : eaux usées, égouts.
\item \cle{hazardous waste} : déchets dangereux.
\end{itemize}
\end{definition}

\begin{attention}[Faux ami majeur : \textbf{to dispose of} $\neq$ \textbf{disposer de}]
\begin{itemize}
\item \eng{to dispose of} = éliminer, jeter. \eng{How do you dispose of old batteries?}
\item \eng{to dispose of} (autre sens) = régler, traiter (un problème). \eng{The court disposed of the case quickly.}
\item \eng{to have at one's disposal} = avoir à sa disposition. \eng{We have a large budget at our disposal.}
\item \textbf{Ne confondez jamais} \eng{dispose of} (éliminer) et \eng{have at disposal} (avoir à disposition).
\end{itemize}
\end{attention}

\begin{exemplebox}[Exemples traduits : élimination]
\eng{The municipality disposes of solid waste in a landfill.} = La municipalité élimine les déchets solides dans une décharge. \\
\eng{Illegal dumping pollutes groundwater.} = Le déversement illégal pollute les nappes phréatiques. \\
\eng{Incinerators reduce waste volume by 90\%.} = Les incinérateurs réduisent le volume des déchets de 90\%. \\
\eng{Hazardous waste must not be mixed with household garbage.} = Les déchets dangereux ne doivent pas être mélangés aux ordures ménagères.
\end{exemplebox}

\courssub{Champ lexical du recyclage (Recycling)}

\begin{definition}[Recycling]
\cle{Recycling} (nom) : the process of treating used materials so that they can be used again. \\
\eng{Plastic bottles are recycled into new products.} = Les bouteilles en plastique sont recyclées en nouveaux produits.
\end{definition}

\begin{definition}[Vocabulaire clé : recyclage]
\begin{itemize}
\item \cle{to recycle} : recycler.
\item \cle{recycling} (nom) : recyclage.
\item \cle{to sort waste} / \cle{to separate waste} : trier les déchets.
\item \cle{reusable} : réutilisable.
\item \cle{biodegradable} : biodégradable.
\item \cle{compost} : compost ; \cle{to compost} : composter.
\item \cle{plastic}, \cle{glass}, \cle{metal}, \cle{paper}, \cle{cardboard} : plastique, verre, métal, papier, carton.
\item \cle{to reuse} : réutiliser.
\item \cle{to reduce} : réduire (à la source).
\item \cle{the three Rs} : Reduce, Reuse, Recycle.
\end{itemize}
\end{definition}

\begin{savaistu}[Royaume-Uni vs États-Unis : cohérence exigée]
Au Royaume-Uni, on dit \eng{rubbish} (ordures), \eng{bin} (poubelle), \eng{dustman} (éboueur). Aux États-Unis, on dit \eng{garbage} / \eng{trash}, \eng{trash can}, \eng{garbage collector}. Au Baccalauréat, les deux variétés sont acceptées \textbf{à condition de rester cohérent} dans toute la composition. Ne mélangez pas \eng{bin} et \eng{trash can} dans le même texte.
\end{savaistu}

\begin{center}
\begin{tabularx}{\linewidth}{|X|X|X|}
\hline
\rowcolor{popgreenL} \textbf{Action} & \textbf{Verbe} & \textbf{Exemple} \\ \hline
Trier & to sort, to separate & \eng{Sort your waste into three bins.} \\ \hline
Recycler & to recycle & \eng{Glass is 100\% recyclable.} \\ \hline
Réutiliser & to reuse & \eng{Reuse plastic bags instead of throwing them away.} \\ \hline
Réduire & to reduce & \eng{Reduce packaging by buying in bulk.} \\ \hline
Composter & to compost & \eng{Compost kitchen scraps for the garden.} \\ \hline
\end{tabularx}
\end{center}

\courssub{Collocations et expressions utiles}

Les collocations (associations de mots naturelles) donnent de la fluidité et de la précision à votre écriture.

\begin{center}
\begin{tabularx}{\linewidth}{|X|X|}
\hline
\rowcolor{poporangeL} \textbf{Collocation} & \textbf{Traduction} \\ \hline
to collect garbage / rubbish / waste & collecter les ordures \\ \hline
to dispose of waste properly & éliminer les déchets correctement \\ \hline
to sort / separate waste & trier les déchets \\ \hline
to reduce pollution / waste & réduire la pollution / les déchets \\ \hline
to protect the environment & protéger l'environnement \\ \hline
to raise awareness & sensibiliser \\ \hline
public health & santé publique \\ \hline
waste management & gestion des déchets \\ \hline
recycling plant / centre & centre de tri / de recyclage \\ \hline
landfill site & site d'enfouissement \\ \hline
\end{tabularx}
\end{center}

\begin{exemplebox}[Expressions idiomatiques et slogans]
\eng{A clean city is a healthy city.} = Une ville propre est une ville en bonne santé. \\
\eng{Waste not, want not.} = Qui ne gaspille pas ne manque de rien (proverbe). \\
\eng{The three Rs: Reduce, Reuse, Recycle.} = Les trois R : Réduire, Réutiliser, Recycler.
\end{exemplebox}

\courssub{Exercice résolu : identifier et corriger les erreurs}

\begin{exoresolu}[Correction d'erreurs fréquentes]
Corrigez les phrases suivantes (une ou plusieurs erreurs par phrase).
\begin{enumerate}
\item The garbages are collected every morning.
\item We must dispose the waste in the landfill.
\item The municipality has many dustbins at their disposal for the population.
\item People should recycle their garbages to protect the environment.
\item Open dumping is a serious problem in ours cities.
\end{enumerate}
\tcblower
\textbf{Solution détaillée :}
\begin{enumerate}
\item \eng{The garbage is collected every morning.} — \eng{garbage} est indénombrable $\rightarrow$ singulier (\eng{is}), pas de \eng{-s}.
\item \eng{We must dispose \textbf{of} the waste in the landfill.} — \eng{dispose} s'emploie avec \eng{of} (éliminer).
\item \eng{The municipality has many dustbins \textbf{at its disposal} for the population.} — \eng{at one's disposal} = à sa disposition ; \eng{their} $\rightarrow$ \eng{its} (municipality = singulier).
\item \eng{People should recycle \textbf{their waste} to protect the environment.} — \eng{garbages} n'existe pas ; on utilise \eng{waste} ou \eng{rubbish}.
\item \eng{Open dumping is a serious problem in \textbf{our} cities.} — \eng{ours} est pronom possessif (mine, yours, \textbf{ours}) ; adjectif possessif = \eng{our}.
\end{enumerate}
\end{exoresolu}

\courssub{Schéma du cycle de gestion des déchets}

\begin{popfigure}
\begin{tikzpicture}[
  node distance=1.2cm and 1.5cm,
  box/.style={draw=popblue, thick, rounded corners, fill=popblueL, minimum width=2.2cm, minimum height=1cm, align=center, font=\small},
  arrow/.style={-Stealth, thick, popdark}
]
\node[box, align=center] (household){Household\\Ménage};
\node[box, right=of household, align=center] (dustbin){Dustbin /\\Trash can};
\node[box, right=of dustbin, align=center] (truck){Garbage truck\\Camion-poubelle};
\node[box, below right=of truck, align=center] (landfill){Landfill\\Décharge};
\node[box, above right=of truck, align=center] (plant){Recycling plant\\Centre de tri};
\node[box, right=of plant, align=center] (new){New products\\Nouveaux produits};

\draw[arrow] (household) -- (dustbin);
\draw[arrow] (dustbin) -- (truck);
\draw[arrow] (truck) -- node[above, sloped, font=\tiny, text=popdark] {unsorted} (landfill);
\draw[arrow] (truck) -- node[above, sloped, font=\tiny, text=popdark] {sorted} (plant);
\draw[arrow] (plant) -- (new);
\end{tikzpicture}
\legende{Cycle simplifié : du ménage à la décharge ou au recyclage}
\end{popfigure}

\begin{methode}[Mémoriser le vocabulaire par champs]
\begin{enumerate}
\item Créez trois colonnes dans votre cahier : \textbf{Collection}, \textbf{Disposal}, \textbf{Recycling}.
\item Y rangez chaque mot nouveau avec sa traduction et une phrase exemple personnelle.
\item Révisez en couvrant la colonne anglaise et en retrouvant les mots à partir du français.
\item Utilisez au moins 5 mots de chaque colonne dans votre prochaine composition.
\end{enumerate}
\end{methode}

\coursec{Analyser le sujet et planifier}

Après avoir acquis le \cle{vocabulaire thématique} indispensable (garbage, waste, disposal, recycling, landfill, incineration, composting, collection, segregation, sustainability), nous abordons maintenant l'étape cruciale qui conditionne la réussite de toute composition : \textbf{l'analyse du sujet et la planification}. Trop d'élèves se lancent dans la rédaction sans avoir décortiqué la consigne, ce qui entraîne des hors-sujet, des répétitions ou une structure confuse. Une composition bien planifiée, c'est déjà la moitié du travail accompli.

\courssub{Lire le sujet : identifier les mots-clés et la tâche}

La première règle d'or est de \textbf{lire le sujet au moins deux fois}. Lors de la première lecture, on saisit le thème global. Lors de la seconde, on \cle{souligne les mots-clés} qui dictent le contenu, la forme et les contraintes.

\begin{textebox}[Sujet type d'examen (Bac Blanc)]
Write a composition of about 250–300 words on the following topic:
\textbf{“Garbage collection and disposal in our cities remain a major challenge. Discuss the causes of this problem and suggest practical solutions to improve the situation.”}
\end{textebox}

\begin{methode}[Analyser un sujet de composition en 3 temps]
\begin{enumerate}
    \item \textbf{Souligner les mots-clés} (thème, verbes d'action, contraintes de longueur).
    \item \textbf{Identifier le type de sujet} (description, argumentation, proposition de solutions, narration).
    \item \textbf{Transformer le sujet en question problématique} pour guider la réflexion.
\end{enumerate}
\end{methode}

Appliquons cette méthode au sujet ci-dessus :

\begin{center}
\begin{tabularx}{\linewidth}{|X|X|X|}
\hline
\rowcolor{popblueL} \textbf{Mot-clé / Segment} & \textbf{Rôle} & \textbf{Conséquence pour la rédaction} \\
\hline
\textbf{Garbage collection and disposal} & \cle{Thème central} & Le vocabulaire appris (landfill, trucks, bins, sorting) est obligatoire. \\
\hline
\textbf{in our cities} & \cle{Contexte spatial} & Ancrer les exemples au Cameroun (Yaoundé, Douala, Buea, Bamenda). \\
\hline
\textbf{remain a major challenge} & \cle{Constat / Problème} & L'introduction doit présenter ce constat (present perfect / present simple). \\
\hline
\textbf{Discuss the causes} & \cle{Tâche 1 (Analyse)} & Paragraphes du corps 1 (et 2) : expliquer \emph{why} (population growth, lack of bins, bad habits, funding). \\
\hline
\textbf{suggest practical solutions} & \cle{Tâche 2 (Proposition)} & Paragraphes du corps 3 (et 4) : proposer \emph{how} (recycling plants, education, fines, PPP). \\
\hline
\textbf{250–300 words} & \cle{Contrainte de longueur} & Ni trop court (sanction), ni trop long (perte de temps, risque d'erreurs). \\
\hline
\end{tabularx}
\end{center}

\begin{attention}[Piège fréquent : ignorer une partie du sujet]
Si vous ne traitez que les \textbf{causes} et oubliez les \textbf{solutions} (ou l'inverse), vous êtes \textbf{hors-sujet partiel}. La note de \emph{Content/Task Achievement} s'effondre (souvent divisée par 2). \textbf{Cochez chaque mot-clé de la consigne sur votre brouillon} au fur et à mesure que vous le traitez dans votre plan.
\end{attention}

\courssub{Identifier le type de sujet et formuler la problématique}

Tous les sujets de composition ne se valent pas. En Terminale, on rencontre trois grands types. Les reconnaître oriente instantanément la structure du plan.

\begin{definition}[Les 3 types de sujets de composition (Writing)]
\begin{itemize}
    \item \textbf{Type A — Description / Exposition :} \eng{Describe how garbage is collected in your neighbourhood.} $\rightarrow$ Plan chronologique ou thématique (étapes, acteurs, matériel).
    \item \textbf{Type B — Argumentation (Pour / Contre) :} \eng{Should the government ban plastic bags completely? Give reasons.} $\rightarrow$ Plan dialectique (Thèse / Antithèse / Synthèse) ou arguments groupés.
    \item \textbf{Type C — Problème / Solutions (Problem-Solution) :} \emph{Notre sujet type.} $\rightarrow$ Plan : Introduction (problème) $\rightarrow$ Causes $\rightarrow$ Solutions $\rightarrow$ Conclusion.
\end{itemize}
\end{definition}

Une fois le type identifié (ici Type C), on formule la \cle{problématique} : une question unique qui résume la tension du sujet et que la composition va répondre.

\begin{exemplebox}[Du sujet à la problématique]
\begin{itemize}
    \item \textbf{Sujet :} \eng{Garbage collection and disposal in our cities remain a major challenge. Discuss the causes of this problem and suggest practical solutions to improve the situation.}
    \item \textbf{Problématique :} \eng{How can our cities effectively tackle the causes of poor waste management to ensure a cleaner environment?} (« Comment nos villes peuvent-elles s'attaquer efficacement aux causes de la mauvaise gestion des déchets pour assurer un environnement plus propre ? »)
\end{itemize}
\end{exemplebox}

\begin{aretenir}[La problématique : votre boussole]
La problématique ne s'écrit \textbf{pas} dans la copie finale (sauf si le sujet demande explicitement \eng{Write an essay answering the following question}). Elle reste sur votre \textbf{brouillon}. Elle vous oblige à ne pas raconter n'importe quoi : chaque paragraphe doit apporter un élément de réponse à cette question.
\end{aretenir}

\courssub{Le Brainstorming (Remue-méninges) : 5 minutes chrono}

Ne cherchez pas la phrase parfaite. Notez \textbf{tout} ce qui vous passe par la tête : idées, mots de vocabulaire, exemples locaux, connecteurs, arguments contraires. C'est le « déversoir » cognitif.

\begin{experience}[Brainstorming guidé : « Garbage in Yaoundé »]
Sur votre brouillon, faites deux colonnes : \textbf{CAUSES} | \textbf{SOLUTIONS}. Remplissez en 5 minutes.

\begin{center}
\begin{tabularx}{\linewidth}{|X|X|}
\hline
\rowcolor{popgreenL} \textbf{CAUSES (Why?)} & \textbf{SOLUTIONS (How?)} \\
\hline
Rapid urbanisation / population growth & Build modern landfills / recycling plants \\
Lack of enough trash bins (dustbins) & Provide more bins + regular collection \\
Irregular collection trucks (breakdown, fuel) & Public-Private Partnerships (PPP) for fleet \\
Poor road access in neighbourhoods & Community-based collection (tricycles) \\
Citizens' bad habits (littering, burning) & Sensitisation campaigns in schools/media \\
No sorting at source (mixed waste) & Mandatory segregation (organic / plastic) \\
Corruption / mismanagement of funds & Transparency, audits, citizen monitoring \\
\hline
\end{tabularx}
\end{center}
\end{experience}

\begin{attention}[Erreur de francophone : écrire le brouillon en français]
Si vous faites votre brainstorming et votre plan en français, vous perdrez un temps précieux à \textbf{traduire} pendant la rédaction, ce qui génère des erreurs de structure (calques syntaxiques). \textbf{Pensez et notez directement en anglais} (mots-clés, chunks : \eng{rapid urbanisation}, \eng{irregular collection}, \eng{sorting at source}).
\end{attention}

\courssub{Sélectionner et classer : de la masse au plan (Outline)}

Vous avez 10-15 idées. Impossible de tout dire en 250 mots. Il faut \textbf{sélectionner les 3 ou 4 idées les plus fortes} (celles que vous pouvez développer avec du vocabulaire précis et des exemples concrets) et les \textbf{organiser logiquement}.

Pour un sujet \textbf{Problème / Solutions}, la structure standard (et la plus sûre au Bac) est :

\begin{propriete}[Structure type « Problem – Solution » (4 paragraphes corps)]
\begin{tabular}{|p{2.5cm}|p{6cm}|p{5cm}|}
\hline
\rowcolor{popblueL} \textbf{Partie} & \textbf{Contenu} & \textbf{Exemple d'idée directrice (Topic Sentence)} \\
\hline
\textbf{Intro} ($\approx$ 40 mots) & Accroche + Constat + Problématique + Annonce du plan & \eng{Urban waste management is a burning issue in Cameroonian cities like Yaoundé. This essay examines the root causes and proposes concrete remedies.} \\
\hline
\textbf{Body 1 : Causes} ($\approx$ 80 mots) & 2 causes principales + exemples / explications & \eng{The primary causes are rapid urbanisation outpacing infrastructure and citizens' poor disposal habits.} \\
\hline
\textbf{Body 2 : Causes (suite) ou Transition} ($\approx$ 60 mots) & Cause structurelle (financement, voirie) ou transition vers solutions & \eng{Moreover, insufficient funding and poor road networks hinder regular collection services.} \\
\hline
\textbf{Body 3 : Solutions} ($\approx$ 80 mots) & 2 solutions concrètes + modalités & \eng{To address this, authorities must invest in recycling plants and enforce waste sorting at source.} \\
\hline
\textbf{Body 4 : Solutions (suite)} ($\approx$ 60 mots) & Solution éducative / citoyenne + exemple local & \eng{Equally important, civic education and community clean-up days can change mindsets.} \\
\hline
\textbf{Conclusion} ($\approx$ 40 mots) & Résumé + Ouverture / Appel à l'action & \eng{In conclusion, solving the waste crisis requires both state investment and citizen responsibility. A clean city is everyone's duty.} \\
\hline
\end{tabular}
\end{propriete}

\begin{savaistu}[Au Bac : l'organisation rapporte gros]
La grille officielle d'évaluation du Baccalauréat (MINESEC) réserve souvent \textbf{4 à 6 points sur 20} à l'\emph{Organization / Coherence and Cohesion} (paragraphing, linking words, logical progression). Une composition avec un anglais moyen mais un \textbf{plan parfait et visible} (paragraphes sautés, connecteurs) obtiendra une meilleure note qu'une composition en anglais parfait mais sans structure apparente. \textbf{Le plan, c'est des points gratuits.}
\end{savaistu}

\courssub{Rédiger le plan écrit (Outline) sur le brouillon}

Ne gardez pas le plan dans la tête. \textbf{Écrivez-le} sur la marge gauche de votre brouillon. C'est votre contrat avec vous-même. Utilisez des abréviations, des mots-clés anglais, des flèches.

\begin{exemplebox}[Modèle de plan écrit (Outline) — Brouillon]
\textbf{Topic:} Garbage collection challenges \& solutions.
\textbf{Thesis/Problem:} Major challenge in cities (Yaoundé/Douala).
\textbf{Plan:}
\begin{enumerate}
    \item \textbf{Intro:} Hook (visual: piles of trash), Context, Thesis: Causes = urban growth + habits; Solutions = infrastructure + education.
    \item \textbf{Body 1 (Causes):} 1. Rapid urbanisation $\rightarrow$ infrastructure gap. Ex: Yaoundé pop. doubled, trucks same. 2. Bad habits: littering, burning plastic. Connectors: \eng{Firstly, Furthermore}.
    \item \textbf{Body 2 (Causes - Structural):} 3. Funding / Logistics. Trucks broken, fuel scarce. Roads bad in neighbourhoods. Connector: \eng{Moreover}.
    \item \textbf{Body 3 (Solutions - Institutional):} 1. Govt/PPP: Modern landfills, recycling plants (plastic $\rightarrow$ pellets). 2. Policy: Mandatory sorting (organic/plastic), fines. Connectors: \eng{To solve this, Consequently}.
    \item \textbf{Body 4 (Solutions - Civic):} 3. Education: Schools, media (CRTV, social media). 4. Community action: \eng{Clean-up Saturdays}, youth groups. Connector: \eng{Additionally, In the long run}.
    \item \textbf{Conclusion:} Sum up (Infrastructure + Behaviour). Final thought: \eng{Shared responsibility}. Call to action.
\end{enumerate}
\textbf{Word count check:} $40 + 80 + 60 + 80 + 60 + 40 \approx 360$ (trop long $\rightarrow$ couper exemples dans B2/B4). \textbf{Target: 280 words.}
\end{exemplebox}

\begin{methode}[Check-list avant de rédiger (30 secondes)]
\begin{enumerate}
    \item Ai-je répondu à \textbf{TOUTES} les parties de la consigne ? (Causes \checkmark / Solutions \checkmark)
    \item Mon plan a-t-il \textbf{4 paragraphes de corps} distincts ? (Pas un gros bloc).
    \item Ai-je noté les \textbf{connecteurs logiques} pour chaque transition ?
    \item Ai-je prévu du \textbf{vocabulaire précis} (landfill, segregation, PPP, littering) dans chaque paragraphe ?
    \item Le compte-mots estimé est-il dans la fourchette \textbf{250–300} ?
\end{enumerate}
\end{methode}

\courssub{Illustration : Le processus complet de la planification}

\begin{popfigure}
\begin{tikzpicture}[
    node distance=1.2cm and 1.5cm,
    box/.style={draw=popblue, thick, rounded corners, fill=popblueL, text width=2.8cm, align=center, minimum height=1.6cm, font=\small},
    arrow/.style={-Stealth, thick, popdark}
]
\node[box, align=center] (sujet){\textbf{SUJET}\\\eng{Write a composition...}};
\node[box, below=of sujet, align=center] (brain){\textbf{BRAINSTORMING}\\\textit{5 min}\\Causes / Solutions / Vocab / Exemples locaux};
\node[box, below=of brain, align=center] (select){\textbf{SÉLECTION}\\\textit{2 min}\\Choisir 3-4 idées fortes\\Grouper : Causes // Solutions};
\node[box, below=of select, align=center] (plan){\textbf{PLAN (OUTLINE)}\\\textit{3 min}\\Intro / Body 1-4 / Conclusion\\Topic sentences + Connecteurs};
\node[box, below=of plan, align=center] (draft){\textbf{RÉDACTION}\\\textit{45-50 min}\\Suivre le plan\\Compter les mots (\approx 280)};
\node[box, below=of draft, align=center] (proof){\textbf{RELECTURE}\\\textit{5-10 min}\\Grammaire / Orthographe\\Cohérence / Longueur};

\draw[arrow] (sujet) -- (brain);
\draw[arrow] (brain) -- (select);
\draw[arrow] (select) -- (plan);
\draw[arrow] (plan) -- (draft);
\draw[arrow] (draft) -- (proof);
\end{tikzpicture}
\legende{Le flux de travail idéal : du sujet à la relecture. Respecter les temps indicatifs est vital pour finir à l'heure.}
\end{popfigure}

\courssub{Exercice résolu : De l'analyse au plan}

\begin{exoresolu}[Analyser et planifier]
\textbf{Sujet :} \eng{“Many households in your area do not sort their waste. Write a composition explaining why this happens and proposing ways to encourage recycling.”} (250–300 words)

\textbf{Étape 1 : Mots-clés \& Tâche}
\begin{itemize}
    \item Thème : \eng{Household waste sorting / recycling} (Tri sélectif / recyclage ménager).
    \item Contexte : \eng{in your area} (votre quartier/ville $\rightarrow$ exemples camerounais).
    \item Constat : \eng{do not sort} (ne trient pas).
    \item Tâches : 1. \eng{Explaining why} (Causes). 2. \eng{Proposing ways to encourage} (Solutions/Incitation).
    \item Longueur : 250-300 mots.
\end{itemize}

\textbf{Étape 2 : Type \& Problématique}
\begin{itemize}
    \item Type : \textbf{Problem-Solution} (Pourquoi pas trier ? Comment inciter ?).
    \item Problématique : \eng{Why do households neglect waste sorting and how can authorities and communities foster a recycling culture?}
\end{itemize}

\textbf{Étape 3 : Brainstorming (Extrait)}
\begin{center}
\begin{tabular}{|l|l|}
\hline
\rowcolor{poporangeL} \textbf{WHY (Causes)} & \textbf{HOW TO ENCOURAGE (Solutions)} \\
\hline
No bins provided (only one bin) & Provide colour-coded bins (free/subsidised) \\
Ignorance: don't know how/why & Education: schools, TV, radio, chiefs \\
No recycling plant nearby (no outlet) & Build local recycling units (jobs!) \\
No incentive / no penalty & Incentives: tax rebate, cash for plastic \\
Habit / Laziness / \eng{Not my job} & Fines for non-compliance + peer pressure \\
\hline
\end{tabular}
\end{center}

\textbf{Étape 4 : Plan Final (Outline)}

\begin{tabularx}{\linewidth}{|X|}
\hline
\rowcolor{popgreenL} \textbf{PLAN DÉTAILLÉ} \\
\hline
\textbf{Intro (35 mots):} Hook: Mountains of mixed waste. Context: Sorting rare in [Ville]. Thesis: Causes = lack of infrastructure + awareness. Solutions = equipment + incentives + education. \\
\hline
\textbf{Body 1 - Cause 1: Infrastructure (75 mots):} Absence of separate bins. One truck takes all. Ex: \eng{In Mvog-Ada quarter, only one rusty bin exists.} Connector: \eng{Firstly, As a result}. \\
\hline
\textbf{Body 2 - Cause 2: Awareness \& Habit (65 mots):} People don't know \eng{biodegradable vs non-biodegradable}. Burning plastic = toxic. Cultural habit: \eng{throwing away}. Connector: \eng{Secondly, Furthermore}. \\
\hline
\textbf{Body 3 - Solution 1: Infrastructure \& Incentives (80 mots):} HYSACAM / Council provide 3 bins (organic, plastic, other). \eng{Deposit-return scheme} for plastic bottles (50 FCFA/bottle). Create local jobs. Connector: \eng{To remedy this, Consequently}. \\
\hline
\textbf{Body 4 - Solution 2: Education \& Law (70 mots):} School curriculum + media campaigns in \eng{Pidgin/French/English}. Bye-laws: fines for mixed waste. \eng{Quarter chiefs} enforce. Connector: \eng{Equally important, In the long term}. \\
\hline
\textbf{Conclusion (35 mots):} Sorting is a chain: citizen sorts $\rightarrow$ truck collects separately $\rightarrow$ factory recycles. Needs political will + civic duty. \eng{A cleaner Cameroon starts in my kitchen.} \\
\hline
\textbf{Total estimé :} $\approx$ 295 mots. \textbf{Parfait.}
\end{tabularx}
\tcblower
\textbf{Analyse de la solution :} Le plan respecte la consigne (2 causes, 2 solutions), ancre les exemples au Cameroun (\eng{Mvog-Ada}, \eng{HYSACAM} — \textit{Hygiène et Salubrité du Cameroun}, \eng{Quarter chiefs}, \eng{Pidgin}), prévoit les connecteurs et le vocabulaire cible (\eng{deposit-return scheme}, \eng{biodegradable}, \eng{bye-laws}). Le compte-mots est sécurisé.
\end{exoresolu}

\begin{exercice}[À vous de jouer : Analyse et Plan]
\textbf{Sujet :} \eng{“The government has decided to ban single-use plastic bags. Some people support this decision while others oppose it. Write a composition discussing both views and giving your own opinion.”} (250–300 words)

\textbf{Consignes :}
\begin{enumerate}
    \item Soulignez les mots-clés et identifiez le \textbf{type de sujet}.
    \item Formulez la \textbf{problématique} en anglais.
    \item Faites un \textbf{brainstorming} (5 idées pour, 5 idées contre) en anglais.
    \item Rédigez le \textbf{plan détaillé (Outline)} avec \emph{Topic Sentences} et \emph{Connecteurs} prévus pour 4 paragraphes de corps (2 pour, 2 contre / ou 3 arguments + opinion).
    \item Estimez le nombre de mots par paragraphe.
\end{enumerate}
\end{exercice}

\begin{corrige}[Exercice : Analyse et Plan]
\textbf{1. Mots-clés / Type :} \eng{Ban single-use plastic bags} (Thème), \eng{Discuss both views} (Tâche 1), \eng{Give your own opinion} (Tâche 2). Type : \textbf{Argumentation (Discussion + Opinion)}.
\textbf{2. Problématique :} \eng{Is the ban on single-use plastic bags an effective environmental measure or an economic burden for citizens and traders?}
\textbf{3. Brainstorming (Extrait) :}
\begin{itemize}
    \item \textbf{Pour (Pro-ban):} Less litter (visual pollution), protects wildlife (choking animals), reduces flooding (blocked drains), promotes alternatives (paper, baskets), long-term health (microplastics).
    \item \textbf{Contre (Anti-ban):} Cost of alternatives for poor families, job losses in plastic factories, enforcement corruption (bribes at market), lack of cheap alternatives available, hygiene for raw meat/food.
\end{itemize}
\textbf{4. Plan (Structure Discussion + Opinion) :}
\begin{itemize}
    \item \textbf{Intro:} Context (ban announced), Controversy, Thesis: \eng{While the ban protects the environment, it hurts the informal economy; a phased approach with subsidies is best.}
    \item \textbf{Body 1 (View A - Pro):} Environmental benefits. \eng{Firstly, plastic bags clog drains causing floods in Douala. Secondly, they kill livestock eating them.} Connectors: \eng{Proponents argue that, Moreover}.
    \item \textbf{Body 2 (View A - Pro):} Health/Long term. \eng{Microplastics in food chain. Rwanda/Kenya success stories.} Connector: \eng{In addition}.
    \item \textbf{Body 3 (View B - Con):} Economic hardship. \eng{Market women (bayam-sellam — \textit{petits commerçants ambulants}) cannot afford paper bags. Plastic factory workers laid off.} Connector: \eng{However, Opponents claim that}.
    \item \textbf{Body 4 (View B - Con / Nuance):} Enforcement issues. \eng{Corruption at checkpoints. No affordable woven baskets mass-produced.} Connector: \eng{Furthermore, Besides}.
    \item \textbf{Conclusion + Opinion:} \eng{In my view, the ban is necessary but must be gradual. Subsidise alternatives, support factories to convert, educate.} Connectors: \eng{To conclude, I firmly believe that}.
\end{itemize}
\textbf{5. Mots :} Intro (40) + B1 (70) + B2 (60) + B3 (70) + B4 (60) + Conc (40) = 340 $\rightarrow$ \textbf{Réduire B2/B4} pour viser 290.
\end{corrige}

\begin{aretenir}[Règle d'or de la planification]
\textbf{Zero Plan = Zero Structure = Low Score.} Prenez \textbf{10 minutes} sur les 2h (ou 1h30) pour : Lire $\rightarrow$ Analyser $\rightarrow$ Brainstormer $\rightarrow$ Sélectionner $\rightarrow$ Planifier. Ces 10 minutes vous feront gagner 30 minutes de rédaction hésitante et vous assureront les points d'\emph{Organization}. \textbf{Le plan se fait en anglais, sur le brouillon, avec les connecteurs notés.}
\end{aretenir}

\coursec{La structure de la composition}

Dans la section précédente, vous avez appris à \emph{analyser le sujet} (identifier le thème, le type de production attendue, le destinataire) et à \emph{planifier} vos idées sous forme de plan. Il est maintenant temps de transformer ce plan en une composition réelle, organisée et lisible. En anglais comme en français, toute composition argumentative ou explicative suit une architecture en trois temps : une \cle{introduction}, un \cle{body} (développement) et une \cle{conclusion}. Mais l'anglais scolaire exige une rigueur particulière sur un point : \emph{un paragraphe = une seule idée principale}. C'est cette architecture que nous allons construire ensemble, brique par brique.

\begin{popfigure}
\begin{tikzpicture}[node distance=0.35cm]
\node[draw=popblue, fill=popblueL, rounded corners, text width=11.5cm, align=center, font=\bfseries] (intro) {INTRODUCTION (3--4 sentences)\\ \mdseries Hook $\rightarrow$ Context $\rightarrow$ Question $\rightarrow$ Plan};
\node[draw=popgreen, fill=popgreenL, rounded corners, text width=9cm, align=center, below=of intro, font=\bfseries] (p1) {BODY -- Paragraph 1\\ \mdseries The problem: uncollected garbage, blocked gutters, diseases};
\node[draw=poporange, fill=poporangeL, rounded corners, text width=7.5cm, align=center, below=of p1, font=\bfseries] (p2) {Paragraph 2\\ \mdseries The causes: lack of dustbins, few trucks, bad habits};
\node[draw=poppurple, fill=poppurpleL, rounded corners, text width=7.5cm, align=center, below=of p2, font=\bfseries] (p3) {Paragraph 3\\ \mdseries The solutions: collection, sorting, recycling, fines};
\node[draw=popgold, fill=popgoldL, rounded corners, text width=6cm, align=center, below=of p3, font=\bfseries] (concl) {CONCLUSION (2--3 sentences)\\ \mdseries Summary + opinion or recommendation};
\draw[-{Stealth}, thick, popmuted] (intro) -- (p1);
\draw[-{Stealth}, thick, popmuted] (p1) -- (p2);
\draw[-{Stealth}, thick, popmuted] (p2) -- (p3);
\draw[-{Stealth}, thick, popmuted] (p3) -- (concl);
\end{tikzpicture}
\legende{La composition en entonnoir : on part du général (introduction), on détaille (développement), puis on resserre (conclusion).}
\end{popfigure}

\courssub{L'introduction : capter, situer, annoncer}

\textbf{Le rôle de l'introduction}\par

L'introduction compte \cle{3 à 4 phrases} et remplit quatre fonctions, dans cet ordre :

\begin{enumerate}
\item \textbf{The hook} (l'amorce) : une phrase qui attire l'attention du lecteur.
\item \textbf{The context} (la présentation du problème) : on situe le sujet dans la réalité.
\item \textbf{The question} (la problématique) : une question qui annonce l'enjeu.
\item \textbf{The plan} (l'annonce du plan) : une phrase qui indique ce que la composition va traiter.
\end{enumerate}

\begin{methode}[Rédiger une introduction en quatre coups : Hook $\rightarrow$ Context $\rightarrow$ Question $\rightarrow$ Plan]
\begin{enumerate}
\item \textbf{Hook} : ouvrez par une question, un fait marquant, une statistique ou un proverbe.
\item \textbf{Context} : décrivez la situation en une phrase (\eng{This situation threatens our health.}).
\item \textbf{Question} : formulez la problématique (\eng{How can we solve this problem?}).
\item \textbf{Plan} : annoncez le contenu (\eng{This essay will examine the causes and propose solutions.}).
\end{enumerate}
Modèle complet : \eng{Every day, tons of waste pile up in our streets. This situation threatens our health. How can we solve this problem? This essay will examine the causes and propose solutions.} « Chaque jour, des tonnes de déchets s'entassent dans nos rues. Cette situation menace notre santé. Comment résoudre ce problème ? Cette dissertation examinera les causes et proposera des solutions. »
\end{methode}

\textbf{Les quatre types d'amorces (hooks)}\par

Le choix de l'amorce donne le ton de toute la composition. Voici les quatre types les plus efficaces au lycée :

\begin{center}
\begin{tabularx}{\linewidth}{|l|X|X|}
\hline
\rowcolor{popblueL} \textbf{Type d'amorce} & \textbf{Exemple en anglais} & \textbf{Traduction} \\
\hline
Question & \eng{Have you ever walked past a mountain of garbage on your way to school?} & Avez-vous déjà passé devant une montagne d'ordures en allant à l'école ? \\
\hline
Fait marquant & \eng{In many quarters of Douala, garbage stays uncollected for weeks.} & Dans de nombreux quartiers de Douala, les ordures restent non ramassées pendant des semaines. \\
\hline
Statistique & \eng{A large city can produce thousands of tons of waste every single day.} & Une grande ville peut produire des milliers de tonnes de déchets chaque jour. \\
\hline
Proverbe & \eng{“Cleanliness is next to godliness,” says an old English proverb.} & « La propreté est proche de la sainteté », dit un vieux proverbe anglais. \\
\hline
\end{tabularx}
\end{center}

\begin{textebox}[Un proverbe comme amorce]
“Cleanliness is next to godliness.”

This old English proverb means that being clean is almost as important as being morally good. It is often used as a hook in compositions about hygiene, health and the environment. A student writing about garbage collection in Yaoundé could begin: “Cleanliness is next to godliness,” says an old proverb. Yet in many of our neighbourhoods, heaps of garbage lie in the streets for days. This situation is dangerous for our health. What can be done to improve garbage collection and recycling services? This essay will first describe the problem, then explain its causes, and finally suggest practical solutions.
\end{textebox}

Glossaire : \emph{godliness} (nom) : sainteté, piété ; \emph{heap} (nom) : tas, monceau ; \emph{neighbourhood} (nom) : quartier ; \emph{yet} (conjonction) : pourtant, cependant.

\begin{attention}[Statistiques : restez vague si vous n'êtes pas sûr]
N'inventez jamais un chiffre précis que vous ne pouvez pas vérifier (\eng{73\% of Cameroonians...}). Préférez une formulation générale : \eng{tons of waste}, \eng{thousands of people}, \eng{every year, many children fall ill}. Un correcteur pardonne une statistique vague, mais sanctionne un chiffre faux présenté comme un fait.
\end{attention}

\courssub{Le développement : un paragraphe, une idée}

\begin{propriete}[Règle d'or : one paragraph, one main idea]
En anglais scolaire, chaque paragraphe du développement ne défend qu'\textbf{une seule idée principale}. Ce paragraphe s'organise ainsi :
\begin{enumerate}
\item une \cle{topic sentence} (phrase d'ouverture) qui annonce clairement l'idée du paragraphe ;
\item des \emph{explications} qui développent cette idée ;
\item des \emph{exemples} concrets qui l'illustrent.
\end{enumerate}
Toute phrase qui ne sert pas l'idée principale doit être supprimée ou déplacée dans un autre paragraphe.
\end{propriete}

Comparaison avec le français : la dissertation française accepte parfois des paragraphes longs qui enchaînent plusieurs arguments. La composition anglaise, elle, exige une \emph{topic sentence} explicite dès la première ligne : le lecteur doit connaître l'idée du paragraphe avant même de le lire. C'est une différence de culture d'écriture : l'anglais va \emph{droit au but}.

\begin{exemplebox}[Topic sentences et phrases de conclusion : modèles traduits]
\eng{The first problem is the irregular collection of garbage.} « Le premier problème est la collecte irrégulière des ordures. »

\eng{There are several causes of this situation.} « Il y a plusieurs causes à cette situation. »

\eng{Fortunately, many solutions can be found.} « Heureusement, de nombreuses solutions peuvent être trouvées. »

\eng{In conclusion, recycling is the key to a cleaner city.} « En conclusion, le recyclage est la clé d'une ville plus propre. »
\end{exemplebox}

\textbf{Le plan type en trois paragraphes}\par

Pour un sujet sur \emph{garbage collection, disposal and recycling services}, le développement le plus naturel suit la logique \emph{problème $\rightarrow$ causes $\rightarrow$ solutions} :

\begin{center}
\begin{tabularx}{\linewidth}{|l|X|X|}
\hline
\rowcolor{popblueL} \textbf{Paragraphe} & \textbf{Idée principale (topic sentence)} & \textbf{Exemples à développer} \\
\hline
P1 : the problem & \eng{The first problem is the irregular collection of garbage.} & uncollected garbage, blocked gutters, bad smells, diseases like cholera and malaria \\
\hline
P2 : the causes & \eng{This situation has several causes.} & lack of dustbins, few garbage trucks, people's bad habits (throwing waste anywhere) \\
\hline
P3 : the solutions & \eng{Fortunately, practical solutions exist.} & regular collection, sorting, recycling, fines for polluters, awareness campaigns in schools \\
\hline
\end{tabularx}
\end{center}

\textbf{Anatomie d'un paragraphe complet}\par

Voici le paragraphe 1 rédigé selon la règle \emph{topic sentence + explications + exemples} :

\begin{exemplebox}[Paragraph 1 -- The problem]
\eng{The first problem is the irregular collection of garbage. In many neighbourhoods, waste is not collected for several days or even weeks. As a result, heaps of rubbish pile up in the streets and block the gutters. When the rains come, the dirty water cannot flow away, and it floods the houses. This stagnant water attracts mosquitoes, which spread malaria, and the rotting garbage can cause diseases such as cholera.}

« Le premier problème est la collecte irrégulière des ordures. Dans de nombreux quartiers, les déchets ne sont pas ramassés pendant plusieurs jours, voire plusieurs semaines. En conséquence, des tas d'ordures s'entassent dans les rues et bouchent les caniveaux. Quand les pluies arrivent, l'eau sale ne peut pas s'écouler et elle inonde les maisons. Cette eau stagnante attire les moustiques, qui propagent le paludisme, et les ordures en décomposition peuvent provoquer des maladies comme le choléra. »
\end{exemplebox}

Observez la mécanique : la première phrase annonce l'idée (\emph{irregular collection}), les trois phrases suivantes expliquent les conséquences (\emph{blocked gutters, floods}), et les deux dernières donnent des exemples précis (\emph{malaria, cholera}). Rien ne sort du sujet du paragraphe.

\textbf{La présentation visuelle des paragraphes}\par

En anglais, les paragraphes doivent être \emph{visuellement séparés}. Deux conventions sont acceptées :

\begin{itemize}
\item l'\cle{alinéa} (indentation) : la première ligne de chaque paragraphe commence en retrait ;
\item la \cle{ligne sautée} : on laisse une ligne vide entre deux paragraphes.
\end{itemize}

Choisissez \emph{une seule} convention et gardez-la du début à la fin. Ne mélangez jamais les deux.

\courssub{La conclusion : résumer et recommander}

La conclusion compte \cle{2 à 3 phrases} et remplit deux fonctions :

\begin{enumerate}
\item \textbf{résumer} les idées principales du développement (en d'autres mots, sans copier-coller) ;
\item donner une \textbf{opinion personnelle} ou une \textbf{recommandation} finale.
\end{enumerate}

\begin{exemplebox}[Modèle de conclusion]
\eng{In conclusion, garbage collection is a serious problem caused by the lack of dustbins and by people's bad habits. However, regular collection, sorting and recycling can make our cities cleaner. In my opinion, every citizen should take part in keeping our environment healthy.}

« En conclusion, la collecte des ordures est un problème sérieux causé par le manque de poubelles et par les mauvaises habitudes des gens. Cependant, une collecte régulière, le tri et le recyclage peuvent rendre nos villes plus propres. À mon avis, chaque citoyen devrait participer à la protection de notre environnement. »
\end{exemplebox}

\begin{attention}[Jamais d'idée nouvelle dans la conclusion]
La conclusion ne doit \textbf{rien introduire de nouveau} : ni exemple inédit, ni cause supplémentaire, ni solution non traitée dans le développement. Elle se contente de \emph{résumer} et de \emph{recommander}. Erreur fréquente : terminer par \eng{Moreover, the government should build more factories.} — si cette idée n'apparaît nulle part dans le corps, elle n'a pas sa place ici.

Autre piège de francophone : pour exprimer « selon moi », n'écrivez jamais \eng{according to me} (calque du français, incorrect en anglais) ; dites \eng{in my opinion} ou \eng{I believe that}. Notez aussi qu'après la virgule, \eng{in my opinion} s'écrit avec une minuscule : \eng{In conclusion, recycling is essential. In my opinion, everyone should take part.}
\end{attention}

\begin{exoresolu}[Remettre une composition dans l'ordre]
Énoncé : les cinq phrases suivantes forment l'introduction et la conclusion d'une composition sur le recyclage, mais elles sont mélangées. Remettez-les dans l'ordre logique et justifiez.

(a) \eng{In conclusion, recycling is the key to a cleaner city.}

(b) \eng{“Cleanliness is next to godliness,” says an old English proverb.}

(c) \eng{This essay will describe the problem and propose solutions.}

(d) \eng{Yet our streets are full of uncollected garbage.}

(e) \eng{I strongly recommend that every school should organise recycling campaigns.}

\tcblower
Solution détaillée : l'ordre correct est \textbf{(b) $\rightarrow$ (d) $\rightarrow$ (c) $\rightarrow$ (a) $\rightarrow$ (e)}.

\begin{itemize}
\item (b) est le \emph{hook} : le proverbe ouvre la composition.
\item (d) est le \emph{contexte} : le mot \eng{yet} (« pourtant ») crée un contraste avec le proverbe, donc (d) suit nécessairement (b).
\item (c) est l'\emph{annonce du plan} : \eng{this essay will...} ferme toujours l'introduction.
\item (a) ouvre la conclusion : le connecteur \eng{in conclusion} est un signal évident, et la phrase résume l'idée principale.
\item (e) est la \emph{recommandation} finale : \eng{I strongly recommend} donne l'opinion personnelle, sans idée nouvelle (les campagnes de recyclage ont été traitées dans le développement).
\end{itemize}
\end{exoresolu}

\begin{aretenir}[L'essentiel de la structure]
\begin{itemize}
\item \textbf{Introduction} (3--4 phrases) : hook $\rightarrow$ contexte $\rightarrow$ problématique $\rightarrow$ annonce du plan.
\item \textbf{Développement} (2--3 paragraphes) : un paragraphe = une idée ; topic sentence + explications + exemples. Plan type : problème $\rightarrow$ causes $\rightarrow$ solutions.
\item \textbf{Conclusion} (2--3 phrases) : résumé + opinion ou recommandation ; jamais d'idée nouvelle.
\item Séparez visuellement les paragraphes : alinéa \emph{ou} ligne sautée, mais pas les deux.
\end{itemize}
\end{aretenir}

Maintenant que l'architecture est solide, il reste à la faire vivre : la section suivante vous donnera les \emph{connecteurs logiques} et les \emph{temps verbaux} qui relient ces paragraphes entre eux et donnent à votre composition sa fluidité.

\coursec{Connecteurs logiques et temps verbaux}

Dans la section précédente, nous avons appris à construire le squelette d'une composition : une introduction qui présente le sujet, un développement en deux ou trois paragraphes, et une conclusion qui propose une solution ou une opinion. Mais un squelette sans articulations ne peut pas bouger ! Ce qui relie vos idées entre elles et rend votre texte fluide, ce sont les \cle{connecteurs logiques} (\eng{link words} ou \eng{connectors}). Et ce qui donne de la précision à vos phrases, ce sont les \cle{temps verbaux} bien choisis. Cette section vous donne les deux outils indispensables pour transformer une liste d'idées en une vraie composition digne d'un élève de Terminale.

\courssub{Pourquoi les connecteurs sont-ils essentiels ?}

Observez ces deux versions du même paragraphe. Laquelle est la meilleure ?

\begin{exemplebox}[Deux versions à comparer]
\textbf{Version A (sans connecteurs) :} \eng{Garbage is not collected regularly in our neighbourhood. Diseases spread quickly. Recycling creates jobs. The council does not provide enough dustbins. The problem remains serious.}

\textbf{Version B (avec connecteurs) :} \eng{Because garbage is not collected regularly in our neighbourhood, diseases spread quickly. Moreover, recycling creates jobs. However, the council does not provide enough dustbins. Therefore, the problem remains serious.}
\end{exemplebox}

La version B est clairement supérieure : le lecteur comprend immédiatement les relations logiques entre les idées (cause, addition, opposition, conséquence). À l'examen, l'utilisation variée et correcte des connecteurs est un critère d'évaluation explicite : un texte sans connecteurs perd des points en \emph{coherence and cohesion}.

\begin{exemplebox}[Exemples traduits]
\eng{Because garbage is not collected regularly, diseases spread quickly.} = « Parce que les ordures ne sont pas collectées régulièrement, les maladies se propagent vite. »

\eng{Moreover, recycling creates jobs.} = « De plus, le recyclage crée des emplois. »

\eng{Waste is piling up in our streets; therefore, the council must act immediately.} = « Les déchets s'accumulent dans nos rues ; par conséquent, la mairie doit agir immédiatement. »
\end{exemplebox}

\courssub{Les grandes familles de connecteurs}

\textbf{1. Les connecteurs d'addition}\par

Ils servent à ajouter une idée, un argument ou un exemple supplémentaire.

\begin{center}
\begin{tabularx}{\linewidth}{|l|X|X|}
\hline
\rowcolor{popblueL} Connecteur & Exemple anglais & Traduction \\
\hline
\eng{first of all} & \eng{First of all, we must reduce the waste we produce.} & Tout d'abord, nous devons réduire les déchets que nous produisons. \\
\hline
\eng{moreover} & \eng{Moreover, garbage attracts rats and flies.} & De plus, les ordures attirent les rats et les mouches. \\
\hline
\eng{furthermore} & \eng{Furthermore, landfills pollute the soil and the water.} & En outre, les décharges polluent le sol et l'eau. \\
\hline
\eng{in addition} & \eng{In addition, the council should build a recycling centre.} & De plus, la mairie devrait construire un centre de recyclage. \\
\hline
\eng{also} & \eng{Recycling also saves natural resources.} & Le recyclage économise aussi les ressources naturelles. \\
\hline
\eng{besides} & \eng{Besides, plastic bags block the gutters.} & En outre, les sachets plastiques bouchent les caniveaux. \\
\hline
\end{tabularx}
\end{center}

\begin{attention}[Position de \eng{also}]
Contrairement à \eng{moreover} et \eng{furthermore} qui se placent en tête de phrase suivis d'une virgule, \eng{also} se place \textbf{au milieu de la phrase}, avant le verbe principal (ou après \eng{be}) : \eng{Recycling \textbf{also} creates jobs.} Dans une rédaction soignée, évitez de commencer vos phrases par \eng{also}. De même, \eng{besides} en tête de phrase signifie « en outre », mais attention au faux ami : \eng{beside} (sans \emph{s}) signifie « à côté de » !
\end{attention}

\textbf{2. Les connecteurs de cause et de conséquence}\par

C'est la famille la plus utile pour un sujet comme la gestion des déchets, car vous devez constamment expliquer \emph{pourquoi} le problème existe et \emph{quelles} en sont les conséquences.

\begin{center}
\begin{tabularx}{\linewidth}{|l|X|X|}
\hline
\rowcolor{popgreenL} Connecteur & Exemple anglais & Traduction \\
\hline
\eng{because} & \eng{People dump waste in the river because there are no dustbins.} & Les gens jettent les déchets dans la rivière parce qu'il n'y a pas de poubelles. \\
\hline
\eng{since / as} & \eng{Since the population is growing, waste is increasing too.} & Puisque la population augmente, les déchets augmentent aussi. \\
\hline
\eng{therefore} & \eng{The situation is dangerous; therefore, we must react.} & La situation est dangereuse ; par conséquent, nous devons réagir. \\
\hline
\eng{consequently} & \eng{Rubbish blocks the drains; consequently, floods occur every rainy season.} & Les ordures bouchent les caniveaux ; par conséquent, des inondations surviennent chaque saison des pluies. \\
\hline
\eng{as a result} & \eng{The trucks broke down. As a result, garbage was not collected for two weeks.} & Les camions sont tombés en panne. En conséquence, les ordures n'ont pas été collectées pendant deux semaines. \\
\hline
\eng{that is why} & \eng{Plastic does not decompose. That is why we must recycle it.} & Le plastique ne se décompose pas. C'est pourquoi nous devons le recycler. \\
\hline
\eng{so} & \eng{The bins were full, so people threw rubbish on the ground.} & Les poubelles étaient pleines, alors les gens ont jeté les ordures par terre. \\
\hline
\end{tabularx}
\end{center}

\begin{propriete}[Cause et conséquence : deux constructions à ne pas mélanger]
\begin{itemize}
\item \eng{because / since / as} + \textbf{proposition} (sujet + verbe) introduisent la \textbf{cause} : \eng{Because it rains, the gutters overflow.}
\item \eng{therefore / consequently / as a result / that is why} introduisent la \textbf{conséquence}, généralement en début de phrase suivis d'une virgule : \eng{It rains heavily; therefore, the gutters overflow.}
\item \textbf{Interdit} : on ne combine jamais \eng{because} et \eng{so} dans la même phrase comme en français familier. On dit \eng{Because it rains, the gutters overflow} OU \eng{It rains, so the gutters overflow}, mais jamais \eng{Because it rains, so the gutters overflow}.
\end{itemize}
\end{propriete}

\begin{attention}[\eng{because} ou \eng{because of} ?]
Erreur très fréquente chez les francophones ! \eng{because} est suivi d'une \textbf{proposition} (sujet + verbe), tandis que \eng{because of} est suivi d'un \textbf{nom ou groupe nominal} :
\begin{itemize}
\item \eng{The collection stopped \textbf{because} the trucks broke down.} (sujet + verbe)
\item \eng{The collection stopped \textbf{because of} the breakdown of the trucks.} (nom)
\item \eng{\textbf{Because of} the rain, the streets are flooded.} mais \eng{\textbf{Because} it rains, the streets are flooded.}
\end{itemize}
Ne dites jamais \eng{because of it rains} !
\end{attention}

\textbf{3. Les connecteurs d'opposition}\par

Indispensables pour nuancer, présenter les deux faces d'un problème ou exprimer un regret.

\begin{center}
\begin{tabularx}{\linewidth}{|l|X|X|}
\hline
\rowcolor{poporangeL} Connecteur & Exemple anglais & Traduction \\
\hline
\eng{however} & \eng{The council promises action. However, nothing has changed.} & La mairie promet d'agir. Cependant, rien n'a changé. \\
\hline
\eng{nevertheless} & \eng{The task is huge; nevertheless, volunteers keep working.} & La tâche est énorme ; néanmoins, les volontaires continuent à travailler. \\
\hline
\eng{whereas} & \eng{In Buea, waste is sorted, whereas in many towns it is not.} & À Buea, les déchets sont triés, alors que dans beaucoup de villes ils ne le sont pas. \\
\hline
\eng{although} & \eng{Although recycling is expensive, it creates jobs.} & Bien que le recyclage coûte cher, il crée des emplois. \\
\hline
\eng{on the other hand} & \eng{Burning waste is cheap. On the other hand, it pollutes the air.} & Brûler les déchets coûte peu. En revanche, cela pollue l'air. \\
\hline
\eng{unfortunately} & \eng{Unfortunately, many people still throw plastic into the rivers.} & Malheureusement, beaucoup de gens jettent encore du plastique dans les rivières. \\
\hline
\end{tabularx}
\end{center}

\begin{attention}[La position de \eng{however} et le piège \eng{although\ldots but}]
\eng{however} ne se place pas comme « cependant » n'importe où dans la phrase. En anglais soigné, il s'emploie \textbf{en début de phrase, suivi d'une virgule} : \eng{However, the problem remains serious.} (« Cependant, le problème reste sérieux. »). Évitez de le coincer entre le sujet et le verbe comme on le fait parfois en français. De même, \eng{although} introduit une subordonnée et ne peut pas être suivi de \eng{but} : on dit \eng{Although recycling is expensive, it creates jobs} et jamais \eng{Although recycling is expensive, but it creates jobs}.
\end{attention}

\textbf{4. Les connecteurs d'exemple et d'ordre}\par

Pour illustrer vos arguments et organiser clairement votre développement.

\begin{center}
\begin{tabularx}{\linewidth}{|l|X|X|}
\hline
\rowcolor{poppurpleL} Connecteur & Exemple anglais & Traduction \\
\hline
\eng{for example / for instance} & \eng{Many materials can be recycled, for example plastic and paper.} & Beaucoup de matériaux peuvent être recyclés, par exemple le plastique et le papier. \\
\hline
\eng{such as} & \eng{Organic waste, such as banana peels, can become compost.} & Les déchets organiques, tels que les épluchures de bananes, peuvent devenir du compost. \\
\hline
\eng{firstly / secondly} & \eng{Firstly, we should sort our waste. Secondly, we should compost organic matter.} & Premièrement, nous devrions trier nos déchets. Deuxièmement, nous devrions composter la matière organique. \\
\hline
\eng{finally} & \eng{Finally, the council should punish illegal dumping.} & Enfin, la mairie devrait punir les dépôts sauvages. \\
\hline
\eng{to conclude} & \eng{To conclude, keeping our town clean is everybody's duty.} & Pour conclure, garder notre ville propre est le devoir de chacun. \\
\hline
\end{tabularx}
\end{center}

\begin{propriete}[\eng{such as} contre \eng{for example}]
\eng{such as} s'insère directement après le nom qu'il illustre, sans virgule obligatoire : \eng{Materials such as glass and metal can be recycled.} \eng{for example} et \eng{for instance} sont des expressions indépendantes, encadrées de virgules ou placées en tête de phrase : \eng{Many solutions exist. For instance, households can sort their waste.}
\end{propriete}

\textbf{5. Les connecteurs de but : un bonus pour briller}\par

Pour exprimer le \textbf{but} d'une action (« pour », « afin de »), utilisez \eng{to} + base verbale, ou les formes plus soutenues \eng{in order to} et \eng{so as to} :

\begin{exemplebox}[Exprimer le but]
\eng{The council has bought new trucks \textbf{in order to} improve the collection service.} = « La mairie a acheté de nouveaux camions afin d'améliorer le service de collecte. »

\eng{We sort our waste \textbf{so as to} make recycling easier.} = « Nous trions nos déchets afin de faciliter le recyclage. »

\eng{People should reuse plastic bags \textbf{to} reduce pollution.} = « Les gens devraient réutiliser les sachets plastiques pour réduire la pollution. »
\end{exemplebox}

\begin{attention}[\eng{for} n'exprime pas le but devant un verbe]
Ne traduisez jamais « pour + infinitif » par \eng{for} + verbe : on dit \eng{We recycle \textbf{to protect} the environment} et non \eng{for protect} ni \eng{for protecting}. \eng{for} + nom exprime la destination : \eng{a centre \textbf{for recycling}} (« un centre de recyclage »).
\end{attention}

\courssub{Carte mentale : choisir le bon connecteur}

\begin{popfigure}
\begin{center}
\begin{tikzpicture}[mindmap, grow cyclic, every node/.style={concept, align=center, font=\small},
root concept/.append style={concept color=popblue, font=\bfseries},
level 1/.append style={level distance=3.4cm, sibling angle=72},
level 2/.append style={level distance=2.4cm, sibling angle=45, font=\scriptsize}]
\node[root concept]{Link words}
  child[concept color=popgreen] { node {Addition}
      child { node {moreover} } child { node {in addition} } }
  child[concept color=poporange] { node {Cause}
      child { node {because} } child { node {since} } }
  child[concept color=poppink] { node {Conséquence}
      child { node {therefore} } child { node {as a result} } }
  child[concept color=popgold] { node {Opposition}
      child { node {however} } child { node {although} } }
  child[concept color=poppurple] { node {Exemple / ordre}
      child { node {for example} } child { node {to conclude} } };
\end{tikzpicture}
\end{center}
\legende{Carte mentale des connecteurs logiques classés par fonction}
\end{popfigure}

\begin{methode}[Varier ses connecteurs : la règle d'or]
\begin{enumerate}
\item Avant de rédiger, préparez au brouillon une liste de \textbf{8 à 10 connecteurs variés} (au moins un par fonction : addition, cause, conséquence, opposition, exemple, ordre, conclusion).
\item \textbf{N'utilisez jamais deux fois le même connecteur} dans une composition : si vous avez écrit \eng{moreover} au paragraphe 1, utilisez \eng{furthermore} ou \eng{in addition} au paragraphe 2.
\item Cochez chaque connecteur de votre liste au fur et à mesure que vous l'employez.
\item À la relecture, vérifiez que chaque paragraphe contient au moins deux connecteurs et que la ponctuation est correcte (virgule après \eng{however}, \eng{therefore}, \eng{for example} en tête de phrase).
\end{enumerate}
\end{methode}

\courssub{Les temps verbaux de la composition sur l'environnement}

Un sujet comme \eng{garbage collection, disposal and recycling services} appelle quatre outils verbaux précis. Observons d'abord un court texte.

\begin{textebox}[Our town is getting dirty — extrait]
\eng{Every morning, garbage collectors work in the streets of our town, but their efforts are not enough. Nowadays, waste is piling up in many neighbourhoods because the population is growing fast. Rubbish is dumped in illegal landfills, and plastic bags are thrown into the rivers. As a result, diseases such as cholera spread during the rainy season. Fortunately, solutions exist. Organic waste can be composted, and plastic can be recycled into new products. The council should provide more dustbins, and we must all sort our waste. If everyone makes an effort, our town will become clean and healthy again.}

\textbf{Glossaire :} \emph{to pile up} (v.) : s'accumuler ; \emph{landfill} (n.) : décharge ; \emph{to spread} (v.) : se propager ; \emph{dustbin} (n.) : poubelle ; \emph{to sort} (v.) : trier.

\textbf{Questions :} 1. Why is waste piling up? 2. What happens during the rainy season? 3. What should the council do?
\end{textebox}

\textbf{1. Le présent simple : les faits généraux et les habitudes}\par

\begin{propriete}[Présent simple : forme et emploi]
\textbf{Forme} : base verbale (+ \emph{s} à la 3\textsuperscript{e} personne du singulier). Négation et question avec \eng{do/does}.

\textbf{Emploi dans ce type de composition} : vérités générales, faits permanents, habitudes et routines des services de collecte.
\begin{itemize}
\item \eng{Garbage collectors work every morning.} = « Les éboueurs travaillent chaque matin. »
\item \eng{Plastic does not decompose naturally.} = « Le plastique ne se décompose pas naturellement. »
\item \eng{The truck comes twice a week.} = « Le camion passe deux fois par semaine. »
\end{itemize}
\textbf{Marqueurs typiques} : \eng{every day, always, usually, twice a week, in general}.
\end{propriete}

\textbf{2. Le présent continu : la situation actuelle, en évolution}\par

\begin{propriete}[Présent continu : forme et emploi]
\textbf{Forme} : \eng{be} (au présent) + verbe en \emph{-ing}.

\textbf{Emploi} : une situation en cours \emph{maintenant}, souvent inquiétante ou en train de changer — parfait pour décrire l'état actuel de votre quartier.
\begin{itemize}
\item \eng{Waste is piling up in our streets.} = « Les déchets s'accumulent dans nos rues. »
\item \eng{The population is growing, and more rubbish is being produced.} = « La population augmente, et davantage d'ordures sont produites. »
\end{itemize}
\textbf{Marqueurs typiques} : \eng{nowadays, at present, these days, currently}.
\end{propriete}

\begin{attention}[Présent simple ou présent continu ?]
Comparez : \eng{Garbage collectors \textbf{work} every morning} (habitude régulière) et \eng{Garbage \textbf{is piling up} these days} (situation actuelle temporaire et alarmante). Ne mettez jamais un verbe d'état au continu : on dit \eng{Plastic \textbf{pollutes} the environment} pour une vérité générale ; le continu ne s'emploie que si l'on insiste sur un phénomène en cours d'évolution : \eng{Plastic is increasingly polluting our oceans.}
\end{attention}

\textbf{3. La voix passive : l'outil star du thème}\par

\begin{propriete}[La voix passive : \eng{be} + participe passé]
\textbf{Forme} : auxiliaire \eng{be} (conjugué au temps voulu) + \textbf{participe passé} du verbe : \eng{is collected, was dumped, can be recycled, should be provided}.

\textbf{Emploi} : la voix passive est idéale quand \textbf{l'agent est évident, inconnu ou peu important} ; l'essentiel est l'action ou son résultat. C'est exactement le cas pour les services de déchets !
\begin{itemize}
\item \eng{Waste is collected twice a week.} = « Les déchets sont collectés deux fois par semaine. »
\item \eng{Rubbish is dumped in landfills.} = « Les ordures sont déversées dans des décharges. »
\item \eng{Plastic can be recycled.} = « Le plastique peut être recyclé. »
\item \eng{More dustbins should be provided by the council.} = « Plus de poubelles devraient être fournies par la mairie. »
\end{itemize}
\textbf{Transformation} : \eng{The council collects waste} (actif) $\rightarrow$ \eng{Waste is collected (by the council)} (passif). Le complément d'objet direct devient sujet ; l'ancien sujet devient complément d'agent introduit par \eng{by}, souvent omis.
\end{propriete}

\begin{center}
\begin{tabular}{|l|l|l|}
\hline
\rowcolor{popblueL} Temps / modal & Actif & Passif \\
\hline
présent simple & \eng{collects} & \eng{is collected} \\
\hline
présent continu & \eng{is collecting} & \eng{is being collected} \\
\hline
prétérit & \eng{collected} & \eng{was collected} \\
\hline
futur & \eng{will collect} & \eng{will be collected} \\
\hline
\eng{can / should / must} & \eng{can recycle} & \eng{can be recycled} \\
\hline
\end{tabular}
\end{center}

\begin{attention}[Pièges de la voix passive]
\begin{itemize}
\item N'oubliez jamais l'auxiliaire \eng{be} : \eng{Waste \textbf{is} collected} et non \eng{Waste collected}.
\item Le participe passé des verbes irréguliers doit être exact : \eng{is \textbf{thrown}} (throw–threw–thrown), \eng{is \textbf{burnt}} (burn–burnt–burnt), \eng{is \textbf{sorted}} (sort, régulier).
\item Ne confondez pas avec le français « se + verbe » : « le plastique se recycle » se traduit par la voix passive \eng{plastic can be recycled} ou par l'adjectif \eng{plastic is recyclable}, jamais par un pronom réfléchi.
\end{itemize}
\end{attention}

\textbf{4. Futur et modaux : formuler des recommandations}\par

La conclusion d'une composition sur l'environnement propose presque toujours des solutions. Les modaux sont vos meilleurs alliés.

\begin{center}
\begin{tabularx}{\linewidth}{|l|X|X|}
\hline
\rowcolor{popgreenL} Outil & Exemple anglais & Traduction \\
\hline
\eng{should} (conseil) & \eng{The council should provide more dustbins.} & La mairie devrait fournir plus de poubelles. \\
\hline
\eng{must} (obligation forte) & \eng{We must sort our waste.} & Nous devons trier nos déchets. \\
\hline
\eng{will} (prédiction, promesse) & \eng{If we recycle, our town will become cleaner.} & Si nous recyclons, notre ville deviendra plus propre. \\
\hline
\eng{can} (possibilité) & \eng{Organic waste can be turned into compost.} & Les déchets organiques peuvent être transformés en compost. \\
\hline
\eng{ought to} (devoir moral) & \eng{Every citizen ought to keep the gutters clean.} & Tout citoyen devrait garder les caniveaux propres. \\
\hline
\end{tabularx}
\end{center}

\begin{attention}[Après un modal, la base verbale !]
Après \eng{should, must, can, will}, le verbe est toujours à la \textbf{base verbale}, sans \eng{to} et sans \emph{s} : \eng{We must \textbf{sort} our waste} et non \eng{We must to sort} ni \eng{We must sorts}. Et au passif : \eng{Waste must \textbf{be sorted}} (modal + \eng{be} + participe passé). Seul \eng{ought to} fait exception : il est suivi de \eng{to} + base verbale (\eng{ought to keep}).
\end{attention}

\courssub{Exercice résolu}

\begin{exoresolu}[Connecteurs et temps : à vous de jouer]
\textbf{Énoncé.} Complétez ce paragraphe avec les connecteurs et les formes verbales appropriés : \eng{however, because, moreover, for example, that is why} ; mettez les verbes entre parenthèses au temps et à la voix corrects.

\eng{Garbage is a serious problem in our town. (1) \ldots\ the population (grow) \ldots\ fast, more and more waste (produce) \ldots\ every day. (2) \ldots, the collection service is not regular. (3) \ldots, waste (pile up) \ldots\ in the streets these days. (4) \ldots, in the market area, rubbish (not collect) \ldots\ for two weeks last month. (5) \ldots, the council should buy new trucks and employ more workers.}

\tcblower

\textbf{Solution détaillée.}
\begin{enumerate}
\item \eng{\textbf{Because} the population \textbf{is growing} fast, more and more waste \textbf{is produced} every day.} — Cause : \eng{because} + proposition ; la croissance est en cours (présent continu) ; la production de déchets est un fait régulier subi (présent simple passif : \eng{is produced}).
\item \eng{\textbf{Moreover}, the collection service is not regular.} — Addition d'un argument ; virgule après le connecteur.
\item \eng{\textbf{As a result}, waste \textbf{is piling up} in the streets these days.} — Conséquence logique ; \eng{these days} impose le présent continu. (On pouvait aussi écrire \eng{Therefore} ou \eng{Consequently}.)
\item \eng{\textbf{For example}, in the market area, rubbish \textbf{was not collected} for two weeks last month.} — Illustration ; \eng{last month} impose le prétérit, à la voix passive car l'agent est évident.
\item \eng{\textbf{That is why}, the council should buy new trucks and employ more workers.} — Conclusion-recommandation avec le modal \eng{should} + base verbale. On aurait aussi pu utiliser \eng{however} plus haut pour opposer les promesses de la mairie à la réalité : \eng{The council promises action; \textbf{however}, the service remains irregular.}
\end{enumerate}
\end{exoresolu}

\begin{exercice}[À votre tour]
Réécrivez les phrases suivantes en utilisant le connecteur ou la voix indiqués entre parenthèses :
\begin{enumerate}
\item \eng{People throw plastic into the river. Fish die.} (\eng{because} / \eng{as a result})
\item \eng{The council collects our waste twice a week.} (voix passive)
\item \eng{Recycling is expensive. It creates many jobs.} (\eng{although})
\item \eng{The trucks are old. The service is irregular.} (\eng{since} / \eng{consequently})
\item \eng{The government must build more recycling centres.} (voix passive avec \eng{must})
\end{enumerate}
\end{exercice}

\begin{corrige}[Exercice « À votre tour »]
\begin{enumerate}
\item \eng{Because people throw plastic into the river, fish die.} / \eng{People throw plastic into the river; as a result, fish die.}
\item \eng{Our waste is collected twice a week (by the council).}
\item \eng{Although recycling is expensive, it creates many jobs.}
\item \eng{Since the trucks are old, the service is irregular.} / \eng{The trucks are old; consequently, the service is irregular.}
\item \eng{More recycling centres must be built (by the government).}
\end{enumerate}
\end{corrige}

\begin{aretenir}[L'essentiel de la section]
\begin{itemize}
\item Les connecteurs donnent la \textbf{logique} de votre texte : addition (\eng{moreover, furthermore, in addition}), cause (\eng{because, since, as}), conséquence (\eng{therefore, consequently, as a result, that is why}), opposition (\eng{however, nevertheless, although, whereas}), exemple et ordre (\eng{for example, such as, firstly, to conclude}), but (\eng{to, in order to, so as to}).
\item \eng{because} + proposition ; \eng{because of} + nom. Jamais \eng{because\ldots so} ni \eng{although\ldots but} dans la même phrase.
\item \eng{however} : début de phrase + virgule.
\item Présent simple = faits et habitudes ; présent continu = situation actuelle ; voix passive (\eng{be} + participe passé) = l'action compte plus que l'agent ; modaux + base verbale = recommandations.
\item Variez vos connecteurs : jamais deux fois le même dans une composition.
\end{itemize}
\end{aretenir}

\coursec{Modèle commenté}

Dans la section précédente, nous avons rassemblé nos outils : les \cle{connecteurs logiques} (\eng{first of all}, \eng{as a result}, \eng{moreover}, \eng{in conclusion}) et les \cle{temps verbaux} adaptés à ce type de sujet. Il est temps maintenant de voir ces outils \emph{en action} dans une composition complète et réussie. Lire et décortiquer un bon modèle est la méthode la plus efficace pour progresser en \eng{writing} : on n'apprend pas à écrire seulement en écrivant, mais aussi en observant comment un texte bien construit fonctionne de l'intérieur.

\courssub{Le sujet et le modèle}

Voici le sujet type que l'on pourrait vous donner à l'examen :

\begin{exemplebox}[Sujet de composition]
\eng{Write a composition (about 250 words) on how garbage collection and recycling can be improved in your town.}

« Rédigez une composition (environ 250 mots) sur la manière dont la collecte des ordures et le recyclage peuvent être améliorés dans votre ville. »
\end{exemplebox}

Lisez maintenant le modèle ci-dessous une première fois, sans vous arrêter sur les mots inconnus : cherchez seulement le sens général et les trois grandes parties.

\begin{textebox}[Model composition — “A clean city is a healthy city”]
In many Cameroonian towns, heaps of uncollected garbage can be seen along the streets. This situation is dangerous for public health. How can waste management be improved? This essay will first describe the problem, then suggest practical solutions.

First of all, garbage is not collected regularly, so it piles up and blocks the gutters. As a result, during the rainy season, floods spread diseases such as cholera and malaria. Moreover, many people throw rubbish anywhere because there are not enough dustbins.

However, solutions exist. The city council should provide more garbage trucks and organise door-to-door collection twice a week. In addition, citizens must be educated to sort their waste: plastic, glass and metal can be recycled, while organic waste can be turned into compost. For example, in some Nigerian cities, recycling cooperatives have created jobs for young people.

In conclusion, better garbage collection and recycling are possible if the authorities and the population work together. A clean city is a healthy city.
\end{textebox}

\textbf{Glossaire utile}\par
\begin{itemize}
\item \eng{heaps of uncollected garbage} : des tas d'ordures non ramassées (\eng{heap} : tas, nom)
\item \eng{waste management} : la gestion des déchets
\item \eng{to pile up} : s'entasser, s'accumuler (phrasal verb)
\item \eng{gutters} : les caniveaux, les gouttières
\item \eng{to spread diseases} : propager des maladies
\item \eng{dustbins} : les poubelles (anglais britannique ; \eng{trash cans} en anglais américain)
\item \eng{door-to-door collection} : la collecte porte-à-porte
\item \eng{to sort waste} : trier les déchets
\item \eng{compost} : le compost (déchets organiques décomposés utilisés comme engrais)
\end{itemize}

\begin{aretenir}[Les trois parties repérées]
\begin{enumerate}
\item \textbf{Introduction} (paragraphe 1) : présentation du problème + annonce du plan.
\item \textbf{Développement} (paragraphes 2 et 3) : le problème en détail, puis les solutions.
\item \textbf{Conclusion} (paragraphe 4) : synthèse + formule finale percutante.
\end{enumerate}
Remarquez que le plan annoncé dans l'introduction (\eng{first describe the problem, then suggest practical solutions}) est \emph{exactement} suivi : un paragraphe pour le problème, un paragraphe pour les solutions. C'est la marque d'une composition maîtrisée.
\end{aretenir}

\courssub{Commentaire de l'introduction}

Une bonne introduction de composition argumentative contient trois éléments, dans cet ordre :

\begin{enumerate}
\item \textbf{Une amorce} (\eng{hook}) : un fait, une observation ou une question qui capte l'attention du lecteur. Ici : \eng{In many Cameroonian towns, heaps of uncollected garbage can be seen along the streets.} L'auteur part d'un fait concret et visuel que tout lecteur camerounais reconnaît immédiatement.
\item \textbf{La problématique} : la question à laquelle la composition va répondre. Ici, elle est formulée explicitement : \eng{How can waste management be improved?} (« Comment la gestion des déchets peut-elle être améliorée ? »). C'est une reformulation du sujet sous forme de question — excellente technique.
\item \textbf{L'annonce du plan} : \eng{This essay will first describe the problem, then suggest practical solutions.} Le lecteur sait désormais exactement ce qui l'attend.
\end{enumerate}

\begin{exemplebox}[Phrases d'introduction à réutiliser]
\eng{In many African cities, waste has become a major problem.} « Dans de nombreuses villes africaines, les déchets sont devenus un problème majeur. » (amorce par un fait)

\eng{What can be done to keep our town clean?} « Que peut-on faire pour garder notre ville propre ? » (problématique)

\eng{This essay will first analyse the causes of the problem and then propose solutions.} « Cette dissertation analysera d'abord les causes du problème, puis proposera des solutions. » (annonce du plan)
\end{exemplebox}

\courssub{Commentaire du développement}

Chaque paragraphe du développement est construit autour d'une \cle{topic sentence} (phrase d'ouverture qui annonce l'idée principale du paragraphe), suivie d'explications et d'exemples précis.

\textbf{Paragraphe 2 — le problème.} Topic sentence : \eng{First of all, garbage is not collected regularly, so it piles up and blocks the gutters.} L'idée est ensuite développée avec une conséquence (\eng{as a result} + les inondations et les maladies : \eng{cholera and malaria}) et une cause supplémentaire (\eng{moreover} + le manque de poubelles : \eng{there are not enough dustbins}).

\textbf{Paragraphe 3 — les solutions.} Topic sentence : \eng{However, solutions exist.} — courte, claire, efficace. Suivent deux solutions concrètes : la responsabilité des autorités (\eng{the city council should provide more garbage trucks}) et celle des citoyens (\eng{citizens must be educated to sort their waste}). Le paragraphe se termine par un exemple précis et crédible : \eng{in some Nigerian cities, recycling cooperatives have created jobs for young people.}

\begin{exemplebox}[Exemples clés traduits]
\eng{Garbage is not collected regularly, so it piles up.} « Les ordures ne sont pas collectées régulièrement, donc elles s'entassent. »

\eng{Citizens must be educated to sort their waste.} « Les citoyens doivent être sensibilisés au tri de leurs déchets. »

\eng{Organic waste can be turned into compost.} « Les déchets organiques peuvent être transformés en compost. »
\end{exemplebox}

\begin{popfigure}
\begin{tikzpicture}[node distance=0.45cm, every node/.style={font=\small}]
\node[draw=popblue, fill=popblueL, rounded corners, text width=5.6cm, align=center] (intro) {\textbf{§1 Introduction}\\hook : heaps of garbage\\question : how to improve?\\plan : problem + solutions};
\node[draw=poporange, fill=poporangeL, rounded corners, text width=5.6cm, align=center, below=of intro] (dev1) {\textbf{§2 Problem}\\topic sentence : not collected\\examples : floods, cholera,\\blocked gutters, no dustbins};
\node[draw=popgreen, fill=popgreenL, rounded corners, text width=5.6cm, align=center, below=of dev1] (dev2) {\textbf{§3 Solutions}\\topic sentence : solutions exist\\examples : trucks, sorting,\\recycling cooperatives};
\node[draw=poppurple, fill=poppurpleL, rounded corners, text width=5.6cm, align=center, below=of dev2] (concl) {\textbf{§4 Conclusion}\\summary : work together\\final message : clean = healthy};
\node[draw=popblue, fill=white, rounded corners, text width=4.1cm, align=center, right=1.1cm of intro] (f1) {\textbf{Function}\\catch the reader's\ and announce\\the plan};
\node[draw=poporange, fill=white, rounded corners, text width=4.1cm, align=center, right=1.1cm of dev1] (f2) {\textbf{Function}\\describe the problem\\with precise examples};
\node[draw=popgreen, fill=white, rounded corners, text width=4.1cm, align=center, right=1.1cm of dev2] (f3) {\textbf{Function}\\propose realistic\\solutions with an example};
\node[draw=poppurple, fill=white, rounded corners, text width=4.1cm, align=center, right=1.1cm of concl] (f4) {\textbf{Function}\\summarise and give\\a personal recommendation};
\draw[-{Stealth}, thick, popblue] (intro) -- (f1);
\draw[-{Stealth}, thick, poporange] (dev1) -- (f2);
\draw[-{Stealth}, thick, popgreen] (dev2) -- (f3);
\draw[-{Stealth}, thick, poppurple] (concl) -- (f4);
\draw[-{Stealth}, thick, popmuted] (intro) -- (dev1);
\draw[-{Stealth}, thick, popmuted] (dev1) -- (dev2);
\draw[-{Stealth}, thick, popmuted] (dev2) -- (concl);
\end{tikzpicture}
\legende{Schéma annoté du modèle : chaque paragraphe est relié à sa fonction dans la composition.}
\end{popfigure}

\courssub{Commentaire de la conclusion}

La conclusion du modèle fait deux choses, en deux phrases seulement :

\begin{enumerate}
\item \textbf{Une synthèse} qui répond à la problématique : \eng{better garbage collection and recycling are possible if the authorities and the population work together.} (« une meilleure collecte et un meilleur recyclage sont possibles si les autorités et la population travaillent ensemble »). Notez le conditionnel avec \eng{if} : la solution est présentée comme réaliste mais conditionnée à un effort commun.
\item \textbf{Une recommandation finale mémorable} : \eng{A clean city is a healthy city.} Cette formule courte et symétrique fonctionne comme un slogan : elle reste dans la tête du correcteur. C'est une excellente manière de terminer.
\end{enumerate}

\begin{attention}[Ce qu'il ne faut PAS faire dans une conclusion]
\begin{itemize}
\item Ne pas introduire une idée totalement nouvelle dans la conclusion : elle doit \emph{conclure}, pas ouvrir un nouveau débat.
\item Ne pas recopier l'introduction mot pour mot : reformulez avec d'autres termes.
\item Ne pas terminer par \eng{that is all} ou \eng{I have finished} : c'est enfantin et hors du genre académique.
\end{itemize}
\end{attention}

\courssub{Analyse des temps et de la voix passive}

Le modèle utilise trois choix grammaticaux typiques de ce type de sujet :

\begin{center}
\begin{tabularx}{\linewidth}{|X|X|X|}
\hline
\rowcolor{popblueL} Choix grammatical & Exemple du modèle & Pourquoi ce choix ? \\
\hline
Présent simple & \eng{garbage is not collected regularly} & décrire une situation générale, une habitude \\
\hline
Voix passive (présent) & \eng{can be seen}, \eng{can be recycled}, \eng{must be educated} & l'important est l'action, pas celui qui la fait \\
\hline
Modaux \eng{should} / \eng{must} / \eng{can} & \eng{the city council should provide} & exprimer la recommandation et la possibilité \\
\hline
Present perfect & \eng{cooperatives have created jobs} & un résultat récent visible aujourd'hui \\
\hline
Futur avec \eng{will} & \eng{this essay will first describe} & annoncer le plan dans l'introduction \\
\hline
\end{tabularx}
\end{center}

\begin{propriete}[La voix passive, reine des sujets de société]
Forme : \textbf{be} (au temps voulu) + \textbf{participe passé}. Avec un modal : \textbf{modal + be + participe passé}.

\eng{Plastic can be recycled.} « Le plastique peut être recyclé. »

\eng{Rubbish is thrown everywhere.} « Les ordures sont jetées partout. »

Emploi : on utilise le passif quand l'agent est inconnu, évident ou sans importance. C'est exactement le cas dans les compositions sur l'environnement : on s'intéresse aux \emph{actions} (collecter, trier, recycler), pas aux personnes.

Comparaison avec le français : le français préfère souvent « on » (\emph{on jette les ordures partout}) ; l'anglais préfère le passif (\eng{rubbish is thrown everywhere}). Traduire « on » par \eng{people} ou \eng{they} partout est une erreur de style.
\end{propriete}

\begin{attention}[Pièges fréquents des francophones]
\begin{itemize}
\item \eng{the garbages} : faux ! \eng{garbage} et \eng{rubbish} sont \emph{indénombrables} : jamais de \eng{-s}, jamais de \eng{a}. Dites \eng{some garbage}, \eng{a lot of rubbish}.
\item \eng{informations}, \eng{advices} : même piège, ces mots sont indénombrables en anglais.
\item \eng{to sensibilize citizens} : ce verbe n'existe pas ! Dites \eng{to educate citizens}, \eng{to raise awareness among citizens}, \eng{to sensitise} (britannique, avec \eng{-ise}).
\item \eng{the town hall must to provide} : après un modal, la base verbale \emph{sans} \eng{to} : \eng{must provide}.
\end{itemize}
\end{attention}

\courssub{Méthode : exploiter un modèle sans le copier}

\begin{methode}[Comment travailler avec un modèle de composition]
\begin{enumerate}
\item \textbf{Lire pour le sens} : une première lecture globale, sans dictionnaire, pour saisir le message général.
\item \textbf{Identifier la structure} : numérotez les paragraphes et donnez à chacun une étiquette (introduction, problème, solutions, conclusion).
\item \textbf{Relever les connecteurs} : soulignez-les et classez-les (addition, opposition, cause, conséquence, conclusion).
\item \textbf{Relever le vocabulaire thématique} : constituez votre propre liste personnelle (ici : \eng{garbage, dustbin, rubbish, waste}).
\end{enumerate}
\end{methode}


\coursec{Erreurs à éviter et grille d'évaluation}

Vous avez maintenant un modèle de composition commenté sous les yeux. Mais un bon modèle ne suffit pas : au baccalauréat, la différence entre une note moyenne et une excellente note se joue souvent sur la \cle{relecture} et sur l'élimination des erreurs « bêtes » qui coûtent cher. Cette dernière partie vous donne d'abord la liste noire des erreurs les plus fréquentes des candidats francophones sur ce sujet, puis une stratégie de relecture professionnelle, et enfin la grille d'évaluation utilisée par les correcteurs — pour que vous sachiez exactement comment votre copie sera notée.

\courssub{Les erreurs de lexique à bannir}

Commençons par observer quatre phrases extraites de copies réelles d'élèves. Saurez-vous repérer l'erreur dans chacune ?

\begin{exemplebox}[Observez et corrigez]
\begin{enumerate}
\item \eng{The garbages are everywhere in our quarter.}
\item \eng{Children throw the garbage on the floor in the street.}
\item \eng{We waste a lot of waste every day.}
\item \eng{The government should to provide more dustbins.}
\end{enumerate}
\end{exemplebox}

\begin{corrige}[Observation]
\begin{enumerate}
\item \eng{Garbage is everywhere in our quarter.} — \eng{garbage} est indénombrable : pas de pluriel, verbe au singulier.
\item \eng{Children throw garbage on the ground in the street.} — \eng{floor} = le sol d'une maison ; dehors, on dit \eng{the ground}.
\item \eng{We waste a lot of food every day.} ou \eng{We produce a lot of waste.} — ne pas confondre le nom et le verbe.
\item \eng{The government should provide more dustbins.} — après \eng{should}, infinitif sans \eng{to}.
\end{enumerate}
\end{corrige}

\begin{propriete}[Waste, garbage, rubbish : les mots des déchets]
\begin{itemize}
\item \cle{garbage} (américain) et \cle{rubbish} (britannique) : \emph{indénombrables}. Jamais de \eng{-s}, verbe toujours au \textbf{singulier}. \eng{Rubbish is collected twice a week.} « Les ordures sont ramassées deux fois par semaine. »
\item \cle{waste} (nom) : « les déchets », indénombrable. \eng{Plastic waste pollutes our rivers.} « Les déchets plastiques polluent nos rivières. »
\item \cle{to waste} (verbe) : « gaspiller ». \eng{Do not waste water.} « Ne gaspillez pas l'eau. »
\item \cle{litter} : « les détritus » abandonnés dans la rue (indénombrable) ; \cle{to litter} : « jeter des détritus ». \eng{Do not drop litter on the ground.} « Ne jetez pas de détritus par terre. »
\item \cle{ground} : le sol dehors ; \cle{floor} : le sol à l'intérieur.
\end{itemize}
\end{propriete}

\begin{attention}[Faux amis et confusions fréquentes]
\begin{itemize}
\item \eng{waste} (nom) = déchets ; \eng{to waste} (verbe) = gaspiller. \eng{Throwing away good food is a waste.} « Jeter de la nourriture encore bonne est un gaspillage. »
\item \eng{to throw away} = jeter (aux ordures) ; \eng{to throw up} = vomir ! Attention au sens : la particule change tout dans les phrasal verbs.
\item \eng{a bin / a dustbin} = une poubelle ; on ne dit pas \eng{a trash} pour l'objet (\eng{trash} = les déchets, indénombrable, surtout américain).
\item \eng{to recycle} (verbe) / \eng{recycling} (nom) : \eng{We must recycle.} « Nous devons recycler. » / \eng{Recycling is important.} « Le recyclage est important. »
\item \eng{to sort} = trier ; \eng{sorting} = le tri. \eng{We should sort our waste at home.} « Nous devrions trier nos déchets à la maison. »
\end{itemize}
\end{attention}

\courssub{Les erreurs de grammaire qui coûtent des points}

\begin{attention}[Après should, must, can : infinitif SANS to]
Les auxiliaires modaux \eng{should}, \eng{must}, \eng{can}, \eng{may} sont toujours suivis de la \textbf{base verbale}, jamais de \eng{to}.\par
Faux : \eng{The government should to provide more trucks.}\par
Juste : \eng{The government should provide more trucks.} « Le gouvernement devrait fournir plus de camions. »\par
Faux : \eng{We must recycling our waste.}\par
Juste : \eng{We must recycle our waste.} « Nous devons recycler nos déchets. »\par
C'est l'erreur numéro un des francophones, car en français on dit « devrait \emph{fournir} » avec l'infinitif complet, et la tentation de calquer avec \eng{to} est permanente.
\end{attention}

\begin{attention}[People est déjà pluriel]
\eng{People} signifie « les gens » et se construit toujours avec un verbe au \textbf{pluriel}, sans \eng{-s} : \eng{People throw garbage everywhere.} « Les gens jettent des ordures partout. »\par
La forme \eng{peoples} existe mais signifie « les peuples » (les nations) : \eng{the peoples of Africa} « les peuples d'Afrique ». Ne l'utilisez jamais pour dire « les gens ».\par
Même vigilance avec \eng{information} : indénombrable en anglais. Faux : \eng{informations}. Juste : \eng{The council gives useful information about recycling.} « La municipalité donne des informations utiles sur le recyclage. »
\end{attention}

\begin{propriete}[Le -s de la 3e personne du singulier : la règle d'or]
Au présent simple, à la 3\up{e} personne du singulier (\eng{he, she, it} ou un nom singulier comme \eng{the council, the government}), le verbe prend un \eng{-s}.\par
Faux : \eng{The council collect the garbage every Monday.}\par
Juste : \eng{The council collects the garbage every Monday.} « La municipalité ramasse les ordures chaque lundi. »\par
Faux : \eng{Garbage are everywhere.}\par
Juste : \eng{Garbage is everywhere.} « Les ordures sont partout. » (sujet indénombrable = singulier)\par
Astuce de relecture : soulignez chaque sujet singulier et vérifiez que son verbe porte bien son \eng{-s}.
\end{propriete}

\begin{propriete}[Once, twice, three times : compter les fréquences]
Pour exprimer la fréquence, l'anglais n'utilise pas \eng{one time} ou \eng{two times} mais des mots spéciaux :
\begin{itemize}
\item \cle{once} = une fois : \eng{The truck comes once a week.} « Le camion passe une fois par semaine. »
\item \cle{twice} = deux fois : \eng{Rubbish is collected twice a month.} « Les ordures sont ramassées deux fois par mois. »
\item \cle{three times}, \cle{four times}... = trois fois, quatre fois (à partir de trois, on utilise \eng{times}).
\end{itemize}
\end{propriete}

\begin{center}
\begin{tabularx}{\linewidth}{|X|X|X|}
\hline
\rowcolor{popblueL} Erreur fréquente (FAUX) & Correction (JUSTE) & Traduction \\
\hline
\eng{the garbages} & \eng{garbage} & les ordures (indénombrable) \\
\hline
\eng{on the floor} (dehors) & \eng{on the ground} & par terre \\
\hline
\eng{peoples throw} & \eng{people throw} & les gens jettent \\
\hline
\eng{informations} & \eng{information} & des informations \\
\hline
\eng{should to provide} & \eng{should provide} & devrait fournir \\
\hline
\eng{the council collect} & \eng{the council collects} & la municipalité ramasse \\
\hline
\eng{we must recycling} & \eng{we must recycle} & nous devons recycler \\
\hline
\eng{the garbage are} & \eng{garbage is} & les ordures sont \\
\hline
\eng{one time a week} & \eng{once a week} & une fois par semaine \\
\hline
\end{tabularx}
\end{center}

\courssub{Les erreurs de structure, de ponctuation et de majuscules}

Au-delà des mots, les correcteurs sanctionnent sévèrement les défauts d'organisation. Voici les quatre fautes de structure à éviter absolument :

\begin{itemize}
\item \textbf{Pas de paragraphes} : un seul bloc de texte donne l'impression d'un désordre total. Une composition = au minimum trois paragraphes visibles (introduction, développement, conclusion), séparés par un saut de ligne ou un alinéa.
\item \textbf{Pas d'introduction} : commencer directement par \eng{Firstly, garbage is dangerous...} sans phrase d'accroche ni annonce de plan prive le lecteur de tout repère.
\item \textbf{Idées répétées} : dire deux fois la même chose avec des mots différents ne fait pas deux arguments. Chaque paragraphe du développement doit apporter une idée \emph{nouvelle} (problème, causes, solutions).
\item \textbf{Hors sujet} : si le sujet porte sur \emph{recycling services}, ne rédigez pas une page entière sur la pollution de l'air. Relisez le sujet après chaque paragraphe et demandez-vous : « Est-ce que je réponds encore à la question ? »
\end{itemize}

Côté ponctuation et majuscules, deux erreurs « de francophone » reviennent sans cesse :

\begin{exemplebox}[Ponctuation et majuscules]
Faux : \eng{i think recycling is important because it protect the environment}\par
Juste : \eng{I think recycling is important because it protects the environment.}\par
« Je pense que le recyclage est important parce qu'il protège l'environnement. »\par
Règles : le pronom \eng{I} prend \textbf{toujours} une majuscule, même au milieu de la phrase ; chaque phrase se termine par un point ; chaque phrase nouvelle commence par une majuscule.
\end{exemplebox}

\courssub{La relecture professionnelle : trois passes et la méthode C.O.P.S.}

Relire ne veut pas dire « regarder sa copie une dernière fois ». Une relecture efficace se fait en \textbf{trois passes séparées}, chacune avec un objectif unique :

\begin{enumerate}
\item \textbf{Passe 1 — Sens et organisation} : lisez votre texte comme un étranger. L'introduction annonce-t-elle le sujet ? Chaque paragraphe a-t-il une idée claire ? Les connecteurs (\eng{firstly, moreover, however, to conclude}) sont-ils présents ? Y a-t-il du hors sujet ?
\item \textbf{Passe 2 — Grammaire et temps} : vérifiez chaque verbe : accord avec le sujet (\eng{-s} à la 3\up{e} personne), modal suivi de la base verbale, cohérence des temps (présent simple pour les vérités générales, futur ou conditionnel pour les solutions).
\item \textbf{Passe 3 — Orthographe et ponctuation} : chassez les fautes de mots (\eng{environment}, \eng{government}, \eng{recycling}), les points manquants, les majuscules oubliées.
\end{enumerate}

\begin{methode}[La méthode C.O.P.S. pour ne rien oublier]
Retenez l'acronyme anglais \cle{C.O.P.S.} et vérifiez chaque lettre séparément avant de rendre votre copie :
\begin{enumerate}
\item \textbf{C — Capitalization} (majuscules) : début de phrase, \eng{I}, noms propres (\eng{Yaounde, Cameroon}).
\item \textbf{O — Organization} (organisation) : introduction, paragraphes, connecteurs, conclusion.
\item \textbf{P — Punctuation} (ponctuation) : un point à chaque phrase, des virgules après les connecteurs en début de phrase (\eng{Firstly, ...}).
\item \textbf{S — Spelling} (orthographe) : mots du vocabulaire du thème et verbes irréguliers.
\end{enumerate}
Consacrez les cinq dernières minutes de l'épreuve exclusivement à cette vérification : elle peut vous faire gagner deux points.
\end{methode}

\begin{exoresolu}[Corrigez ce paragraphe d'élève]
Énoncé — Voici le début d'une copie. Relevez les six erreurs et corrigez-les.\par
\eng{In our town, the garbages is a big problem. Peoples throw waste on the floor and the council collect it only one time a week. i think the government should to build a recycling centre.}\tcblower
Solution détaillée :
\begin{enumerate}
\item \eng{the garbages is} $\rightarrow$ \eng{garbage is} (indénombrable, sans pluriel).
\item \eng{Peoples throw} $\rightarrow$ \eng{People throw} (déjà pluriel).
\item \eng{on the floor} $\rightarrow$ \eng{on the ground} (à l'extérieur).
\item \eng{the council collect} $\rightarrow$ \eng{the council collects} (\eng{-s} de la 3\up{e} personne).
\item \eng{only one time} $\rightarrow$ \eng{only once} (« une fois » se dit \eng{once}).
\item \eng{i think} $\rightarrow$ \eng{I think} (majuscule obligatoire) ; \eng{should to build} $\rightarrow$ \eng{should build} (base verbale après modal).
\end{enumerate}
Version corrigée : \eng{In our town, garbage is a big problem. People throw waste on the ground and the council collects it only once a week. I think the government should build a recycling centre.}\par
« Dans notre ville, les ordures sont un gros problème. Les gens jettent des déchets par terre et la municipalité ne les ramasse qu'une fois par semaine. Je pense que le gouvernement devrait construire un centre de recyclage. »
\end{exoresolu}

\courssub{La grille d'évaluation : comment votre copie sera notée}

Au baccalauréat, l'expression écrite est notée sur des critères précis. Connaître cette grille, c'est connaître les attentes du correcteur : vous pouvez rédiger \emph{pour} ces critères.

\begin{propriete}[Les quatre critères de la grille d'évaluation (20 points)]
\begin{itemize}
\item \cle{Content} (contenu et pertinence) — 5 pts : le sujet est traité, les idées sont pertinentes, développées et illustrées d'exemples.
\item \cle{Organisation} — 5 pts : trois parties visibles (introduction, développement, conclusion), paragraphes clairs, connecteurs logiques variés.
\item \cle{Vocabulary} (vocabulaire) — 5 pts : lexique du thème riche et précis (\eng{waste, dustbin, landfill, to recycle, to sort, collection}), pas de répétitions.
\item \cle{Accuracy} (correction de la langue) — 5 pts : grammaire (accords, temps, modaux), orthographe, ponctuation, majuscules.
\end{itemize}
\end{propriete}

\begin{popfigure}
\begin{tikzpicture}[scale=1]
\fill[popblue] (0,0) rectangle (12,1);
\node[text=white, font=\bfseries] at (2.6,0.5) {Criterion};
\node[text=white, font=\bfseries] at (5.6,0.5) {Excellent (4--5)};
\node[text=white, font=\bfseries] at (8.6,0.5) {Good (2.5--3.5)};
\node[text=white, font=\bfseries] at (11,0.5) {Weak (0--2)};
\fill[popblueL] (0,-1.4) rectangle (12,0);
\node[text=popdark, font=\bfseries, align=center] at (2.6,-0.7) {Content};
\node[text=popdark, align=center, font=\footnotesize] at (5.6,-0.7) {Topic fully treated,\\relevant ideas, examples};
\node[text=popdark, align=center, font=\footnotesize] at (8.6,-0.7) {Topic treated,\\few ideas developed};
\node[text=popdark, align=center, font=\footnotesize] at (11,-0.7) {Off topic or\\very poor ideas};
\fill[white] (0,-2.8) rectangle (12,-1.4);
\node[text=popdark, font=\bfseries, align=center] at (2.6,-2.1) {Organisation};
\node[text=popdark, align=center, font=\footnotesize] at (5.6,-2.1) {3 clear parts,\\varied linkers};
\node[text=popdark, align=center, font=\footnotesize] at (8.6,-2.1) {Parts visible,\\few linkers};
\node[text=popdark, align=center, font=\footnotesize] at (11,-2.1) {No paragraphs,\\no linkers};
\fill[popblueL] (0,-4.2) rectangle (12,-2.8);
\node[text=popdark, font=\bfseries, align=center] at (2.6,-3.5) {Vocabulary};
\node[text=popdark, align=center, font=\footnotesize] at (5.6,-3.5) {Rich, precise\\theme lexis};
\node[text=popdark, align=center, font=\footnotesize] at (8.6,-3.5) {Correct but\\repetitive};
\node[text=popdark, align=center, font=\footnotesize] at (11,-3.5) {Very limited,\\wrong words};
\fill[white] (0,-5.6) rectangle (12,-4.2);
\node[text=popdark, font=\bfseries, align=center] at (2.6,-4.9) {Accuracy};
\node[text=popdark, align=center, font=\footnotesize] at (5.6,-4.9) {Almost no errors\\of grammar, spelling};
\node[text=popdark, align=center, font=\footnotesize] at (8.6,-4.9) {Errors but\\meaning clear};
\node[text=popdark, align=center, font=\footnotesize] at (11,-4.9) {Frequent errors,\\hard to understand};
\draw[popline] (0,0) rectangle (12,-5.6);
\draw[popline] (4.2,0) -- (4.2,-5.6);
\draw[popline] (7.2,0) -- (7.2,-5.6);
\draw[popline] (9.8,0) -- (9.8,-5.6);
\draw[popline] (0,-1.4) -- (12,-1.4);
\draw[popline] (0,-2.8) -- (12,-2.8);
\draw[popline] (0,-4.2) -- (12,-4.2);
\end{tikzpicture}
\legende{Grille d'évaluation type de la composition : quatre critères, cinq points chacun.}
\end{popfigure}

\begin{savaistu}[Pourquoi « accuracy » et pas « correction » ?]
Dans les grilles officielles d'expression écrite, le dernier critère s'appelle \eng{accuracy} (« exactitude, correction ») et non \eng{correction}. Le mot anglais \eng{accuracy} vient du latin \emph{accurare}, « prendre soin de » : le correcteur ne cherche pas la perfection absolue, mais une langue \emph{soignée}. Quelques erreurs isolées ne vous font pas perdre tous vos points si le sens reste clair — c'est l'accumulation des fautes qui est sanctionnée. D'où l'importance de la relecture : éliminer les erreurs faciles à repérer suffit souvent à passer de « Weak » à « Good ».
\end{savaistu}

\courssub{Auto-évaluation et évaluation par les pairs}

La grille n'est pas réservée aux correcteurs : utilisez-la \emph{avant} de rendre votre copie.

\begin{experience}[Évaluation par les pairs en classe]
\begin{enumerate}
\item Rédigez votre composition au brouillon, puis effectuez vos trois passes de relecture (méthode C.O.P.S.).
\item Échangez votre brouillon avec un voisin. Chacun note la copie de l'autre avec la grille : un score sur 5 pour chaque critère, plus \textbf{un point fort} et \textbf{un conseil d'amélioration} écrits en anglais, par exemple \eng{Good linkers, but check the -s of the third person.}
\item Récupérez votre brouillon, lisez les remarques, puis rédigez la \textbf{copie finale} en intégrant les corrections.
\item Auto-évaluation finale : posez-vous les quatre questions — \eng{Did I answer the question? Did I organize my ideas? Did I use rich vocabulary? Did I check my grammar and spelling?} Si une réponse est « no », corrigez avant de rendre.
\end{enumerate}
\end{experience}

\begin{exercice}[À vous de jouer]
Voici trois phrases de copies. Corrigez chacune, puis attribuez-vous une note sur 5 au critère \emph{Accuracy} pour l'ensemble.
\begin{enumerate}
\item \eng{The government must to sensitize the peoples about recycling.}
\item \eng{Garbage collection are very important for our healths.}
\item \eng{i think we should sorting our waste at home}
\end{enumerate}
\end{exercice}

\begin{corrige}[Exercice — À vous de jouer]
\begin{enumerate}
\item \eng{The government must sensitize the people about recycling.} — modal + base verbale ; \eng{people} sans \eng{-s}. « Le gouvernement doit sensibiliser les gens au recyclage. »
\item \eng{Garbage collection is very important for our health.} — sujet singulier ; \eng{health} est indénombrable. « Le ramassage des ordures est très important pour notre santé. »
\item \eng{I think we should sort our waste at home.} — majuscule à \eng{I} ; \eng{should} + base verbale ; point final. « Je pense que nous devrions trier nos déchets à la maison. »
\end{enumerate}
\end{corrige}

\begin{aretenir}[L'essentiel de la section]
\begin{itemize}
\item Indénombrables à surveiller : \eng{garbage, rubbish, waste, litter, information, health} — jamais de \eng{-s}, verbe au singulier.
\item Après \eng{should, must, can} : base verbale \textbf{sans} \eng{to}.
\item \eng{People} = pluriel (« les gens ») ; \eng{I} prend toujours une majuscule ; « une fois » = \eng{once}, « deux fois » = \eng{twice}.
\item Relecture en trois passes, guidée par la méthode \cle{C.O.P.S.} (Capitalization, Organization, Punctuation, Spelling).
\item Grille d'évaluation : \emph{Content, Organisation, Vocabulary, Accuracy} — 5 points chacun. Évaluez-vous et faites évaluer votre brouillon par un pair avant la copie finale.
\end{itemize}
\end{aretenir}

\newpage

\coursec{Activité d'intégration}

\courssub{Objectifs de l'activité}

À l'issue de cette activité d'intégration, tu seras capable de :
\begin{itemize}
\item analyser un sujet de composition (mots-clés, contraintes de longueur, type de texte attendu) ;
\item organiser tes idées en un plan clair : introduction, développement en trois paragraphes, conclusion ;
\item rédiger une composition de 250 à 300 mots sur l'amélioration des services de \cle{garbage collection}, \cle{disposal} et \cle{recycling} dans ta ville ;
\item utiliser correctement les connecteurs logiques (\eng{firstly, moreover, consequently, in conclusion}), les temps verbaux adaptés (présent simple, present perfect, futur avec \eng{will} et \eng{be going to}) et le vocabulaire thématique de la leçon ;
\item relire, corriger et évaluer la production d'un camarade à l'aide d'une grille d'évaluation.
\end{itemize}

\courssub{Situation}

Tu es élève en Terminale A au lycée de Buea, dans la région du Sud-Ouest du Cameroun. Depuis plusieurs mois, les habitants de ton quartier se plaignent : les poubelles débordent, les ordures s'accumulent au bord des routes, les caniveaux sont bouchés et les pluies provoquent des inondations. Le club d'anglais de ton lycée (\emph{English Language Club}) a décidé de réagir. À l'occasion de la \emph{World Environment Day}, le club organise un concours de rédaction. La meilleure composition sera lue à la radio locale et envoyée au maire de la ville (\emph{the Mayor}) avec une lettre officielle.

Voici le sujet du concours, tel qu'il a été affiché au tableau du club :

\begin{textebox}[English Language Club — Essay Competition]
\textbf{Topic:} Write a composition (250--300 words) on how garbage collection, disposal and recycling services can be improved in your town.

In your composition, you should:
\begin{itemize}
\item describe the current waste management situation in your town;
\item explain the causes of the problem and its effects on people's health and the environment;
\item suggest practical solutions, including recycling, and say what the council and the population should do.
\end{itemize}

The best essay will be read on Ocean City Radio and forwarded to the Mayor of Buea. Deadline: Friday 14 June. Good luck!
\end{textebox}

Ton professeur d'anglais t'encourage à participer. Tu vas donc rédiger ta composition en suivant toutes les étapes de la leçon : analyse du sujet, brainstorming, plan, rédaction, relecture.

\courssub{Consignes}

\textbf{Étape 1 — Analyse du sujet (travail en groupes de 3, 5 minutes).}
Lis attentivement le sujet du concours. Souligne les mots-clés et réponds aux questions suivantes :
\begin{enumerate}
\item Quel est le thème général de la composition ?
\item Quelle longueur est exigée ? Que risques-tu si tu écris 120 mots ou 500 mots ?
\item Quelles sont les trois tâches précises demandées dans la consigne ?
\item Qui va lire ta composition ? Cela influence-t-il le ton (familier ou formel) ?
\end{enumerate}

\textbf{Étape 2 — Brainstorming (5 minutes).}
Toujours en groupe, note en vrac, en anglais, toutes les idées et tous les mots de vocabulaire qui te viennent à l'esprit sur les déchets dans ta ville : problèmes, causes, conséquences, solutions. Classe ensuite tes idées en trois colonnes : \emph{Problem}, \emph{Causes and effects}, \emph{Solutions}.

\textbf{Étape 3 — Construction du plan commun (10 minutes).}
Avec ton groupe, rédigez un plan détaillé en suivant cette architecture :
\begin{enumerate}
\item \textbf{Introduction :} accroche (une observation sur ta ville), présentation du sujet, annonce du plan.
\item \textbf{Paragraphe 1 :} description du problème actuel (collecte irrégulière, décharges sauvages, caniveaux bouchés).
\item \textbf{Paragraphe 2 :} causes et effets (manque de camions et de poubelles, comportements des habitants ; maladies, pollution, inondations).
\item \textbf{Paragraphe 3 :} solutions (collecte régulière, tri et recyclage, sensibilisation, amendes, implication de la population).
\item \textbf{Conclusion :} résumé bref, appel à l'action, phrase d'ouverture optimiste.
\end{enumerate}

\textbf{Étape 4 — Rédaction individuelle (35 minutes).}
Chaque élève rédige maintenant \emph{sa propre} composition de 250 à 300 mots, en s'appuyant sur le plan du groupe, la banque de connecteurs et le vocabulaire de la leçon. Soigne l'écriture, la ponctuation et les paragraphes. Indique le nombre de mots à la fin.

\textbf{Étape 5 — Correction croisée (15 minutes).}
Échange ta copie avec un groupe voisin. À l'aide de la grille d'évaluation de la leçon (contenu, organisation, vocabulaire, exactitude linguistique), note la copie sur 20 et souligne en vert les réussites, en rouge les erreurs. Rédige un commentaire bienveillant de deux phrases.

\textbf{Étape 6 — Mise en commun (10 minutes).}
Chaque groupe choisit la meilleure introduction et la meilleure conclusion corrigées. Le professeur les écrit au tableau et la classe commente collectivement les points forts.

\courssub{Ressources à mobiliser}

\begin{aretenir}[Vocabulaire utile]
\begin{itemize}
\item \cle{garbage / rubbish / waste} : les déchets (indénombrables) ; \cle{garbage collection} : la collecte des ordures ; \cle{dustbin / trash can} : la poubelle ; \cle{dump} : la décharge sauvage ; \cle{landfill} : la décharge contrôlée.
\item \cle{to dispose of waste} : éliminer les déchets ; \cle{disposal} : l'élimination ; \cle{to recycle} : recycler ; \cle{recycling bin} : bac de tri ; \cle{plastic bottles, cans, paper, glass} : matières recyclables.
\item \cle{to litter} : jeter des ordures (dans la rue) ; \cle{blocked gutters} : caniveaux bouchés ; \cle{floods} : inondations ; \cle{diseases} : maladies ; \cle{to raise awareness} : sensibiliser ; \cle{fine} : amende ; \cle{the council / the mayor} : la mairie / le maire.
\end{itemize}
\end{aretenir}

\begin{propriete}[Connecteurs et temps verbaux]
\begin{itemize}
\item \textbf{Ordonner :} \eng{firstly, secondly, finally} ; \textbf{ajouter :} \eng{moreover, in addition, furthermore} ; \textbf{exprimer la cause :} \eng{because, since, as} ; \textbf{exprimer la conséquence :} \eng{consequently, as a result, therefore} ; \textbf{opposer :} \eng{however, whereas, although} ; \textbf{conclure :} \eng{in conclusion, to sum up}.
\item \textbf{Temps :} le \emph{present simple} pour décrire la situation actuelle (\eng{Garbage piles up in our streets.}) ; le \emph{present perfect} pour un bilan d'expérience jusqu'à aujourd'hui (\eng{The situation has become worse.}) ; \eng{will} et \eng{be going to} pour les solutions et prévisions (\eng{The council will provide more bins.}) ; le modal \eng{should} pour les recommandations (\eng{The mayor should organise weekly collections.}) et \eng{would} pour les conséquences hypothétiques (\eng{Recycling would create jobs.}).
\end{itemize}
\end{propriete}

\begin{methode}[Rédiger une composition en cinq temps]
\begin{enumerate}
\item \textbf{Analyser} le sujet : souligne les mots-clés, repère la longueur exigée et les tâches à accomplir.
\item \textbf{Brainstormer} : note toutes tes idées en vrac, puis classe-les par paragraphe.
\item \textbf{Planifier} : rédige un plan (introduction, trois paragraphes, conclusion) avant d'écrire la première phrase.
\item \textbf{Rédiger} : une idée principale par paragraphe, un connecteur en tête de paragraphe, des phrases complètes (sujet + verbe).
\item \textbf{Relire} : vérifie les accords (s de la 3\up{e} personne), les temps, l'orthographe, la ponctuation et le nombre de mots.
\end{enumerate}
\end{methode}

\begin{attention}[Erreurs à éviter]
Ne traduis pas mot à mot depuis le français : on dit \eng{to throw rubbish away} et non « to throw garbages » (\eng{garbage} est indénombrable). Attention à \eng{informations} et \eng{advices} : ce sont des faux pluriels ; écris \eng{information} et \eng{advice}. N'oublie pas le \textbf{s} de la troisième personne au présent simple : \eng{the council collects}, \eng{recycling protects}. Évite les phrases sans verbe et les paragraphes d'une seule ligne. Enfin, \eng{to sensitize/sensitise} existe, mais \eng{to raise awareness} est plus naturel pour « sensibiliser ».
\end{attention}

\courssub{Proposition de production et corrigé commenté}

\begin{exoresolu}[Modèle de composition : “Clean Our City!”]
\textbf{Énoncé.} Write a composition (250--300 words) on how garbage collection, disposal and recycling services can be improved in your town.
\tcblower
\textbf{Production modèle (environ 280 mots) :}

\eng{Every morning, the inhabitants of Buea wake up to the same sad sight: overflowing dustbins, plastic bags flying in the wind and gutters blocked with rubbish. Waste management has become one of the most serious problems in our town, and it is high time we found solutions.}

\eng{Firstly, the current situation is alarming. Garbage is collected irregularly, sometimes only once a week, so mountains of waste pile up in the streets and markets. Moreover, many residents simply dump their rubbish in open spaces or in streams, because there are not enough public bins.}

\eng{This situation has worrying causes and effects. The council does not have enough trucks and workers, and the population has never been properly sensitised about hygiene. As a result, blocked gutters cause floods during the rainy season, and dirty water spreads diseases such as cholera and malaria. Our beautiful town, the former capital of West Cameroon, is slowly becoming unhealthy.}

\eng{However, practical solutions exist. The council should organise regular collections, at least twice a week, and place more dustbins in busy areas. In addition, a recycling programme could be launched: plastic bottles, cans and paper should be sorted and sold to recycling companies, which would create jobs for young people. Finally, people who litter should pay a fine, and schools and radio stations should raise awareness about keeping the environment clean.}

\eng{In conclusion, improving garbage collection, disposal and recycling is the duty of both the council and the population. If everyone plays their part, Buea will become a clean, healthy and attractive town again.}

\textbf{Commentaire des choix de langue :}
\begin{itemize}
\item \textbf{Introduction :} une accroche concrète (\eng{Every morning...}) situe le problème ; l'expression \eng{it is high time we found} (suivi du prétérit) annonce le propos avec élégance.
\item \textbf{Organisation :} cinq paragraphes bien distincts, chacun introduit par un connecteur visible (\eng{Firstly, Moreover, As a result, However, In addition, Finally, In conclusion}).
\item \textbf{Temps verbaux :} présent simple pour les faits (\eng{is collected, pile up}), modal \eng{should} pour les recommandations, \eng{would} pour les conséquences hypothétiques (\eng{would create jobs}), futur \eng{will become} pour la conclusion optimiste.
\item \textbf{Vocabulaire :} lexique précis de la leçon (\eng{overflowing dustbins, waste management, recycling programme, to raise awareness, to litter, fine}) et touche locale (\eng{Buea, the rainy season, cholera and malaria}) qui ancre la composition dans le réel.
\item \textbf{Exactitude :} \eng{garbage} et \eng{rubbish} restent au singulier ; le \textbf{s} de la troisième personne est respecté (\eng{spreads, causes, plays}).
\end{itemize}
\end{exoresolu}

\courssub{Grille d'évaluation de la composition}

Utilise cette grille pour la correction croisée (étape 5) : elle reprend les quatre critères de l'épreuve d'expression écrite du baccalauréat.

\begin{center}
\begin{tabularx}{\linewidth}{|X|X|c|}
\hline
\rowcolor{popblueL} \textbf{Critère} & \textbf{Ce qui est attendu} & \textbf{Note} \\
\hline
Contenu et respect du sujet & Les trois tâches sont traitées ; les idées sont pertinentes ; la longueur (250--300 mots) est respectée. & /5 \\
\hline
Organisation & Plan visible (introduction, 3 paragraphes, conclusion) ; connecteurs variés et bien placés. & /5 \\
\hline
Vocabulaire & Lexique riche et précis du thème (waste, disposal, recycling) ; pas de répétitions excessives. & /5 \\
\hline
Exactitude linguistique & Grammaire, temps verbaux, orthographe, ponctuation ; phrases complètes. & /5 \\
\hline
\rowcolor{popgreenL} \textbf{Total} & & \textbf{/20} \\
\hline
\end{tabularx}
\end{center}

\courssub{Bilan de l'activité}

Cette activité t'a permis de mobiliser l'ensemble des savoir-faire de la leçon dans une tâche authentique : analyser un sujet avant d'écrire, transformer un brainstorming en plan structuré, rédiger une composition complète respectant une longueur imposée, puis évaluer objectivement un texte à l'aide d'une grille. Tu as également découvert qu'une bonne composition repose sur trois piliers : un \textbf{plan} clair, des \textbf{connecteurs} qui guident le lecteur et une \textbf{langue exacte}. Ces compétences te serviront à l'examen du baccalauréat, où l'expression écrite est notée sur le contenu, l'organisation, le vocabulaire et l'exactitude — exactement les quatre critères de ta grille. Garde en tête la règle d'or du rédacteur : \eng{Plan first, write second, revise always.}

\newpage

\coursec{Méthodes et exercices résolus}

\coursec{Lesson 36 — Writing: Writes compositions on garbage collection, disposal, and recycling services}

\courssub{Savoir-faire 1 : Planifier et rédiger une composition sur le thème}

\begin{methode}[Planifier et rédiger une composition sur garbage collection, disposal and recycling]
Une composition réussie se prépare \textbf{AVANT} d'écrire. Voici la démarche complète en cinq étapes.
\begin{enumerate}
\item \textbf{Lire le sujet deux fois} et souligner les mots-clés : le thème (\eng{garbage collection}, \eng{waste disposal}, \eng{recycling}), le type de texte demandé (\eng{essay}, \eng{article}, \eng{letter}), le nombre de mots (souvent 250--300 words au Bac).
\item \textbf{Faire un brainstorming} : noter en vrac, en anglais, toutes les idées et tout le vocabulaire liés au sujet (\eng{dustbins}, \eng{landfills}, \eng{plastic bottles}, \eng{to sort waste}, \eng{to pollute}, \eng{to protect the environment}...).
\item \textbf{Classer les idées en trois parties} : ce qui pose le problème (introduction), les arguments et exemples (développement, 2 ou 3 paragraphes), la solution ou l'opinion personnelle (conclusion).
\item \textbf{Rédiger au brouillon} en suivant le plan : une idée principale par paragraphe, introduite par un connecteur (\eng{First}, \eng{Moreover}, \eng{Finally}...). Utiliser des temps cohérents : \cle{present simple} pour les faits généraux, \cle{present perfect} pour les bilans, \cle{futur} ou \cle{conditionnel} pour les propositions.
\item \textbf{Relire et corriger} : vérifier l'accord sujet-verbe, les temps, l'orthographe (\eng{environment}, \eng{government}, \eng{which}), les prépositions (\eng{to dispose of}, \eng{to depend on}) et compter les mots.
\end{enumerate}
Réflexe d'examen : consacrer environ 10 minutes à la planification, 35 minutes à la rédaction et 10 minutes à la relecture.
\end{methode}

\begin{exoresolu}[Analyse d'un sujet type Bac]
Énoncé — Read the following examination topic and answer the questions.\par
\eng{“In many Cameroonian towns, garbage is not collected regularly and waste is dumped anywhere. Write a composition of about 250 words in which you explain the causes of this problem and suggest solutions, including recycling services.”}\par
1. What type of text is required? 2. What are the two tasks you must complete? 3. Give five English words or expressions you could use in this composition. 4. Propose a three-part plan.
\tcblower
Solution détaillée.\par
1. Type de texte : le verbe \eng{write a composition} indique un \textbf{essay} (essai argumentatif), pas une lettre ni un dialogue. Il faut donc une introduction, un développement et une conclusion, sans formule d'appel ni de politesse.\par
2. Les deux tâches : le sujet contient deux verbes d'instruction — \eng{explain the causes} (expliquer les causes : \eng{irregular collection}, \eng{lack of dustbins}, \eng{population growth}) et \eng{suggest solutions} (proposer des solutions : \eng{regular collection}, \eng{recycling services}, \eng{public awareness}). Oublier l'une des deux tâches fait perdre la moitié des points de contenu.\par
3. Vocabulaire exploitable : \eng{garbage collection} (ramassage des ordures), \eng{to dump waste} (déposer des déchets sauvagement), \eng{landfill} (décharge), \eng{to recycle plastic bottles} (recycler les bouteilles en plastique), \eng{public awareness campaign} (campagne de sensibilisation).\par
4. Plan en trois parties :\par
— \textbf{Introduction} : présenter le problème (\eng{uncollected garbage in towns like Douala and Yaoundé}) et annoncer le plan.\par
— \textbf{Développement} : paragraphe 1 = causes (\eng{few garbage trucks}, \eng{no dustbins in some neighbourhoods}, \eng{carelessness of citizens}) ; paragraphe 2 = solutions (\eng{regular collection services}, \eng{sorting and recycling plastic and metal}, \eng{fining illegal dumping}).\par
— \textbf{Conclusion} : opinion personnelle et appel à l'action (\eng{a clean town is everybody's responsibility}).\par
Conclusion de méthode : deux minutes d'analyse du sujet évitent le hors-sujet, première cause d'échec à l'épreuve d'expression écrite.
\end{exoresolu}

\courssub{Savoir-faire 2 : Organiser les idées avec des connecteurs}

\begin{methode}[Construire la structure introduction — développement — conclusion et enchaîner avec des connecteurs]
\begin{enumerate}
\item \textbf{Introduction (3--4 phrases)} : une phrase d'accroche (fait général ou question), la présentation du problème, l'annonce du plan. Modèle : \eng{Nowadays, waste management has become a serious problem in our towns. This essay will examine the causes of poor garbage collection and suggest possible solutions.}
\item \textbf{Développement (2--3 paragraphes)} : chaque paragraphe commence par une phrase d'idée (\cle{topic sentence}), suivie d'explications et d'un exemple concret (Douala, Buea, a market, a school...).
\item \textbf{Conclusion (2--3 phrases)} : résumé des idées + opinion personnelle ou recommandation. Ne jamais introduire d'idée nouvelle dans la conclusion.
\item \textbf{Placer les connecteurs} au début des phrases ou des paragraphes, suivis d'une virgule.
\end{enumerate}
Connecteurs à retenir :
\begin{center}
\begin{tabularx}{\linewidth}{|X|X|X|}
\hline
\rowcolor{popblueL} Fonction & English & Français \\
\hline
Ajout & First of all, Moreover, Furthermore, In addition & D'abord, De plus, En outre \\
\hline
Cause & Because of, Due to, As, Since & À cause de, Comme, Puisque \\
\hline
Conséquence & Therefore, As a result, Consequently & Donc, Par conséquent \\
\hline
Opposition & However, Whereas, On the other hand, Although & Cependant, Tandis que, Bien que \\
\hline
Exemple & For example, For instance, Such as & Par exemple, Tel que \\
\hline
Conclusion & To conclude, In conclusion, All in all & Pour conclure, En somme \\
\hline
\end{tabularx}
\end{center}
\end{methode}

\begin{exoresolu}[Relier des phrases avec le bon connecteur]
Énoncé — Complete each sentence with a suitable linker from the box: \eng{However — Moreover — As a result — For instance — To conclude}.\par
1. \eng{Garbage is not collected regularly in our neighbourhood. \dots, heaps of rubbish can be seen near the market.}\par
2. \eng{Recycling creates jobs for young people. \dots, it reduces the amount of waste in landfills.}\par
3. \eng{Many materials can be recycled. \dots, plastic bottles, metal cans and old newspapers.}\par
4. \eng{The city council has bought new garbage trucks. \dots, the collection service is still irregular.}\par
5. \eng{\dots, everybody must take part in keeping our town clean.}
\tcblower
Solution détaillée.\par
1. \eng{As a result} : la deuxième phrase exprime la \textbf{conséquence} de la première (pas de ramassage $\rightarrow$ tas d'ordures). \eng{As a result} est suivi d'une virgule.\par
2. \eng{Moreover} : on \textbf{ajoute} un deuxième avantage du recyclage (emplois + réduction des déchets). \eng{Furthermore} ou \eng{In addition} seraient aussi corrects.\par
3. \eng{For instance} : on donne des \textbf{exemples} de matériaux recyclables. Attention : \eng{such as} serait possible mais sans virgule et sans point avant ; ici la structure avec point exige \eng{For instance}.\par
4. \eng{However} : il y a \textbf{opposition} entre l'achat de camions et le service toujours irrégulier. Ne pas confondre avec \eng{therefore} (conséquence), qui renverserait le sens.\par
5. \eng{To conclude} : phrase finale de recommandation, typique d'une conclusion.\par
Conclusion : le choix du connecteur dépend de la relation logique entre les deux idées (cause, conséquence, ajout, opposition, exemple) ; toujours se demander « quel est le rapport entre ces deux phrases ? » avant de choisir.
\end{exoresolu}

\courssub{Savoir-faire 3 : Employer les bons temps verbaux et le bon vocabulaire}

\begin{methode}[Choisir les temps et le lexique adaptés au thème de l'environnement]
\begin{enumerate}
\item \textbf{Present simple} : faits généraux, habitudes, vérités. \eng{Garbage collection takes place twice a week in our district.} « Le ramassage des ordures a lieu deux fois par semaine dans notre quartier. »
\item \textbf{Present continuous} : situation en cours, évolution. \eng{More and more people are throwing plastic bags into rivers.} « De plus en plus de gens jettent des sachets plastiques dans les rivières. »
\item \textbf{Present perfect} : bilan d'actions récentes avec résultat présent. \eng{The council has recently created a recycling centre.} « La mairie a récemment créé un centre de recyclage. »
\item \textbf{Futur (will / be going to)} et \textbf{conditionnel (should, must, could)} : prédictions et recommandations. \eng{The government should invest in recycling services.} « Le gouvernement devrait investir dans les services de recyclage. »
\item \textbf{Voix passive} : très utile pour ce thème (on s'intéresse à l'action, pas à l'auteur). \eng{Tons of waste are dumped in landfills every day.} « Des tonnes de déchets sont déversées dans des décharges chaque jour. »
\end{enumerate}
Lexique-clé : \cle{garbage}/\cle{rubbish}/\cle{waste} (déchets), \cle{dustbin}/\cle{bin} (poubelle), \cle{garbage truck} (camion à ordures), \cle{landfill} (décharge), \cle{to sort waste} (trier les déchets), \cle{to recycle} (recycler), \cle{recycling bin} (bac de tri), \cle{reusable} (réutilisable), \cle{to pollute} / \cle{pollution} (polluer / pollution), \cle{environmentally friendly} (écologique).
\end{methode}

\begin{exoresolu}[Temps verbaux et voix passive]
Énoncé — Put the verbs in brackets in the correct tense or form.\par
1. \eng{Every Monday, a garbage truck (collect) the waste in our street.}\par
2. \eng{Look! The workers (empty) the dustbins near the stadium.}\par
3. \eng{Since last year, our school (recycle) more than two tons of paper.}\par
4. \eng{If nothing is done, the pollution of our rivers (increase).}\par
5. \eng{Most plastic bags (throw) away after a single use.}
\tcblower
Solution détaillée.\par
1. \eng{collects} : present simple, car \eng{every Monday} indique une habitude ; à la 3\textsuperscript{e} personne du singulier, on ajoute -s (\eng{a garbage truck collects}).\par
2. \eng{are emptying} : present continuous, signalé par \eng{Look!} qui annonce une action en train de se dérouler ; forme : be + V-ing, attention à \eng{the workers are} (pluriel).\par
3. \eng{has recycled} : present perfect, justifié par \eng{since last year} (bilan d'une période qui continue jusqu'à maintenant) ; auxiliaire \eng{has} car le sujet \eng{our school} est singulier.\par
4. \eng{will increase} : dans une conditionnelle de type 1, si + present simple $\rightarrow$ futur avec \eng{will} ; prédiction sur l'avenir.\par
5. \eng{are thrown} : voix passive au present simple (be + participe passé) ; le sujet \eng{plastic bags} subit l'action. Participe passé irrégulier : throw — threw — thrown. Ne pas écrire \eng{are threw}.\par
Conclusion : repérer les marqueurs temporels (every Monday, Look!, since, if) permet de choisir le temps sans hésitation.
\end{exoresolu}

\courssub{Savoir-faire 4 : Rédiger à partir d'un modèle commenté}

\begin{methode}[Exploiter un modèle de composition et la grille d'évaluation]
\begin{enumerate}
\item \textbf{Lire le modèle} en repérant les trois parties et en soulignant les connecteurs et les temps utilisés.
\item \textbf{Relever les « briques » réutilisables} : phrases d'introduction (\eng{Nowadays...}), d'annonce de plan (\eng{This essay will examine...}), d'opinion (\eng{In my opinion...}), de conclusion (\eng{To conclude...}).
\item \textbf{Adapter, ne pas recopier} : garder la structure et les formules, changer le contenu selon le sujet donné.
\item \textbf{Auto-évaluer sa copie} avec la grille officielle : contenu et respect du sujet (environ 40\%), organisation et connecteurs (environ 20\%), correction grammaticale (environ 20\%), vocabulaire et orthographe (environ 20\%).
\end{enumerate}
\end{methode}

\begin{textebox}[Model composition — “Improving waste management in our town”]
Nowadays, garbage collection and waste disposal have become major problems in many Cameroonian towns. Heaps of rubbish can be seen near markets, schools and even hospitals. This essay will examine the causes of this situation and suggest some solutions, including recycling services.\par
First of all, garbage is not collected regularly because many councils do not have enough trucks or workers. Moreover, some citizens dump their waste anywhere instead of using dustbins. As a result, the streets are dirty and diseases can spread easily.\par
However, several solutions exist. The city council should organise a regular collection service in every neighbourhood. In addition, recycling centres could be created so that plastic bottles, metal cans and paper are sorted and reused. For instance, young people in Buea have started collecting plastic bottles to make paving stones. Finally, public awareness campaigns should teach everybody that a clean environment means good health.\par
To conclude, waste management is everybody's responsibility. If the authorities and the population work together, our towns will be cleaner, healthier and more beautiful.
\end{textebox}

\begin{exoresolu}[Analyse du modèle et rédaction guidée]
Énoncé — A. Answer these questions about the model composition above.\par
1. Underline the sentence that announces the plan in the introduction. 2. Find three linkers used in the essay and state their function. 3. Identify one sentence in the passive voice. 4. What solution is given as a concrete example?\par
B. Write the introduction (3--4 sentences) of a composition on: \eng{“Why should our school set up a recycling service?”}
\tcblower
Solution détaillée.\par
A. 1. La phrase d'annonce du plan : \eng{This essay will examine the causes of this situation and suggest some solutions, including recycling services.} Elle reprend exactement les deux tâches du sujet.\par
2. Trois connecteurs : \eng{First of all} (énumération, premier argument), \eng{As a result} (conséquence : rues sales, maladies), \eng{However} (opposition : passage du problème aux solutions). On peut citer aussi \eng{In addition}, \eng{For instance}, \eng{Finally}, \eng{To conclude}.\par
3. Voix passive : \eng{plastic bottles, metal cans and paper are sorted and reused} (be + participes passés \eng{sorted}, \eng{reused}) ; aussi \eng{heaps of rubbish can be seen}. La voix passive est naturelle ici car l'auteur de l'action importe peu.\par
4. L'exemple concret : des jeunes de Buea collectent des bouteilles plastiques pour fabriquer des pavés (\eng{young people in Buea have started collecting plastic bottles to make paving stones}) — present perfect pour une initiative récente.\par
B. Introduction rédigée (modèle de réponse) :\par
\eng{Nowadays, protecting the environment has become a priority for everyone, including students. Our school produces a large amount of waste every day, such as paper, plastic bottles and food remains. This essay will explain why setting up a recycling service in our school would be useful for the students, the school and the whole community.}\par
Cette introduction respecte la méthode : phrase d'accroche (\eng{Nowadays...}), présentation du problème, annonce du plan avec \eng{This essay will explain...}. Conclusion : un bon modèle sert de gabarit ; on en réutilise la charpente, jamais le contenu mot pour mot.
\end{exoresolu}

\courssub{Savoir-faire 5 : Relire et corriger sa production}

\begin{methode}[La relecture systématique en quatre passages]
\begin{enumerate}
\item \textbf{Passage 1 — Sens et sujet} : ai-je répondu à TOUTES les tâches du sujet ? Chaque paragraphe a-t-il une idée claire ?
\item \textbf{Passage 2 — Grammaire} : vérifier chaque verbe (temps, accord avec le sujet, -s à la 3\textsuperscript{e} personne), chaque nom (singulier/pluriel, a/an), la voix passive (be + participe passé).
\item \textbf{Passage 3 — Vocabulaire et prépositions} : \eng{to dispose of waste}, \eng{to depend on}, \eng{to be responsible for}, \eng{to be interested in} ; chasser les faux amis et les calques du français.
\item \textbf{Passage 4 — Orthographe et présentation} : \eng{environment} (pas \emph{enviroment}), \eng{government} (avec le n), \eng{which} (pas \emph{wich}), \eng{recycling} (orthographe : r-e-c-y-c-l-i-n-g), \eng{separate} (pas \emph{seperate}), \eng{responsibility} (pas \emph{responsability}) ; vérifier les majuscules en début de phrase et compter les mots.
\end{enumerate}
Réflexe : garder toujours 5 à 10 minutes pour cette relecture ; une copie bien relue gagne facilement 2 points.
\end{methode}

\begin{exoresolu}[Corriger un paragraphe d'élève]
Énoncé — The following paragraph contains six mistakes (grammar, spelling, prepositions). Find them and correct them.\par
\eng{“The garbages are not collected regulary in my quarter. Because of this, peoples throws their waste in the river, wich is very dangerous for the health. The governement should to create a recycling service and sensibilise the population.”}
\tcblower
Solution détaillée.\par
1. \eng{The garbages} $\rightarrow$ \eng{Garbage} (ou \eng{The garbage}) : \eng{garbage} est \textbf{indénombrable} en anglais, il ne prend jamais de -s. Comme \eng{information} ou \eng{advice}.\par
2. \eng{regulary} $\rightarrow$ \eng{regularly} : orthographe de l'adverbe (regular + -ly = regularly, avec -al-).\par
3. \eng{in my quarter} $\rightarrow$ \eng{in my neighbourhood} : faux ami ! \eng{quarter} existe mais signifie « quart » ou « trimestre » ; « quartier » se dit \eng{neighbourhood} (ou \eng{district}).\par
4. \eng{peoples throws} $\rightarrow$ \eng{people throw} : \eng{people} est déjà pluriel (« les gens »), donc verbe au pluriel sans -s ; \eng{peoples} ne s'emploie que pour « les peuples ».\par
5. \eng{wich} $\rightarrow$ \eng{which} : orthographe du pronom relatif (w-h-i-c-h).\par
6. \eng{should to create} $\rightarrow$ \eng{should create} : les auxiliaires modaux (should, must, can) sont suivis de la base verbale \textbf{sans to}. Et \eng{sensibilise} $\rightarrow$ \eng{raise the awareness of} ou \eng{sensitize} : calque du français « sensibiliser ».\par
Paragraphe corrigé : \eng{“Garbage is not collected regularly in my neighbourhood. Because of this, people throw their waste into the river, which is very dangerous for health. The government should create a recycling service and raise the awareness of the population.”}\par
Conclusion : la plupart des fautes des candidats francophones sont prévisibles (indénombrables, modaux + to, faux amis) ; une liste personnelle de ses erreurs types est un excellent outil de relecture.
\end{exoresolu}

\begin{attention}[Les 5 erreurs qui coûtent des points au Bac]
\begin{enumerate}
\item \textbf{Les indénombrables avec un -s} : \eng{garbage}, \eng{waste}, \eng{rubbish}, \eng{information}, \eng{advice}, \eng{pollution} ne prennent JAMAIS de -s. On écrit \eng{Waste is everywhere}, pas \emph{wastes are everywhere}. Pour compter : \eng{a piece of waste}, \eng{an item of information}.
\item \textbf{Modal + to} : après \eng{should, must, can, could, may}, la base verbale seule : \eng{The council should recycle more}, jamais \emph{should to recycle}. Mais attention : \eng{ought to recycle} (avec to).
\item \textbf{Faux amis et calques du français} : « sensibiliser » $\ne$ \emph{sensibilise} mais \eng{to raise awareness} ; « quartier » = \eng{neighbourhood}, pas \emph{quarter} ; « une décharge » = \eng{a landfill / a dump}, pas \emph{a discharge} ; « actuellement » = \eng{currently}, pas \emph{actually} (qui signifie « en fait »).
\item \textbf{Voix passive mal formée} : il faut be + participe passé : \eng{Plastic bottles are recycled}, pas \emph{are recycle} ni \emph{are threw}. Apprendre les participes passés irréguliers : throw — threw — thrown ; take — took — taken ; make — made — made.
\item \textbf{Orthographe des mots-clés du thème} : \eng{environment} (avec le n avant -ment), \eng{government} (g-o-v-e-r-n-m-e-n-t), \eng{which} (pas \emph{wich}), \eng{recycling} (r-e-c-y-c-l-i-n-g), \eng{separate} (pas \emph{seperate}), \eng{responsibility} (pas \emph{responsability}). Une faute sur le mot central du sujet, répétée dix fois, donne une très mauvaise impression au correcteur.
\end{enumerate}
\end{attention}

\begin{savaistu}[Gestion des déchets au Cameroun : initiatives locales]
Au Cameroun, plusieurs initiatives citoyennes et municipales émergent pour améliorer la collecte et le recyclage. À Yaoundé, l'association \eng{Youth for Environment} organise des campagnes de nettoyage et de tri dans les quartiers. À Douala, le projet \eng{Plastic Waste to Wealth} transforme les bouteilles plastiques en pavés écologiques, créant des emplois pour les jeunes. À Buea, la mairie a mis en place des \eng{recycling bins} dans les écoles et les marchés. Ces exemples concrets peuvent enrichir vos compositions : citez-les pour illustrer vos arguments (\eng{For instance, in Buea...}).
\end{savaistu}

\begin{experience}[Débat guidé : « Who is responsible for a clean town? »]
\begin{enumerate}
\item En groupes de 4, préparez 3 arguments pour chaque camp : \eng{The local council is responsible} vs \eng{Every citizen is responsible}.
\item Utilisez les connecteurs : \eng{First of all}, \eng{Moreover}, \eng{However}, \eng{For example}, \eng{To conclude}.
\item Un rapporteur par groupe présente la synthèse en 2 minutes.
\item Vote de la classe : quelle responsabilité prime ? Justifiez en anglais.
\end{enumerate}
Objectif : s'entraîner à l'oral (speaking) tout en réinvestissant le vocabulaire et les connecteurs de la leçon.
\end{experience}

\newpage

\coursec{Exercices}

\courssub{Je vérifie mes connaissances}

\begin{exercice}[QCM : la composition et le recyclage]
Pour chaque question, entoure la bonne réponse.
\begin{enumerate}
\item A composition is usually divided into :
\begin{itemize}
\item two parts
\item three parts : introduction, body, conclusion
\item five paragraphs of equal length
\end{itemize}
\item \eng{Garbage} est :
\begin{itemize}
\item un nom dénombrable (on peut dire \eng{garbages})
\item un nom indénombrable (jamais de \emph{s})
\item un adjectif
\end{itemize}
\item To add an idea, the best connector is :
\begin{itemize}
\item however
\item moreover
\item in conclusion
\end{itemize}
\item \eng{Recycling} signifie :
\begin{itemize}
\item l'enfouissement des déchets
\item le tri sélectif
\item le recyclage
\end{itemize}
\end{enumerate}
\end{exercice}

\begin{exercice}[Vrai ou faux — justifie ta réponse]
Réponds par \eng{True} ou \eng{False} et justifie en une phrase en anglais.
\begin{enumerate}
\item \eng{The government should to build more recycling plants.}
\item \eng{Plastic bottles can be recycled into new products.}
\item \eng{First of all} is used to conclude a composition.
\item \eng{Garbage is collected twice a week} is a passive sentence.
\end{enumerate}
\end{exercice}

\begin{exercice}[Vocabulaire : la chaîne des déchets]
Complète chaque définition avec le mot anglais qui convient : \eng{bin, landfill, waste, sorting, compost, collection}.
\begin{enumerate}
\item The container where you put your household refuse : \ldots
\item The place where garbage is buried : \ldots
\item The action of separating paper, plastic and glass : \ldots
\item Organic matter made from decomposed food, used as fertiliser : \ldots
\item The regular removal of garbage from homes by the council : \ldots
\item A general word for things we throw away : \ldots
\end{enumerate}
\end{exercice}

\begin{exercice}[Conjugaison à compléter]
Mets le verbe entre parenthèses à la forme correcte (présent simple, présent en \emph{-ing}, futur avec \eng{will} ou voix passive).
\begin{enumerate}
\item Every morning, the dustmen \ldots\ (collect) the bins in our neighbourhood.
\item Look ! People \ldots\ (burn) garbage in the street ; it is dangerous.
\item If we recycle more, we \ldots\ (reduce) pollution.
\item Last year, thousands of plastic bags \ldots\ (throw) into the Wouri river.
\item Our school \ldots\ (organise) a clean-up day next month.
\end{enumerate}
\end{exercice}

\courssub{Je m'entraîne}

\begin{exercice}[Traduction]
Traduis les phrases suivantes en anglais.
\begin{enumerate}
\item La collecte des ordures est un vrai problème à Douala.
\item Nous devrions trier nos déchets avant de les jeter.
\item Le recyclage du plastique crée des emplois pour les jeunes.
\item Les déchets sont ramassés deux fois par semaine dans mon quartier.
\item En conclusion, chaque citoyen doit participer à la protection de l'environnement.
\end{enumerate}
\end{exercice}

\begin{exercice}[De la voix active à la voix passive]
Transforme les phrases suivantes à la voix passive.
\begin{enumerate}
\item The council collects garbage twice a week.
\item The students cleaned the school yard yesterday.
\item Factories dump chemical waste into the river.
\item The recycling company will employ fifty young people.
\item People throw thousands of bottles away every day.
\end{enumerate}
\end{exercice}

\begin{exercice}[Texte à trous : les connecteurs]
Complète le texte suivant avec les connecteurs appropriés : \eng{first of all, however, moreover, as a result, in conclusion}.
\begin{textebox}[Keeping our city clean]
\ldots , every household should sort its waste into different bins. \ldots , the town council should organise garbage collection more regularly. \ldots , many neighbourhoods are still neglected ; \ldots , garbage accumulates in the streets and causes diseases. \ldots , I believe that keeping our city clean is everyone's responsibility.
\end{textebox}
\end{exercice}

\begin{exercice}[Appariements : moitiés de phrases]
Associe chaque début de phrase (1--5) à sa fin logique (a--e).
\begin{enumerate}
\item If we do not recycle plastic,
\item Garbage collection services should be
\item Burning waste in the open air
\item Thanks to recycling,
\item Before writing your composition,
\end{enumerate}
\begin{itemize}
\item[a.] old paper can be turned into new notebooks.
\item[b.] it will pollute the soil and the rivers for centuries.
\item[c.] make a plan with an introduction, a body and a conclusion.
\item[d.] releases toxic smoke that damages our lungs.
\item[e.] available in every neighbourhood, not only in the city centre.
\end{itemize}
\end{exercice}

\courssub{Je me prépare au Bac}

\begin{exercice}[Compréhension et langue — type Bac]
Lis le texte suivant, puis réponds aux questions.
\begin{textebox}[A new start for waste management in Yaoundé]
For many years, the mountains of garbage along the roads of Yaoundé have been a familiar sight. Blocked gutters, plastic bags flying in the wind and the smell of rotting waste have made life difficult for the inhabitants of the capital. However, things are slowly changing.

Last year, the city council signed an agreement with a private company to improve garbage collection. New trucks now visit the districts three times a week, and large green bins have been placed near markets and schools. Moreover, a recycling centre has been opened in the industrial zone. There, young workers sort plastic, metal and paper, which are then sold to factories and transformed into new products.

“We used to throw everything away,” explains Mama Béa, a shopkeeper at Mokolo market. “Now I keep my plastic bottles separately, and a young man collects them every Friday. He even pays me a small amount for them.”

Nevertheless, many challenges remain. Some residents still dump their waste into rivers, and the collection trucks cannot enter the narrow streets of certain neighbourhoods. As a result, community clean-up campaigns are organised every month, during which students, traders and civil servants work side by side.

If these efforts continue, Yaoundé could become one of the cleanest capitals in Central Africa. But success will depend on the cooperation of every citizen.
\end{textebox}
\textbf{A. Comprehension questions.} Answer in your own words.
\begin{enumerate}
\item What two problems caused by garbage are mentioned in the first paragraph ?
\item What has the city council done to improve the situation ? Give two measures.
\item How does Mama Béa benefit from the new system ?
\item Why are community clean-up campaigns necessary ?
\end{enumerate}
\textbf{B. Language work.}
\begin{enumerate}
\item Find in the text : (a) a connector expressing contrast ; (b) a connector expressing consequence.
\item Put into the passive voice : \eng{A private company collects the garbage.}
\item Find in the text a synonym of : \eng{inhabitants}, \eng{rubbish}.
\end{enumerate}
\end{exercice}

\begin{exercice}[Traduction — type Bac]
Traduis en français le paragraphe suivant, extrait du texte ci-dessus :
\begin{textebox}[Paragraph to translate]
Nevertheless, many challenges remain. Some residents still dump their waste into rivers, and the collection trucks cannot enter the narrow streets of certain neighbourhoods. As a result, community clean-up campaigns are organised every month, during which students, traders and civil servants work side by side.
\end{textebox}
Puis traduis en anglais : « Le recyclage est important parce qu'il protège notre environnement et crée des emplois. »
\end{exercice}

\begin{exercice}[Composition — sujet de Bac]
\textbf{Topic :} \eng{The importance of recycling in our community.}

Write a composition of about \textbf{250 words}. Your composition must include :
\begin{itemize}
\item an \textbf{introduction} presenting the problem of waste in your community ;
\item a \textbf{body} of two or three paragraphs explaining the benefits of recycling (environment, health, jobs) and suggesting practical solutions ;
\item a \textbf{conclusion} giving your personal opinion and a call to action.
\end{itemize}
Use at least \textbf{four different connectors} (for example : \eng{first of all, moreover, however, as a result, in conclusion}) and at least \textbf{two sentences in the passive voice}. Marks will be awarded for : organisation (5), vocabulary (5), grammar accuracy (5), originality and relevance (5). Total : 20 marks.
\end{exercice}

\newpage

\coursec{Corrigés des exercices}

\begin{corrige}[Exercice 1 — QCM : la composition et le recyclage]
\begin{enumerate}
\item \textbf{Bonne réponse : three parts : introduction, body, conclusion.} Une composition anglaise se compose classiquement de trois parties : l'\eng{introduction} (présentation du sujet), le \eng{body} (développement, deux ou trois paragraphes) et la \eng{conclusion} (opinion personnelle). Le nombre de paragraphes du développement n'est pas fixe.
\item \textbf{Bonne réponse : un nom indénombrable (jamais de \emph{s}).} \eng{Garbage} est un \emph{uncountable noun}, comme \eng{waste}, \eng{rubbish}, \eng{information} ou \eng{advice}. On dit \eng{some garbage}, \eng{a lot of garbage}, mais jamais \eng{a garbage} ni \eng{garbages}. Pour compter, on utilise \eng{a piece of rubbish} ou \eng{a bag of garbage}.
\item \textbf{Bonne réponse : moreover.} \eng{Moreover} (« de plus ») ajoute une idée dans le même sens. \eng{However} exprime un contraste et \eng{in conclusion} sert à conclure.
\item \textbf{Bonne réponse : le recyclage.} \eng{Recycling} = le recyclage. L'enfouissement se dit \eng{landfilling} ou \eng{burial}, et le tri sélectif se dit \eng{sorting} ou \eng{separate collection}.
\end{enumerate}
\end{corrige}

\begin{corrige}[Exercice 2 — Vrai ou faux]
\begin{enumerate}
\item \textbf{False.} Après un modal (\eng{should}, \eng{must}, \eng{can}), on utilise toujours la base verbale sans \eng{to}. La phrase correcte est : \eng{The government should build more recycling plants.} (« Le gouvernement devrait construire davantage d'usines de recyclage. »)
\item \textbf{True.} C'est une phrase correcte à la voix passive avec \eng{can be} + participe passé : \eng{Plastic bottles can be recycled into new products.} (« Les bouteilles en plastique peuvent être recyclées en nouveaux produits. »)
\item \textbf{False.} \eng{First of all} (« tout d'abord ») sert à introduire le premier point du développement, pas à conclure. Pour conclure, on utilise \eng{in conclusion}, \eng{to conclude} ou \eng{to sum up}.
\item \textbf{True.} La phrase \eng{Garbage is collected twice a week} est bien à la voix passive : sujet (\eng{garbage}) + auxiliaire \eng{be} conjugué (\eng{is}) + participe passé (\eng{collected}). L'agent n'est pas mentionné car il est évident (les services de la mairie).
\end{enumerate}
\end{corrige}

\begin{corrige}[Exercice 3 — Vocabulaire : la chaîne des déchets]
\begin{enumerate}
\item \eng{bin} — la poubelle, le conteneur.
\item \eng{landfill} — la décharge, le site d'enfouissement.
\item \eng{sorting} — le tri (des déchets).
\item \eng{compost} — le compost (engrais organique).
\item \eng{collection} — la collecte, le ramassage.
\item \eng{waste} — les déchets (mot général, indénombrable).
\end{enumerate}
\textbf{Phrase récapitulative à retenir :} \eng{Household waste is put in bins, collected by the council, sorted, and either recycled, composted or buried in a landfill.} (« Les déchets ménagers sont mis dans des poubelles, ramassés par la mairie, triés, puis soit recyclés, soit compostés, soit enfouis dans une décharge. »)
\end{corrige}

\begin{corrige}[Exercice 4 — Conjugaison à compléter]
\begin{enumerate}
\item \eng{Every morning, the dustmen \textbf{collect} the bins in our neighbourhood.} — Présent simple (habitude signalée par \eng{every morning}) ; sujet pluriel, donc verbe sans \emph{s}.
\item \eng{Look ! People \textbf{are burning} garbage in the street ; it is dangerous.} — Présent en \emph{-ing} (action en cours, signalée par \eng{Look !}) : \eng{be} + verbe-\emph{ing}.
\item \eng{If we recycle more, we \textbf{will reduce} pollution.} — Futur avec \eng{will} : dans une phrase conditionnelle de type 1, \eng{if} + présent simple, puis \eng{will} + base verbale dans la proposition principale.
\item \eng{Last year, thousands of plastic bags \textbf{were thrown} into the Wouri river.} — Voix passive au prétérit : \eng{were} (sujet pluriel) + participe passé irrégulier de \eng{throw} : \emph{throw -- threw -- thrown}.
\item \eng{Our school \textbf{will organise} a clean-up day next month.} — Futur avec \eng{will} (projet annoncé, marqueur \eng{next month}). On accepte aussi \eng{is going to organise}.
\end{enumerate}
\end{corrige}

\begin{corrige}[Exercice 5 — Traduction]
\begin{enumerate}
\item \eng{Garbage collection is a real problem in Douala.} — Attention : \eng{garbage} est indénombrable, donc \eng{is} et non \eng{are}.
\item \eng{We should sort our waste before throwing it away.} — Après la préposition \eng{before}, le verbe se met en \emph{-ing} ; \eng{waste} est indénombrable, donc on le reprend par \eng{it}.
\item \eng{Recycling plastic creates jobs for young people.} — Sujet en \emph{-ing} (\eng{recycling}) ; le verbe prend un \emph{s} au présent simple (sujet de troisième personne du singulier).
\item \eng{Waste is collected twice a week in my neighbourhood.} — Voix passive au présent : \eng{is collected} ; « deux fois par semaine » = \eng{twice a week} (et non \eng{two times}).
\item \eng{In conclusion, every citizen must take part in protecting the environment.} — « participer à » se traduit par \eng{take part in} + verbe en \emph{-ing} ; « doit » = modal \eng{must} + base verbale.
\end{enumerate}
\end{corrige}

\begin{corrige}[Exercice 6 — De la voix active à la voix passive]
\textbf{Rappel de méthode :} le complément d'objet direct devient sujet ; le verbe se met sous la forme \eng{be} (au même temps que le verbe actif) + participe passé ; l'ancien sujet devient complément d'agent introduit par \eng{by} (souvent omis).
\begin{enumerate}
\item \eng{Garbage is collected twice a week (by the council).} — Présent simple : \eng{is collected}.
\item \eng{The school yard was cleaned by the students yesterday.} — Prétérit : \eng{was cleaned} (sujet singulier).
\item \eng{Chemical waste is dumped into the river by factories.} — Présent simple ; \eng{waste} est indénombrable, donc \eng{is}.
\item \eng{Fifty young people will be employed by the recycling company.} — Futur passif : \eng{will be} + participe passé.
\item \eng{Thousands of bottles are thrown away every day.} — Présent passif, sujet pluriel : \eng{are thrown} ; on conserve la particule \eng{away} du verbe à particule \eng{throw away}.
\end{enumerate}
\end{corrige}

\begin{corrige}[Exercice 7 — Texte à trous : les connecteurs]
\begin{textebox}[Keeping our city clean — corrected version]
\textbf{First of all}, every household should sort its waste into different bins. \textbf{Moreover}, the town council should organise garbage collection more regularly. \textbf{However}, many neighbourhoods are still neglected ; \textbf{as a result}, garbage accumulates in the streets and causes diseases. \textbf{In conclusion}, I believe that keeping our city clean is everyone's responsibility.
\end{textebox}
\textbf{Justifications :}
\begin{itemize}
\item \eng{First of all} ouvre le développement (premier argument).
\item \eng{Moreover} ajoute un deuxième argument dans le même sens.
\item \eng{However} introduit un contraste (la réalité négative).
\item \eng{As a result} exprime la conséquence de cette négligence.
\item \eng{In conclusion} annonce l'opinion finale.
\end{itemize}
\end{corrige}

\begin{corrige}[Exercice 8 — Appariements : moitiés de phrases]
\begin{center}
\begin{tabularx}{\linewidth}{|c|X|}
\hline
\rowcolor{popblueL} Association & Phrase complète \\
\hline
1 -- b & \eng{If we do not recycle plastic, it will pollute the soil and the rivers for centuries.} \\
\hline
2 -- e & \eng{Garbage collection services should be available in every neighbourhood, not only in the city centre.} \\
\hline
3 -- d & \eng{Burning waste in the open air releases toxic smoke that damages our lungs.} \\
\hline
4 -- a & \eng{Thanks to recycling, old paper can be turned into new notebooks.} \\
\hline
5 -- c & \eng{Before writing your composition, make a plan with an introduction, a body and a conclusion.} \\
\hline
\end{tabularx}
\end{center}
\textbf{Remarques :} dans la phrase 1, le conditionnel de type 1 impose \eng{will} dans la proposition principale. Dans la phrase 3, le sujet est un gérondif (\eng{burning}), donc le verbe est au singulier : \eng{releases}.
\end{corrige}

\begin{corrige}[Exercice 9 — Compréhension et langue — type Bac]
\textbf{A. Comprehension questions} (réponses modèles ; toute formulation correcte avec les mots de l'élève est acceptée) :
\begin{enumerate}
\item \eng{Garbage blocks the gutters and produces a bad smell of rotting waste ; plastic bags also fly in the wind, which makes life difficult for the inhabitants.} (Deux problèmes suffisent : \eng{blocked gutters} et \eng{the smell of rotting waste}.)
\item \eng{The city council has signed an agreement with a private company : new trucks now collect the garbage three times a week, large green bins have been placed near markets and schools, and a recycling centre has been opened.} (Deux mesures suffisent.)
\item \eng{She keeps her plastic bottles separately and a young man collects them every Friday ; he even pays her a small amount for them, so she earns a little money.}
\item \eng{They are necessary because some residents still dump their waste into rivers and the trucks cannot enter the narrow streets of certain neighbourhoods.}
\end{enumerate}
\textbf{B. Language work :}
\begin{enumerate}
\item (a) Contraste : \eng{however}, \eng{nevertheless} ou \eng{but}. (b) Conséquence : \eng{as a result}.
\item \eng{The garbage is collected by a private company.} (Présent passif : \eng{is collected} ; \eng{garbage} indénombrable, donc singulier.)
\item \eng{inhabitants} = \eng{residents} ; \eng{rubbish} = \eng{garbage} (ou \eng{waste}).
\end{enumerate}
\textbf{Barème indicatif (sur 10) :} compréhension : 4 × 1,5 = 6 pts ; langue : connecteurs 1 pt, voix passive 1,5 pt, synonymes 1,5 pt.
\textbf{Points de vigilance :} répondre par des phrases complètes ; ne pas recopier tout le texte ; vérifier l'accord singulier avec \eng{garbage} et \eng{waste}.
\end{corrige}

\begin{corrige}[Exercice 10 — Traduction — type Bac]
\textbf{Traduction en français (proposition) :}

« Néanmoins, de nombreux défis demeurent. Certains habitants jettent encore leurs déchets dans les rivières, et les camions de ramassage ne peuvent pas pénétrer dans les rues étroites de certains quartiers. Par conséquent, des campagnes communautaires de nettoyage sont organisées chaque mois, au cours desquelles élèves, commerçants et fonctionnaires travaillent côte à côte. »

\textbf{Traduction en anglais :}

\eng{Recycling is important because it protects our environment and creates jobs.}

\textbf{Points de vigilance :} \eng{nevertheless} = « néanmoins » (et non « jamais ») ; \eng{to dump} = « déverser, jeter sauvagement » ; \eng{narrow streets} = « rues étroites » ; \eng{side by side} = « côte à côte ». Dans la phrase à traduire en anglais, \eng{recycling} est un sujet singulier : \eng{is}, \eng{protects}, \eng{creates} (avec \emph{s}). Ne pas écrire \eng{the recycling} : pas d'article devant un nom général indénombrable.
\textbf{Barème indicatif :} 5 pts pour le passage en français (fidélité 3, qualité de la langue 2) ; 3 pts pour la phrase en anglais (grammaire 2, vocabulaire 1).
\end{corrige}

\begin{corrige}[Exercice 11 — Composition — sujet de Bac]
\textbf{Proposition de rédaction modèle (environ 250 mots) :}
\begin{textebox}[The importance of recycling in our community — model answer]
In my community, waste is everywhere : plastic bags fly in the streets, gutters are blocked and rubbish is burnt in the open air. This situation is dangerous for our health and for the environment. Fortunately, recycling offers a simple and effective solution.

First of all, recycling protects the environment. When plastic, paper and metal are sorted and transformed into new products, less waste is thrown into rivers and fields. As a result, the soil remains fertile and the water stays clean. Moreover, recycling improves public health. If garbage is collected and recycled instead of being burnt, toxic smoke is no longer released into the air, and diseases caused by dirty water are reduced.

However, recycling is not only good for nature ; it also creates jobs. In many towns, young people are employed to collect plastic bottles, which are then sold to recycling companies. Old paper can be turned into notebooks, and organic waste can be transformed into compost for farmers. To encourage these activities, our community should place sorting bins in every neighbourhood and organise regular clean-up campaigns in schools and markets.

In conclusion, I strongly believe that recycling is essential for the future of our community. It protects the environment, keeps us healthy and gives work to our young people. Let us all sort our waste today : a clean community is everyone's responsibility.
\end{textebox}
\textbf{Analyse du modèle :}
\begin{itemize}
\item \textbf{Connecteurs utilisés (plus de quatre) :} \eng{first of all, as a result, moreover, however, in conclusion}.
\item \textbf{Voix passive (au moins deux occurrences) :} \eng{less waste is thrown} ; \eng{toxic smoke is no longer released} ; \eng{diseases \ldots\ are reduced} ; \eng{young people are employed} ; \eng{old paper can be turned}.
\item \textbf{Structure :} introduction (problème) ; développement en trois paragraphes (environnement et santé ; emplois ; solutions) ; conclusion (opinion + appel à l'action).
\end{itemize}
\textbf{Barème indicatif (sur 20) :} organisation et respect du plan : 5 pts ; richesse et exactitude du vocabulaire thématique : 5 pts ; correction grammaticale (temps, voix passive, accords) : 5 pts ; originalité et pertinence des arguments : 5 pts.
\textbf{Points de vigilance :} ne pas écrire \eng{garbages} ni \eng{wastes} ; ne pas mettre \eng{to} après \eng{should} ; accorder le verbe avec un sujet indénombrable (\eng{waste is}) ; éviter les phrases trop longues ; compter ses mots et relire avant de rendre la copie.
\end{corrige}

\newpage

\coursec{Fiche bilan}

\begin{popfigure}
\begin{tikzpicture}[mindmap, grow cyclic,
every node/.style={concept, font=\small},
concept color=popgreen,
level 1/.append style={level distance=3.4cm, sibling angle=51.4},
level 2/.append style={level distance=1.9cm, sibling angle=40}]
\node[concept color=popgreen!70, font=\small\bfseries, align=center] {Writing\\Garbage, Disposal\\and Recycling}
child[concept color=popblue!60] { node[align=center] {Vocabulary} 
  child { node[align=center, font=\scriptsize] {waste, bin,\\dump, landfill} }
  child { node[align=center, font=\scriptsize] {recycle, reuse,\\reduce} }
}
child[concept color=poporange!60] { node[align=center] {Analyse\\the topic}
  child { node[align=center, font=\scriptsize] {who? what?\\why? how?} }
  child { node[align=center, font=\scriptsize] {brainstorm\\+ plan} }
}
child[concept color=poppurple!60] { node[align=center] {Structure}
  child { node[align=center, font=\scriptsize] {introduction} }
  child { node[align=center, font=\scriptsize] {body\\(2--3 paragraphs)} }
  child { node[align=center, font=\scriptsize] {conclusion} }
}
child[concept color=popgold!70] { node[align=center] {Connectors}
  child { node[align=center, font=\scriptsize] {first, then,\\finally} }
  child { node[align=center, font=\scriptsize] {however,\\therefore} }
}
child[concept color=poppink!60] { node[align=center] {Tenses}
  child { node[align=center, font=\scriptsize] {simple present\\(facts)} }
  child { node[align=center, font=\scriptsize] {present perfect,\\future, passive} }
}
child[concept color=popblue!40] { node[align=center] {Model essay}
  child { node[align=center, font=\scriptsize] {commented\\example} }
}
child[concept color=poporange!40] { node[align=center] {Proofread}
  child { node[align=center, font=\scriptsize] {checklist\\+ grid} }
};
\end{tikzpicture}
\legende{Carte mentale de la leçon : rédiger une composition sur la collecte, l'élimination et le recyclage des déchets.}
\end{popfigure}

\begin{bilan}
\textbf{1. Lexique clé à connaître par cœur}
\begin{itemize}
\item \eng{garbage / rubbish / waste} : les ordures, les déchets ; \eng{a dustbin / a trash can} : une poubelle ; \eng{a dump / a landfill} : une décharge.
\item \eng{garbage collection} : la collecte des ordures ; \eng{a garbage truck} : un camion poubelle ; \eng{waste disposal} : l'élimination des déchets.
\item \eng{to recycle} : recycler ; \eng{recycling} : le recyclage ; \eng{to reuse} : réutiliser ; \eng{to reduce} : réduire ; \eng{to sort waste} : trier les déchets.
\item \eng{plastic bags / bottles} : sacs / bouteilles en plastique ; \eng{pollution} : la pollution ; \eng{to pollute} : polluer ; \eng{environment} : l'environnement.
\item \eng{to throw away / to dump} : jeter ; \eng{to burn waste} : brûler les déchets ; \eng{compost} : le compost ; \eng{a clean-up campaign} : une campagne de nettoyage.
\item \eng{unsanitary} : insalubre ; \eng{hygiene} : l'hygiène ; \eng{health hazard} : danger pour la santé ; \eng{diseases} : les maladies.
\end{itemize}

\textbf{2. Règles de grammaire à connaître par cœur}
\begin{center}
\begin{tabularx}{\linewidth}{|X|X|}\hline
\rowcolor{popblueL} \textbf{Règle} & \textbf{Exemple} \\ \hline
Simple present : faits généraux et habitudes (3\up{e} pers. : -\textbf{s}) & \eng{The truck collects garbage twice a week.} \\ \hline
Present perfect : bilan d'actions passées au présent (\eng{have/has} + participe) & \eng{The city has improved its recycling services.} \\ \hline
Voix passive : \eng{be} + participe passé (l'agent est souvent omis) & \eng{Plastic bottles are recycled into new products.} \\ \hline
Futur avec \eng{will} / \eng{be going to} : prévisions et projets & \eng{The council will build a new landfill.} \\ \hline
Modaux : \eng{must, should, can} pour obligation et conseil & \eng{We should sort our waste at home.} \\ \hline
\eng{There is / there are} + nom (dénombrable ou indénombrable) & \eng{There is too much garbage in our streets.} \\ \hline
\eng{much} + indénombrable ; \eng{many} + dénombrable & \eng{much waste / many plastic bags} \\ \hline
\end{tabularx}
\end{center}

\textbf{3. Connecteurs logiques indispensables}
\begin{itemize}
\item Ordre : \eng{first(ly), secondly, then, next, finally}.
\item Addition : \eng{moreover, furthermore, in addition, also}.
\item Opposition : \eng{however, whereas, on the other hand, although}.
\item Cause et conséquence : \eng{because, since, therefore, as a result, consequently}.
\item Conclusion : \eng{to conclude, in conclusion, to sum up, all in all}.
\end{itemize}

\textbf{4. Méthode express : rédiger la composition}
\begin{enumerate}
\item \textbf{Lire} le sujet deux fois et souligner les mots-clés (thème, tâche, longueur).
\item \textbf{Brainstormer} : noter 6 à 8 idées et le vocabulaire utile en anglais.
\item \textbf{Planifier} : introduction (présenter le problème), 2--3 paragraphes de développement (collecte, élimination, recyclage + solutions), conclusion (bilan + recommandation).
\item \textbf{Rédiger} au brouillon : une idée principale par paragraphe, reliée par des connecteurs.
\item \textbf{Relire} : accords, temps, orthographe, ponctuation, nombre de mots.
\end{enumerate}
\end{bilan}

\begin{aretenir}[Checklist avant le Bac]
Avant de rendre ma copie, je vérifie que :
\begin{itemize}
\item[$\square$] j'ai bien compris le sujet et je réponds exactement à la tâche demandée ;
\item[$\square$] ma composition a une introduction, un développement (2--3 paragraphes) et une conclusion ;
\item[$\square$] chaque paragraphe contient une idée principale claire ;
\item[$\square$] j'ai utilisé au moins cinq connecteurs logiques variés (\eng{firstly, however, therefore, moreover, in conclusion}) ;
\item[$\square$] j'ai employé le vocabulaire du thème : \eng{garbage collection, waste disposal, recycling, landfill, pollution} ;
\item[$\square$] mes verbes sont conjugués au bon temps (présent pour les faits, futur pour les solutions) ;
\item[$\square$] je n'ai pas oublié le -\textbf{s} de la troisième personne du singulier au présent ;
\item[$\square$] j'ai utilisé la voix passive là où elle est naturelle (\eng{waste is collected}) ;
\item[$\square$] je n'ai pas traduit mot à mot depuis le français (pas de \eng{*the garbages} : \eng{garbage} est indénombrable) ;
\item[$\square$] j'ai respecté le nombre de mots demandé ;
\item[$\square$] j'ai relu ma copie pour corriger l'orthographe et la ponctuation ;
\item[$\square$] mon écriture est lisible et ma présentation soignée.
\end{itemize}
\end{aretenir}

\findecours

\end{document}
